% Citizenship and Human Rights — Anglais 1re Tech — Cours signé OnBuch+
% Programme officiel MINESEC. Compiler avec : tectonic N04-citizenship-and-human-rights.tex
\newcommand{\DOCMATIERE}{Anglais}
\newcommand{\DOCNIVEAU}{1re Tech}
\newcommand{\DOCMODULE}{Anglais – classe de Première technique}
\newcommand{\DOCLECON}{N04}
\newcommand{\DOCTITRE}{Citizenship and Human Rights}
% OnBuch+ — préambule « cours signés » ENRICHI (compilé avec tectonic / XeLaTeX)
% Système visuel « Pop » d'OnBuch+ : fond crème (jamais blanc), bordures encre
% épaisses, ombres « dures » décalées, titres Archivo Black en capitales.
% Version MATHÉMATIQUES : graphes (pgfplots/tikz), boîtes pédagogiques
% (définition, propriété, méthode, attention, le savais-tu, exercice résolu,
% exercice, TP, bilan), page de garde et signature OnBuch+ ancrée en bas.
%
% Macros attendues dans main.tex AVANT \input{preamble} :
%   \DOCMATIERE  \DOCNIVEAU  \DOCMODULE  \DOCLECON (numéro)  \DOCTITRE
\documentclass[11pt]{article}

\usepackage{fontspec}
\usepackage[french]{babel}
\usepackage[a4paper,margin=2cm,top=3.4cm,bottom=2.8cm,headheight=1.4cm,footskip=1.3cm]{geometry}
\usepackage{amsmath,amssymb,amsfonts}
\usepackage{mathtools} % \xmapsto, \coloneqq... souvent utilisés par les modèles
\usepackage{cancel} % simplifications barrées (bilans de forces)
% \abs{z} (module, valeur absolue) et \norm{u} (norme), utilisés par les modèles
\providecommand{\abs}{}\let\abs\relax\DeclarePairedDelimiter{\abs}{\lvert}{\rvert}
\providecommand{\norm}{}\let\norm\relax\DeclarePairedDelimiter{\norm}{\lVert}{\rVert}
\usepackage[table]{xcolor} % [table] : fournit aussi \rowcolors (alternance de lignes)
\usepackage{pagecolor}
\usepackage{tikz}
\usetikzlibrary{arrows.meta,positioning,calc,shapes.geometric,shapes.misc,shapes.arrows,decorations.pathmorphing,decorations.pathreplacing,decorations.markings,patterns,shadows,mindmap,trees,fit,backgrounds,matrix,intersections,angles,quotes}
\usepackage{pgfplots}
\usepackage[version=4]{mhchem} % équations nucléaires \ce{^{235}_{92}U}
\usepackage[european,siunitx]{circuitikz} % schémas électriques
\pgfplotsset{compat=1.18}
\usepgfplotslibrary{fillbetween,groupplots} % aires entre courbes (calcul intégral)
\usepackage{enumitem}
\usepackage{fancyhdr}
% [most] charge toutes les bibliothèques tcolorbox (listings, minted, raster,
% documentation...) et gonfle inutilement la mémoire TeX sur un document long
% (~300 boîtes) : on ne charge que ce qui est réellement utilisé.
\usepackage[skins,breakable]{tcolorbox}
\usepackage{graphicx}
\usepackage{colortbl}
\usepackage{booktabs}
\usepackage{tabularx}
\usepackage{multirow}
\usepackage{diagbox} % case d'en-tête diagonale des tableaux à double entrée
% Colonnes extensibles centrée / gauche / droite (C, L, R), souvent utilisées par les modèles
\newcolumntype{C}{>{\centering\arraybackslash}X}
\newcolumntype{L}{>{\raggedright\arraybackslash}X}
\newcolumntype{R}{>{\raggedleft\arraybackslash}X}
\usepackage{array}
\usepackage{multicol}
\usepackage{siunitx}
% parse-numbers=false : accepte aussi \SI{2,5 \times 10^{-3}}{...}
\sisetup{locale=FR,output-decimal-marker={,},group-minimum-digits=4,parse-numbers=false}
\DeclareSIUnit\ohm{\ensuremath{\symup{\Omega}}} % U+2126 absent de la police
\DeclareSIUnit\division{div} % réglages d'oscilloscope (ms/div, V/div)
\DeclareSIUnit\tr{tr} % tours (vitesse de rotation, ex. \SI{1500}{\tr\per\minute})
\DeclareSIUnit\molar{M} % concentration molaire (ex. \SI{3}{\molar}, \SI{1}{\micro\molar})
\DeclareSIUnit\an{an} % année (le modèle confond parfois avec \year, un compteur TeX) — ex. \SI{5}{\kilo\meter\per\an}
\DeclareSIUnit\FCFA{FCFA} % franc CFA (ex. \SI{500}{\FCFA\per\kilo\gram})
\DeclareSIUnit\degreeBrix{\textdegree Brix} % teneur en sucre d'un sirop/jus (réfractométrie)
\DeclareSIUnit\month{mois} % le modèle utilise parfois \month, un compteur TeX, comme unité de durée
\DeclareSIUnit\microliter{\micro\liter} % le modèle écrit parfois \microliter en un seul token
\DeclareSIUnit\million{millions} % dénombrement (ex. \SI{15}{\million\per\mL} spermatozoïdes)
\usepackage{qrcode}
% \degree : commande fréquemment utilisée par les modèles pour un angle ou une
% température en dehors de \SI{}{} (non fournie par les paquets chargés).
\providecommand{\degree}{^{\circ}}
% \qed (fin de démonstration), utilisé par les modèles sans amsthm
\providecommand{\qed}{\hfill$\square$}
% \barre : double barre des valeurs interdites dans un tableau de variations (tkz-tab)
\providecommand{\barre}{\ensuremath{\|}}
% environnement proof (amsthm non chargé) utilisé par les modèles
\ifdefined\proof\else\newenvironment{proof}[1][Démonstration]{\par\noindent\textit{#1.}\ }{\qed\par}\fi
\usepackage{lastpage}
\usepackage{hyperref}
\hypersetup{hidelinks}

% ============================================================
% Palette « Pop » OnBuch+ — mêmes hex que la classe P (plus_ui.dart)
% ============================================================
\definecolor{poporange}{HTML}{F59321}
\definecolor{popdark}{HTML}{E07A0C}
\definecolor{popcream}{HTML}{FFF6E8}
\definecolor{popink}{HTML}{1C1714}
\definecolor{poppurple}{HTML}{7A5AE0}
\definecolor{popgreen}{HTML}{2E9E6A}
\definecolor{poppink}{HTML}{E0457B}
\definecolor{popblue}{HTML}{2F7DE1}
\definecolor{popgold}{HTML}{C98A12}
\definecolor{popmuted}{HTML}{8A7F76}
\definecolor{popline}{HTML}{EADFD2}
\definecolor{popgray}{HTML}{8A7F76}
\colorlet{popgrey}{popgray}
% Alias tolérants (le modèle nomme parfois les couleurs différemment)
\colorlet{popwhite}{white}
\colorlet{popblack}{popink}
\colorlet{popred}{poppink}
\colorlet{popyellow}{popgold}
% Teintes claires (fonds de tableaux/encadrés) — le modèle nomme parfois
% les couleurs avec un suffixe L (ex. popblueL) sans les avoir définies.
\colorlet{poporangeL}{poporange!20}
\colorlet{popdarkL}{popdark!20}
\colorlet{poppurpleL}{poppurple!20}
\colorlet{popgreenL}{popgreen!20}
\colorlet{poppinkL}{poppink!20}
\colorlet{popblueL}{popblue!20}
\colorlet{popgoldL}{popgold!20}
\colorlet{popredL}{popred!20}
\colorlet{popgrayL}{popgray!20}
% Teintes claires pour fonds de tableaux / zones
\colorlet{popblueL}{popblue!10!white}
\colorlet{poporangeL}{poporange!14!white}
\colorlet{popgreenL}{popgreen!12!white}
\colorlet{poppurpleL}{poppurple!12!white}
\colorlet{poppinkL}{poppink!12!white}
\colorlet{popgoldL}{popgold!14!white}

\pagecolor{popcream}

% ============================================================
% Typographie
% ============================================================
\defaultfontfeatures{Ligatures=TeX}
\setmainfont{PlusJakartaSans-Medium.ttf}[Path=../../../fonts/,BoldFont=PlusJakartaSans-Bold.ttf]
% Maths en sans-serif (Fira Math) avec chiffres et lettres droites de Plus
% Jakarta, pour que les formules se fondent dans le texte.
\usepackage[math-style=ISO,bold-style=ISO]{unicode-math}
\setmathfont{FiraMath-Regular.otf}[Path=../../../fonts/]
\setmathfont{PlusJakartaSans-Medium.ttf}[Path=../../../fonts/,range={up/num,up/{Latin,latin},\mathup}]
\usepackage{icomma} % virgule décimale sans espace en mode math
% Caractères mathématiques Unicode absents des polices de texte (Plus Jakarta,
% Archivo Black) : rendus par leur équivalent mathématique, en texte comme en
% formule — sinon ils s'affichent en carré vide (ex. « Arithmétique dans ℤ »).
\usepackage{newunicodechar}
\newunicodechar{ℝ}{\ensuremath{\mathbb{R}}}
\newunicodechar{ℤ}{\ensuremath{\mathbb{Z}}}
\newunicodechar{ℂ}{\ensuremath{\mathbb{C}}}
\newunicodechar{ℕ}{\ensuremath{\mathbb{N}}}
\newunicodechar{ℚ}{\ensuremath{\mathbb{Q}}}
\newunicodechar{𝔻}{\ensuremath{\mathbb{D}}}
\newunicodechar{α}{\ensuremath{\alpha}}
\newunicodechar{β}{\ensuremath{\beta}}
\newunicodechar{γ}{\ensuremath{\gamma}}
\newunicodechar{δ}{\ensuremath{\delta}}
\newunicodechar{ε}{\ensuremath{\varepsilon}}
\newunicodechar{θ}{\ensuremath{\theta}}
\newunicodechar{λ}{\ensuremath{\lambda}}
\newunicodechar{μ}{\ensuremath{\mu}}
\newunicodechar{σ}{\ensuremath{\sigma}}
\newunicodechar{τ}{\ensuremath{\tau}}
\newunicodechar{φ}{\ensuremath{\varphi}}
\newunicodechar{ψ}{\ensuremath{\psi}}
\newunicodechar{ω}{\ensuremath{\omega}}
\newunicodechar{Σ}{\ensuremath{\Sigma}}
% majuscules produites par \MakeUppercase dans les titres (α-aminés -> Α)
\newunicodechar{Α}{\ensuremath{\alpha}}
\newunicodechar{Β}{\ensuremath{\beta}}
\newunicodechar{Γ}{\ensuremath{\Gamma}}
\newunicodechar{Δ}{\ensuremath{\Delta}}
\newunicodechar{Ε}{\ensuremath{\varepsilon}}
\newunicodechar{Θ}{\ensuremath{\Theta}}
\newunicodechar{Λ}{\ensuremath{\Lambda}}
\newunicodechar{Μ}{\ensuremath{\mu}}
\newunicodechar{Τ}{\ensuremath{\tau}}
\newunicodechar{Φ}{\ensuremath{\Phi}}
\newunicodechar{Ψ}{\ensuremath{\Psi}}
\newunicodechar{Ω}{\ensuremath{\Omega}}
\newunicodechar{↦}{\ensuremath{\mapsto}}
\newunicodechar{∈}{\ensuremath{\in}}
\newunicodechar{∉}{\ensuremath{\notin}}
\newunicodechar{∧}{\ensuremath{\wedge}}
\newunicodechar{⇔}{\ensuremath{\Leftrightarrow}}
\newunicodechar{⟺}{\ensuremath{\Longleftrightarrow}}
\newunicodechar{⇒}{\ensuremath{\Rightarrow}}
\newunicodechar{⟹}{\ensuremath{\Longrightarrow}}
\newunicodechar{∘}{\ensuremath{\circ}}
\newunicodechar{∪}{\ensuremath{\cup}}
\newunicodechar{∩}{\ensuremath{\cap}}
\newunicodechar{≡}{\ensuremath{\equiv}}
\newunicodechar{⊂}{\ensuremath{\subset}}
\newunicodechar{⊥}{\ensuremath{\perp}}
\newunicodechar{∃}{\ensuremath{\exists}}
\newunicodechar{∀}{\ensuremath{\forall}}
\newunicodechar{∛}{\ensuremath{\sqrt[3]{\,}}}
\newunicodechar{∜}{\ensuremath{\sqrt[4]{\,}}}
\newunicodechar{⃗}{\ensuremath{{}^{\to}}}
\newunicodechar{ⁿ}{\ifmmode^{n}\else\textsuperscript{\ensuremath{n}}\fi}
\newunicodechar{ˣ}{\ifmmode^{x}\else\textsuperscript{\ensuremath{x}}\fi}
\newunicodechar{ᵇ}{\ifmmode^{b}\else\textsuperscript{\ensuremath{b}}\fi}
\newunicodechar{ʳ}{\ifmmode^{r}\else\textsuperscript{\ensuremath{r}}\fi}
\newunicodechar{ᵘ}{\ifmmode^{u}\else\textsuperscript{\ensuremath{u}}\fi}
\newunicodechar{ʸ}{\ifmmode^{y}\else\textsuperscript{\ensuremath{y}}\fi}
\newunicodechar{ᵃ}{\ifmmode^{a}\else\textsuperscript{\ensuremath{a}}\fi}
\newunicodechar{ᵉ}{\ifmmode^{e}\else\textsuperscript{\ensuremath{e}}\fi}
\newunicodechar{ᵏ}{\ifmmode^{k}\else\textsuperscript{\ensuremath{k}}\fi}
\newunicodechar{ᵖ}{\ifmmode^{p}\else\textsuperscript{\ensuremath{p}}\fi}
\newunicodechar{ᴺ}{\ifmmode^{N}\else\textsuperscript{\ensuremath{N}}\fi}
\newunicodechar{ᶜ}{\ifmmode^{c}\else\textsuperscript{\ensuremath{c}}\fi}
\newunicodechar{ʰ}{\ifmmode^{h}\else\textsuperscript{\ensuremath{h}}\fi}
\newunicodechar{ᵠ}{\ifmmode^{q}\else\textsuperscript{\ensuremath{q}}\fi}
\newunicodechar{ₙ}{\ifmmode_{n}\else\textsubscript{\ensuremath{n}}\fi}
\newunicodechar{ₐ}{\ifmmode_{a}\else\textsubscript{\ensuremath{a}}\fi}
\newunicodechar{ᵢ}{\ifmmode_{i}\else\textsubscript{\ensuremath{i}}\fi}
\newunicodechar{ᵣ}{\ifmmode_{r}\else\textsubscript{\ensuremath{r}}\fi}
\newunicodechar{ₖ}{\ifmmode_{k}\else\textsubscript{\ensuremath{k}}\fi}
\newunicodechar{ᵦ}{\ifmmode_{\beta}\else\textsubscript{\ensuremath{\beta}}\fi}
\newunicodechar{ₓ}{\ifmmode_{x}\else\textsubscript{\ensuremath{x}}\fi}
\newfontfamily\popdisplay{ArchivoBlack-Regular.ttf}[Path=../../../fonts/]
\newfontfamily\poplogo{SpaceGrotesk-Bold.ttf}[Path=../../../fonts/]
\newfontfamily\popsemi{PlusJakartaSans-SemiBold.ttf}[Path=../../../fonts/]
\newfontfamily\popxbold{PlusJakartaSans-ExtraBold.ttf}[Path=../../../fonts/]
\newfontfamily\popxboldls{PlusJakartaSans-ExtraBold.ttf}[Path=../../../fonts/,LetterSpace=3.0]

\setlist[enumerate]{leftmargin=1.7em,itemsep=5pt,topsep=4pt}
\setlist[itemize]{leftmargin=1.7em,itemsep=4pt,topsep=4pt,label={\color{poporange}\textbullet}}
\setlength{\parindent}{0pt}
\setlength{\parskip}{5pt}
\renewcommand{\arraystretch}{1.25}

% pgfplots : style Pop par défaut
\pgfplotsset{
  every axis/.append style={axis line style={popink,line width=1pt},tick style={popink},
    grid style={popline,line width=0.5pt},label style={font=\small},tick label style={font=\footnotesize}},
  popcurve/.style={poporange,line width=1.8pt,no markers,smooth},
}
% Même style disponible dans un \draw TikZ simple (hors pgfplots)
\tikzset{popcurve/.style={poporange,line width=1.8pt}}
\tikzset{popgrid/.style={popline,very thin}}
\colorlet{popcurve}{poporange} % le nom du style est parfois utilisé comme couleur

% ============================================================
% Pill / squiggle
% ============================================================
\newcommand{\poppill}[3]{%
  \tikz[baseline=-0.55ex]\node[draw=popink,line width=1.2pt,rounded corners=2.6pt,%
    fill=#2,inner xsep=6.5pt,inner ysep=3pt]{%
    \popxboldls\fontsize{7}{7}\selectfont\color{#3}\MakeUppercase{#1}};%
}
\newcommand{\popsquiggle}[1]{%
  \tikz[baseline=-0.4ex]{
    \draw[#1,line width=2.6pt,line cap=round]
      plot [smooth] coordinates {(0,0.09) (0.34,0.01) (0.68,0.21) (1.02,0.01) (1.36,0.10)};
  }%
}
% Mot-clé surligné (marqueur orange sous le texte)
\newcommand{\cle}[1]{{\popxbold\color{popdark}#1}}

% ============================================================
% En-tête / pied
% ============================================================
\pagestyle{fancy}
\fancyhf{}
\renewcommand{\headrulewidth}{0pt}
\renewcommand{\footrulewidth}{0pt}
\lhead{\raisebox{-0.15ex}{\poplogo\fontsize{13}{13}\selectfont\color{popink}OnBuch\color{poporange}+}}
\rhead{\poppill{\DOCMATIERE\ \textbullet\ \DOCNIVEAU\ \textbullet\ Leçon \DOCLECON}{white}{popink}}
\lfoot{\popxbold\fontsize{7.6}{7.6}\selectfont\color{popmuted}\MakeUppercase{Cours signé OnBuch+}\ \textbullet\ {\popsemi\DOCTITRE}}
\rfoot{\tikz[baseline=-0.6ex]\node[rounded corners=8.5pt,fill=popink,minimum height=17pt,minimum width=17pt,inner xsep=5pt,inner sep=0]{\popxbold\fontsize{8}{8}\selectfont\color{popcream}\thepage\,/\,\pageref*{LastPage}};}

% ============================================================
% Titres
% ============================================================
\newcommand{\chaptitre}[2]{%
  \begin{tcolorbox}[enhanced,colback=poporange,colframe=popink,boxrule=2.6pt,arc=11pt,
      drop shadow=popink,left=15pt,right=15pt,top=12pt,bottom=11pt,before skip=3pt,after skip=18pt]
    \poppill{#2}{popink}{popcream}\par\vspace{9pt}
    {\raggedright\hyphenpenalty=10000\exhyphenpenalty=10000\popdisplay\fontsize{22}{24}\selectfont\color{popcream}\MakeUppercase{#1}\par}
  \end{tcolorbox}
}
% Section numérotée : pastille orange avec le numéro + titre Archivo
\newcounter{popsec}
\newcommand{\coursec}[1]{%
  \stepcounter{popsec}\setcounter{popsub}{0}%
  \par\vspace{16pt}\noindent
  \begin{minipage}{\linewidth}
  \tikz[baseline=-0.7ex]\node[draw=popink,line width=1.6pt,fill=poporange,rounded corners=4pt,minimum size=22pt,inner sep=0]{\popdisplay\fontsize{12}{12}\selectfont\color{popcream}\thepopsec};%
  \hspace{8pt}{\popdisplay\fontsize{15}{17}\selectfont\color{popink}\MakeUppercase{#1}}\par
  \vspace{4pt}\noindent{\color{poporange}\rule{36pt}{3.6pt}}
  \end{minipage}\par\nopagebreak\vspace{8pt}
}
\newcounter{popsub}
\newcommand{\courssub}[1]{%
  \stepcounter{popsub}%
  \par\vspace{9pt}\noindent{\popxbold\fontsize{11.5}{13}\selectfont\color{popink}{\color{poporange}\thepopsec.\thepopsub}\hspace{6pt}#1}\par\nopagebreak\vspace{4pt}
}

% ============================================================
% Boîtes pédagogiques — même recette : fond blanc, bordure épaisse colorée,
% ombre dure encre, badge plein. Titre optionnel [#1].
% ============================================================
\tcbset{popbase/.style={breakable,enhanced,colback=white,boxrule=2.3pt,arc=9pt,
  shadow={3pt}{-3pt}{0pt}{popink},left=14pt,right=14pt,top=11pt,bottom=11pt,before skip=11pt,after skip=15pt,
  fontupper=\popsemi\fontsize{10.6}{14.5}\selectfont\color{popink},
  fontlower=\popsemi\fontsize{10.6}{14.5}\selectfont\color{popink}}}
% Coche dessinée (le glyphe ✓ n'existe pas dans les polices du document)
\newcommand{\popcheck}[1]{\tikz[baseline=-0.5ex]\draw[#1,line width=1.6pt,line cap=round,line join=round](-0.13,0.0)--(-0.04,-0.09)--(0.13,0.1);}
\providecommand{\coche}{\popcheck{popgreen}}
\providecommand{\resultat}[1]{\par\smallskip\noindent\fcolorbox{popgreen}{popgreenL}{\parbox{\dimexpr\linewidth-2\fboxsep-2\fboxrule\relax}{\textbf{Résultat.} #1}}\par\smallskip}
\newcommand{\popboxhead}[3]{\poppill{#1}{#2}{white}\ifx\relax#3\relax\else\hspace{6pt}{\popxbold\fontsize{10.6}{12}\selectfont\color{popink}#3}\fi\par\vspace{7pt}}

\newtcolorbox{aretenir}[1][]{popbase,colframe=popgreen,before upper={\popboxhead{À retenir}{popgreen}{#1}}}
\newtcolorbox{exemplebox}[1][]{popbase,colframe=poporange,before upper={\popboxhead{Exemple}{poporange}{#1}}}
\newtcolorbox{definition}[1][]{popbase,colframe=popblue,before upper={\popboxhead{Définition}{popblue}{#1}}}
\newtcolorbox{propriete}[1][]{popbase,colframe=poppurple,before upper={\popboxhead{Propriété}{poppurple}{#1}}}
\newtcolorbox{methode}[1][]{popbase,colframe=popgold,before upper={\popboxhead{Méthode}{popgold}{#1}}}
\newtcolorbox{attention}[1][]{popbase,colframe=poppink,before upper={\popboxhead{Attention}{poppink}{#1}}}
\newtcolorbox{experience}[1][]{popbase,colframe=popblue,colback=popblueL,before upper={\popboxhead{Expérience}{popblue}{#1}}}
% « Le savais-tu ? » : carte violette pleine (contexte, culture, Cameroun)
\newtcolorbox{savaistu}[1][]{popbase,colback=poppurple,colframe=popink,
  fontupper=\popsemi\fontsize{10.4}{14.2}\selectfont\color{white},
  before upper={\poppill{Le savais-tu\,?}{white}{poppurple}\ifx\relax#1\relax\else\hspace{6pt}{\popxbold\color{white}#1}\fi\par\vspace{7pt}}}
% Exercice résolu : énoncé en haut, \tcblower puis la solution sur fond crème
\newcounter{popexo}\newcounter{popexr}
\newtcolorbox{exoresolu}[1][]{popbase,colframe=poporange,colbacklower=popcream,
  segmentation style={popink,line width=1.2pt,dashed},
  before upper={\stepcounter{popexr}\popboxhead{Exercice résolu \thepopexr}{poporange}{#1}},
  before lower={\poppill{Solution}{popink}{popcream}\par\vspace{6pt}}}
% Exercice (non résolu en place) — la correction est dans la partie Corrigés
\newtcolorbox{exercice}[1][]{popbase,colframe=popink,
  before upper={\stepcounter{popexo}\popboxhead{Exercice \thepopexo}{popink}{#1}}}
% Corrigé
\newtcolorbox{corrige}[1][]{popbase,colframe=popgreen,colback=popgreenL,
  before upper={\popboxhead{Corrigé}{popgreen}{#1}}}
% Objectifs / prérequis / bilan
\newtcolorbox{objectifs}{popbase,colframe=popink,colback=poporangeL,
  before upper={\popboxhead{Objectifs de la leçon}{popink}{}}}
\newtcolorbox{prerequis}{popbase,colframe=popink,colback=popblueL,
  before upper={\popboxhead{Prérequis}{popblue}{}}}
\newtcolorbox{bilan}{popbase,colframe=popink,colback=white,boxrule=2.8pt,
  before upper={\popboxhead{Fiche bilan}{poporange}{L'essentiel à connaître pour l'examen}}}
% Figure encadrée avec légende
\newtcolorbox{popfigure}[1][]{enhanced,breakable=false,colback=white,colframe=popline,boxrule=1.4pt,arc=8pt,
  left=10pt,right=10pt,top=10pt,bottom=8pt,before skip=10pt,after skip=12pt,halign=center,
  lower separated=false,halign lower=center,
  fontlower=\popsemi\fontsize{9}{11.5}\selectfont\color{popmuted},
  #1}
\newcommand{\legende}[1]{\par\vspace{6pt}{\popsemi\fontsize{9}{11.5}\selectfont\color{popmuted}#1}}

% ============================================================
% Signature OnBuch+
% ============================================================
\newcommand{\poplogoinline}[1]{{\poplogo\fontsize{#1}{#1}\selectfont\color{popink}OnBuch\color{poporange}+}}
\newcommand{\signatureonbuch}{%
  \begin{tcolorbox}[enhanced,colback=popink,colframe=popink,boxrule=2.2pt,arc=10pt,
      drop shadow=poporange,left=14pt,right=14pt,top=10pt,bottom=10pt,before skip=0pt,after skip=0pt]
    \tikz[baseline=-0.7ex]\node[circle,fill=popgreen,draw=popcream,line width=1.3pt,minimum size=22pt,inner sep=0]{\popcheck{white}};%
    \hspace{9pt}\begin{minipage}[c]{0.62\linewidth}
      {\popxbold\fontsize{12}{13}\selectfont\color{popcream}Signé\ \textbullet\ {\poplogo OnBuch\color{poporange}+}}\par\vspace{2pt}
      {\raggedright\popsemi\fontsize{8}{10.5}\selectfont\color{popline}\MakeUppercase{Équipe pédagogique OnBuch+}\\\MakeUppercase{Conforme au programme officiel MINESEC}\par}
    \end{minipage}\hfill
    {\poplogo\fontsize{22}{22}\selectfont\color{popcream}OnBuch\color{poporange}+}
  \end{tcolorbox}%
}

% ============================================================
% Page de garde
% #1 titre · #2 matière · #3 niveau · #4 module · #5 n° leçon · #6 durée ·
% #7 description / objectifs (texte ou itemize)
% ============================================================
\newcommand{\pagegarde}[7]{%
  \thispagestyle{empty}
  \newgeometry{a4paper,margin=1.7cm,top=1.6cm,bottom=1.4cm}%
  \begin{tikzpicture}[remember picture,overlay]
    \fill[poporange!18] ([xshift=-3.2cm,yshift=-3.6cm]current page.north east) circle (4.6cm);
    \fill[poppurple!14] ([xshift=2.4cm,yshift=7.6cm]current page.south west) circle (3.8cm);
  \end{tikzpicture}%
  {\poplogo\fontsize{30}{30}\selectfont\color{popink}OnBuch\color{poporange}+}\hfill
  \raisebox{0.6ex}{\poppill{Cours signé \textbullet\ Premium}{popink}{popcream}}\par
  \vspace{1pt}\popsquiggle{poppurple}\par
  \vspace{8pt}
  \begin{tcolorbox}[enhanced,colback=poporange,colframe=popink,boxrule=2.8pt,arc=13pt,
      drop shadow=popink,left=18pt,right=18pt,top=12pt,bottom=13pt,after skip=10pt]
    \poppill{#2 \textbullet\ #3}{popink}{popcream}\hspace{5pt}\poppill{#4}{white}{popink}\par\vspace{8pt}
    {\popdisplay\fontsize{14}{14}\selectfont\color{popink}LEÇON #5}\par\vspace{4pt}
    {\raggedright\hyphenpenalty=10000\exhyphenpenalty=10000\popdisplay\fontsize{25}{27}\selectfont\color{popcream}\MakeUppercase{#1}\par}
    \vspace{8pt}
    {\popxbold\fontsize{9.5}{9.5}\selectfont\color{popink}\MakeUppercase{Durée conseillée : #6}}
  \end{tcolorbox}
  \begin{tcolorbox}[enhanced,colback=white,colframe=popink,boxrule=2.2pt,arc=10pt,
      drop shadow=popink,left=16pt,right=16pt,top=10pt,bottom=10pt,after skip=8pt]
    \poppill{Dans cette leçon}{poppurple}{white}\par\vspace{5pt}
    \popsemi\fontsize{9.3}{12.6}\selectfont\color{popink}#7
  \end{tcolorbox}
  \noindent\begin{minipage}[t]{0.487\linewidth}
    \begin{tcolorbox}[enhanced,colback=popgreen,colframe=popink,boxrule=2.2pt,arc=10pt,
        drop shadow=popink,left=11pt,right=11pt,top=10pt,bottom=10pt,height=3.1cm,valign=center]
      \tcbox[colback=white,colframe=popink,boxrule=1.3pt,arc=5pt,boxsep=0pt,nobeforeafter,
        left=3pt,right=3pt,top=3pt,bottom=3pt,valign=center]{\qrcode[height=1.9cm]{https://play.google.com/store/apps/details?id=com.luvvix.onbuch&hl=fr}}%
      \hspace{8pt}\begin{minipage}[c]{0.5\linewidth}
        {\popdisplay\fontsize{11}{12.5}\selectfont\color{white}\MakeUppercase{Télécharge OnBuch}}\par\vspace{4pt}
        {\raggedright\popsemi\fontsize{8.4}{11}\selectfont\color{white}Gratuit \textbullet\ Google Play\\Tuteur IA, annales, concours, résultats\par}
      \end{minipage}
    \end{tcolorbox}
  \end{minipage}\hfill
  \begin{minipage}[t]{0.487\linewidth}
    \begin{tcolorbox}[enhanced,colback=poppurple,colframe=popink,boxrule=2.2pt,arc=10pt,
        drop shadow=popink,left=11pt,right=11pt,top=10pt,bottom=10pt,height=3.1cm,valign=center]
      \tikz[baseline=-0.7ex]\node[circle,fill=white,draw=popink,line width=1.3pt,minimum size=1.6cm,inner sep=0]{\popdisplay\fontsize{24}{24}\selectfont\color{poppurple}+};%
      \hspace{8pt}\begin{minipage}[c]{0.6\linewidth}
        {\popdisplay\fontsize{11}{12.5}\selectfont\color{white}\MakeUppercase{Passe à OnBuch+}}\par\vspace{4pt}
        {\raggedright\popsemi\fontsize{8.4}{11}\selectfont\color{white}Tous les cours signés, Tuteur IA illimité, fiches PDF — onglet OnBuch+ de l'app\par}
      \end{minipage}
    \end{tcolorbox}
  \end{minipage}
  \vspace{8pt}
  \popcontenu
  \vfill
  \signatureonbuch
  \restoregeometry
  \newpage
}

% Rangée « Ce cours contient » (page de garde)
\newcommand{\poptile}[3]{% #1 couleur #2 gros texte #3 légende
  \begin{tcolorbox}[enhanced,colback=#1,colframe=popink,boxrule=2pt,arc=9pt,shadow={2.5pt}{-2.5pt}{0pt}{popink},
      left=6pt,right=6pt,top=9pt,bottom=9pt,halign=center,valign=center,height=1.75cm,width=\linewidth,nobeforeafter]
    {\popdisplay\fontsize{12.5}{13}\selectfont\color{white}\MakeUppercase{#2}}\par\vspace{3pt}
    {\centering\popsemi\fontsize{7.8}{9.5}\selectfont\color{white}#3\par}
  \end{tcolorbox}}
\newcommand{\popcontenu}{%
  \par\noindent{\popxboldls\fontsize{7.5}{7.5}\selectfont\color{popmuted}CE COURS SIGNÉ CONTIENT}\par\vspace{7pt}
  \noindent\begin{minipage}[t]{0.235\linewidth}\poptile{popblue}{Cours}{complet, illustré, pas à pas}\end{minipage}\hfill
  \begin{minipage}[t]{0.235\linewidth}\poptile{poporange}{Méthodes}{et exercices résolus}\end{minipage}\hfill
  \begin{minipage}[t]{0.235\linewidth}\poptile{poppink}{Exercices}{type examen + corrigés}\end{minipage}\hfill
  \begin{minipage}[t]{0.235\linewidth}\poptile{popgreen}{Fiche bilan}{l'essentiel à retenir}\end{minipage}\par}

% Sommaire
\newtcolorbox{sommaire}{enhanced,colback=white,colframe=popink,boxrule=2.2pt,arc=10pt,
  drop shadow=popink,left=18pt,right=18pt,top=13pt,bottom=13pt,before skip=6pt,after skip=20pt,
  before upper={\poppill{Sommaire}{poppurple}{white}\par\vspace{9pt}\popsemi\fontsize{10.6}{14.5}\selectfont\color{popink}}}

% Signature de fin de cours — poussée tout en bas de la dernière page
\newcommand{\findecours}{%
  \par\vfill\nopagebreak
  \begin{center}\popsquiggle{poporange}\end{center}\vspace{4pt}
  \signatureonbuch
}
\AtBeginDocument{\def\llbracket{\mathopen{[\mkern-3.5mu[}}\def\rrbracket{\mathclose{]\mkern-3.5mu]}}} % intervalles d'entiers (glyphe ⟦ absent de la police)
% Glyphes absents de Fira Math : redéfinis à partir de glyphes disponibles
\AtBeginDocument{%
  \def\setminus{\mathbin{\backslash}}\let\smallsetminus\setminus
  \def\vdots{\mathord{\vbox{\baselineskip3.2pt\lineskiplimit0pt\kern4pt\hbox{.}\hbox{.}\hbox{.}}}}%
  \def\ddots{\mathinner{\mkern1mu\raise6.5pt\hbox{.}\mkern2mu\raise3.6pt\hbox{.}\mkern2mu\raise0.7pt\hbox{.}\mkern1mu}}%
  \def\triangle{\mathord{\scalebox{0.9}{$\symup{\Delta}$}}}%
  \def\nabla{\mathord{\raisebox{\depth}{\scalebox{1}[-1]{$\symup{\Delta}$}}}}%
  \def\checkmark{\popcheck{popdark}}%
  \def\star{\mathbin{\ast}}\def\bigstar{\mathord{\ast}}%
  \def\diamond{\mathbin{\scalebox{0.6}{$\Diamond$}}}%
  \def\bigcup{\mathop{\mathchoice{\raisebox{-0.3em}{\scalebox{1.7}{$\cup$}}}{\scalebox{1.25}{$\cup$}}{\cup}{\cup}}\displaylimits}%
  \def\bigcap{\mathop{\mathchoice{\raisebox{-0.3em}{\scalebox{1.7}{$\cap$}}}{\scalebox{1.25}{$\cap$}}{\cap}{\cap}}\displaylimits}%
  \def\lesseqqgtr{\mathrel{\lessgtr}}\def\bigoplus{\mathop{\oplus}\displaylimits}%
}
\providecommand{\ssi}{si et seulement si} % abréviation fréquente
\AtBeginDocument{\providecommand{\wideparen}{\overparen}} % arc de cercle

\begin{document}

\pagegarde{Citizenship and Human Rights}{Anglais}{1re Tech}{Anglais – classe de Première technique}{N04}{13 séances (fiche de progression harmonisée)}{This lesson develops learners' ability to discuss citizenship, work culture, civic responsibility, democracy, tolerance and basic freedoms in English. Through Cameroonian documents, dialogues and data, learners master key grammar tools (conjunctions, pronouns, prepositions, adverbs, phrasal verbs) and produce a composition on exercising democracy and promoting freedoms.
\vspace{4pt}
\begin{multicols}{2}\small
\begin{itemize}
\item À la fin de cette leçon, je sais échanger des informations en anglais sur la culture du travail et l'excellence professionnelle
\item À la fin de cette leçon, je sais utiliser correctement les conjonctions de coordination et de subordination (since, for, yet, although, in spite of, despite, however, nevertheless...)
\item À la fin de cette leçon, je sais employer le vocabulaire lié au travail, à la citoyenneté, à la démocratie et aux libertés fondamentales
\item À la fin de cette leçon, je sais comprendre un document audio sur la responsabilité civique individuelle
\item À la fin de cette leçon, je sais utiliser les pronoms réciproques et indéfinis (each, each other, one another)
\item À la fin de cette leçon, je sais utiliser les prépositions et locutions prépositionnelles de temps, de lieu et de manière
\end{itemize}\end{multicols}}

\begin{sommaire}
\begin{multicols}{2}
\begin{enumerate}
\item Work Culture and Excellence at Work (Speaking)
\item Promoting Individual Civic Responsibility (Listening)
\item Democracy, Tolerance and Basic Freedoms (Reading)
\item Comparative Adverbs, Adverbs of Negation and Phrasal Verbs
\item Writing: Compositions on Democracy, Tolerance and Freedoms
\item Integration, Evaluation and Remediation
\item Activité d'intégration
\item Méthodes et exercices résolus
\item Exercices
\item Corrigés des exercices
\item Fiche bilan
\end{enumerate}
\end{multicols}
\end{sommaire}

\begin{objectifs}
\begin{itemize}
\item À la fin de cette leçon, je sais échanger des informations en anglais sur la culture du travail et l'excellence professionnelle
\item À la fin de cette leçon, je sais utiliser correctement les conjonctions de coordination et de subordination (since, for, yet, although, in spite of, despite, however, nevertheless...)
\item À la fin de cette leçon, je sais employer le vocabulaire lié au travail, à la citoyenneté, à la démocratie et aux libertés fondamentales
\item À la fin de cette leçon, je sais comprendre un document audio sur la responsabilité civique individuelle
\item À la fin de cette leçon, je sais utiliser les pronoms réciproques et indéfinis (each, each other, one another)
\item À la fin de cette leçon, je sais utiliser les prépositions et locutions prépositionnelles de temps, de lieu et de manière
\item À la fin de cette leçon, je sais comparer les adverbes et employer les adverbes de négation (hardly, scarcely, never)
\item À la fin de cette leçon, je sais identifier et utiliser les phrasal verbs courants
\item À la fin de cette leçon, je sais rédiger une composition argumentée sur l'exercice de la démocratie, la tolérance et la promotion des libertés au Cameroun
\end{itemize}
\end{objectifs}

\begin{prerequis}
\begin{itemize}
\item Basic sentence structure in English (subject-verb-object) seen in Form 5 and Seconde
\item Simple present, present continuous and simple past tenses
\item Basic vocabulary about jobs, school and community life
\item Notions of rights and duties from Citizenship Education (Education à la citoyenneté) in lower classes
\item Reading and listening strategies: skimming and scanning
\end{itemize}
\end{prerequis}

\newpage

\coursec{Work Culture and Excellence at Work (Speaking)}

\textbf{Situation d'accroche.} In Yaoundé, a young technician named Amina works in a carpentry workshop in Nkolbisson. Every morning she arrives early, respects her colleagues, and votes during municipal elections, yet she notices that some workers are paid unfairly and that a neighbour was arrested for criticizing a local official. She asks herself: what does it really mean to be a good citizen in Cameroon? How can she promote excellence at work, defend her rights, and respect the rights of others? This lesson will help learners answer these questions in English.

\textbf{Problématique.} Before Amina can defend rights or exercise democracy, she must first master the language of the workplace: how do we describe good and bad attitudes at work? How do we exchange opinions politely in English? And how do we link our ideas with conjunctions to sound clear and professional? This first section answers these questions through speaking.

\courssub{Brainstorming: What Is a Work Culture?}

Look around you. In some workshops in Douala, workers arrive at 7:30 sharp, greet each other, wear their safety boots, and finish their tasks before leaving. In other places, workers come late, argue, and leave tools everywhere. The difference is not the machines or the money: the difference is the \cle{work culture}.

\begin{definition}[Work culture]
A \cle{work culture} is the set of \textbf{shared values, attitudes, and behaviours} that guide people at work. It answers the question: ``How do we do things here?'' A good work culture includes \cle{punctuality} (ponctualité), \cle{honesty} (honnêteté), \cle{teamwork} (travail d'équipe), \cle{respect} (respect) and \cle{diligence} (application, sérieux).
\end{definition}

\textbf{Brainstorming activity (oral).} In pairs, answer these questions in English:
\begin{enumerate}
\item What time do workers in your neighbourhood start work? Is it important to be on time? Why?
\item Name three behaviours of a good worker and three behaviours of a bad worker.
\item Would you like to work in a team where nobody respects each other? Why not?
\end{enumerate}

\courssub{Vocabulary: Words for Excellence at Work}

Here are the key words you must know and use. Read them aloud, then use each one in a sentence about a Cameroonian workplace.

\begin{center}
\begin{tabularx}{\linewidth}{|l|X|l|}\hline
\rowcolor{popblueL} \textbf{Word} & \textbf{Meaning (French gloss)} & \textbf{Example} \\ \hline
\cle{punctuality} & ponctualité ; being on time & Punctuality is the first rule in our workshop. \\ \hline
\cle{commitment} & engagement ; giving your best & Her commitment to the job impresses everybody. \\ \hline
\cle{integrity} & intégrité ; being honest and moral & A cashier must have integrity with money. \\ \hline
\cle{productivity} & productivité ; quantity and quality of work & Teamwork increases productivity on the farm. \\ \hline
\cle{excellence} & excellence ; doing better than average & We aim at excellence, not just ``good enough''. \\ \hline
\cle{professional ethics} & éthique professionnelle & Professional ethics forbid a nurse from neglecting a patient. \\ \hline
\cle{team spirit} & esprit d'équipe & Team spirit helped the Eneo team repair the line quickly. \\ \hline
\cle{accountability} & responsabilité (reddition de comptes) & Accountability means you accept the consequences of your acts. \\ \hline
\end{tabularx}
\end{center}

\begin{attention}[False friends and tricky words]
Do not confuse \emph{actually} (= in fact, en fait) with \emph{actuellement}. Do not say ``I am \emph{conscient}''; say ``I am \emph{hard-working}'' or ``I take my work seriously''. \emph{Diligence} means careful, steady effort --- it is not only ``speed''.
\end{attention}

\courssub{Good Habits vs Bad Habits at Work}

\begin{popfigure}
\begin{tikzpicture}[font=\small]
% Left column: good habits
\fill[popgreenL,rounded corners=4pt] (-0.2,-0.2) rectangle (6.4,6.2);
\node[anchor=north,font=\bfseries\color{popgreen!60!black}] at (3.1,6.0) {Good work habits};
\node[anchor=north west,text width=5.6cm,align=left] at (0.1,5.3) {
\textbullet\ Arriving on time\\
\textbullet\ Respecting colleagues\\
\textbullet\ Finishing tasks carefully\\
\textbullet\ Wearing safety equipment\\
\textbullet\ Helping teammates\\
\textbullet\ Reporting problems honestly};
% Clock pictogram
\draw[popgreen!60!black,thick] (5.35,4.75) circle (0.22);
\draw[popgreen!60!black,thick] (5.35,4.75) -- (5.35,4.93);
\draw[popgreen!60!black,thick] (5.35,4.75) -- (5.52,4.75);
% Right column: bad habits
\fill[poppinkL,rounded corners=4pt] (7.0,-0.2) rectangle (13.6,6.2);
\node[anchor=north,font=\bfseries\color{poppink!60!black}] at (10.3,6.0) {Bad work habits};
\node[anchor=north west,text width=5.6cm,align=left] at (7.3,5.3) {
\textbullet\ Coming late every day\\
\textbullet\ Insulting or ignoring colleagues\\
\textbullet\ Leaving work unfinished\\
\textbullet\ Refusing safety rules\\
\textbullet\ Working alone, hiding information\\
\textbullet\ Stealing tools or materials};
% Cross pictogram
\draw[poppink!60!black,thick] (12.5,4.6) -- (12.9,5.0);
\draw[poppink!60!black,thick] (12.9,4.6) -- (12.5,5.0);
\end{tikzpicture}
\legende{Two-column chart: good work habits (left, green) versus bad work habits (right, pink). Use it to build sentences: ``A good worker arrives on time, \textbf{but} a lazy worker comes late.''}
\end{popfigure}

\courssub{Exchanging Information: Dialogues and Role Plays}

Speaking about work means \textbf{giving your opinion}, then \textbf{agreeing or disagreeing politely}. Here are the tools.

\begin{aretenir}[Useful expressions for exchanging opinions]
\begin{itemize}
\item \textbf{Giving an opinion:} \emph{I think that... / In my view... / I believe that... / For me, ... / It seems to me that...}
\item \textbf{Agreeing:} \emph{I agree with you. / That's true. / Exactly! / You are right. / I share your point of view.}
\item \textbf{Disagreeing politely:} \emph{I see your point, but... / I'm afraid I don't agree. / That may be true; however, ... / I don't think so, because...}
\end{itemize}
\end{aretenir}

\begin{exemplebox}[Dialogue 1 --- In a carpentry workshop, Nkolbisson (Yaoundé)]
\textbf{Amina:} In my view, punctuality is the most important value in our workshop.\par
\textbf{Bobo:} I agree with you. When we arrive on time, we finish the furniture before the customers come.\par
\textbf{Amina:} Yes, and I also think that team spirit matters. Yesterday, I helped Clément with the heavy boards.\par
\textbf{Bobo:} That's true. However, some workers never help anybody. I'm afraid \textbf{it} is bad for productivity.\par
\textbf{Amina:} You are right. We should talk to them politely, because we are one team.
\end{exemplebox}

\begin{exemplebox}[Dialogue 2 --- In an office in Douala]
\textbf{Mrs Etoundi:} For me, integrity is essential for an accountant.\par
\textbf{Mr Kamga:} I see your point, but I believe that productivity is even more important.\par
\textbf{Mrs Etoundi:} That may be true; however, an accountant without integrity can destroy the whole company.\par
\textbf{Mr Kamga:} Exactly. So we need both: honest workers and hard-working ones.
\end{exemplebox}

\textbf{Role play (in pairs).} Student A is a farm supervisor in Bafoussam; Student B is a new worker. Discuss: working hours, respect, safety, and team spirit. Use at least two opinion expressions, one agreement and one polite disagreement.

\courssub{Grammar Focus 1: Coordinating Conjunctions}

Intuition first. Look at these two short sentences: \emph{Amina arrives early. Amina works hard.} English, like French, prefers to join them: \emph{Amina arrives early \textbf{and} works hard.} The word \emph{and} is a \cle{coordinating conjunction}: it links two elements of \textbf{equal importance}.

\begin{definition}[Coordinating conjunctions]
Coordinating conjunctions join two words, two phrases or two independent clauses of equal rank. Remember them with the word \textbf{FANBOYS}: \cle{for} (car), \cle{and} (et), \cle{nor} (ni), \cle{but} (mais), \cle{or} (ou), \cle{yet} (pourtant), \cle{so} (donc).
\end{definition}

\begin{center}
\begin{tabularx}{\linewidth}{|l|X|X|}\hline
\rowcolor{popblueL} \textbf{Conjunction} & \textbf{Use} & \textbf{Example} \\ \hline
\cle{and} & addition & She is punctual \textbf{and} diligent. \\ \hline
\cle{but} & contrast & He is tired, \textbf{but} he finishes his task. \\ \hline
\cle{or} & choice & You can sweep the floor \textbf{or} arrange the tools. \\ \hline
\cle{so} & result & The machine broke, \textbf{so} we stopped work. \\ \hline
\cle{for} & reason (formal) & She wears gloves, \textbf{for} the chemicals are dangerous. \\ \hline
\cle{yet} & surprising contrast & The salary is low, \textbf{yet} he stays committed. \\ \hline
\end{tabularx}
\end{center}

\begin{attention}[Comma before the conjunction]
When a coordinating conjunction joins two \textbf{full clauses} (each with its own subject), put a comma before it: \emph{The machine broke, so we stopped.} Do not put a comma when it joins two words or phrases: \emph{punctual and diligent}.
\end{attention}

\courssub{Grammar Focus 2: Subordinating Conjunctions of Concession and Reason}

Intuition. Sometimes one idea is \textbf{more important} than the other. \emph{Although Amina is tired, she finishes her work.} The main message is ``she finishes her work''; the clause ``although Amina is tired'' only adds a difficulty. \emph{Although} is a \cle{subordinating conjunction}: it introduces a \textbf{subordinate clause} that depends on the main clause.

\begin{definition}[Subordinating conjunctions: concession and reason]
\begin{itemize}
\item \textbf{Concession} (concession): \cle{although} / \cle{though} (bien que) + clause ; \cle{in spite of} / \cle{despite} (malgré) + noun or \emph{-ing} form.
\item \textbf{Reason} (cause): \cle{because} (parce que), \cle{since} (puisque), \cle{as} (comme, puisque) + clause.
\end{itemize}
\end{definition}

\begin{popfigure}
\begin{tikzpicture}[font=\small, node distance=0.4cm]
\node[draw=poppurple,fill=poppurpleL,rounded corners=3pt,minimum width=5.6cm,minimum height=1cm] (sub) {\textbf{Although} he is tired,};
\node[draw=popblue,fill=popblueL,rounded corners=3pt,minimum width=5.6cm,minimum height=1cm,right=0.8cm of sub] (main) {he works carefully.};
\draw[-{Stealth},thick,popdark] (sub) -- (main);
\node[below=0.15cm of sub,text=poppurple!70!black] {subordinate clause (concession)};
\node[below=0.15cm of main,text=popblue!70!black] {main clause (the real message)};
\node[above=0.25cm of sub,font=\itshape] {subordinating conjunction + subject + verb};
\node[above=0.25cm of main,font=\itshape] {subject + verb};
\end{tikzpicture}
\legende{Structure of a sentence with a subordinating conjunction: the subordinate clause (introduced by \emph{although}) depends on the main clause. A comma separates the two clauses when the subordinate clause comes first.}
\end{popfigure}

\begin{methode}[Building a contrast sentence]
To express a contrast between two facts, follow these steps:
\begin{enumerate}
\item Identify the two facts: \emph{He is tired} / \emph{He works carefully}.
\item Choose the structure: \textbf{Although + clause, + main clause.} $\rightarrow$ \emph{Although he is tired, he works carefully.}
\item Or invert with an adverb: \textbf{main clause; however, + clause.} $\rightarrow$ \emph{He is tired; however, he works carefully.} (Also possible as two sentences: \emph{He is tired. However, he works carefully.})
\item With a noun or \emph{-ing} form, use \textbf{despite / in spite of}: \emph{Despite his tiredness, he works carefully.} / \emph{In spite of being tired, he works carefully.}
\end{enumerate}
\end{methode}

\begin{attention}[The classic trap: although + but]
Never use \emph{but} with \emph{although} in the same sentence. One connector is enough!\par
\emph{Although he is tired, but he works.} \; \textbf{✗}\par
\emph{Although he is tired, he works.} \; \textbf{✓} \quad or \quad \emph{He is tired, but he works.} \; \textbf{✓}
\end{attention}

\begin{exemplebox}[Contrast connectors: however, nevertheless, still]
These connectors also express contrast, but they are \textbf{adverbs}, not conjunctions: they usually start a new sentence or follow a semicolon, with a comma after them.
\begin{itemize}
\item \emph{The work is hard. \textbf{However}, Amina never complains.}
\item \emph{The salary is low; \textbf{nevertheless}, the team stays committed.}
\item \emph{It was raining. \textbf{Still}, the Eneo team repaired the line before night.}
\end{itemize}
Position: beginning of the sentence + comma, or between commas in the middle: \emph{Amina, \textbf{however}, never complains.}
\end{exemplebox}

\courssub{Document Study: Workers in a Douala Factory}

Read the following short text, then answer the questions orally.

\begin{exemplebox}[Text: ``A Morning at the Sawa Factory, Douala'']
At 7:15 every morning, the workers of a soap factory in Bonabéri, Douala, arrive at the gate. Most of them are punctual, and they greet the watchman politely. Inside, the team leader, Mrs Ngo Bell, checks that everybody wears gloves and boots. Although the machines are old, the team maintains them carefully, so accidents are rare. One worker, Mr Fotso, sometimes arrives late; however, he always finishes his quota of boxes. ``Productivity is important,'' says Mrs Ngo Bell, ``but integrity and team spirit are even more important, because a factory is like a family.''
\end{exemplebox}

\textbf{Comprehension questions (oral work):}
\begin{enumerate}
\item What time do the workers arrive? Are they punctual?
\item Find two good work attitudes in the text and one bad habit.
\item Find one coordinating conjunction and one subordinating conjunction in the text. Explain the link each one expresses.
\item Do you agree with Mrs Ngo Bell that ``a factory is like a family''? Give your opinion with \emph{In my view...} or \emph{I think that...}.
\end{enumerate}

\begin{exoresolu}[Analysing a contrast sentence]
\textbf{Statement.} Analyse the structure of this sentence: \emph{``Despite working 8 hours a day, some informal workers in Yaoundé earn less than 2,000 FCFA daily.''} Identify the connector, its type, and the contrast expressed. Then rewrite the sentence with \emph{although}.
\tcblower
\textbf{Solution.}
\begin{enumerate}
\item The connector is \emph{despite} (= malgré). It is a \textbf{preposition of concession}, followed by an \textbf{\emph{-ing} form} (\emph{working}), not by a full clause.
\item The contrast: we expect workers who work 8 hours a day to earn a decent salary; in reality, they earn less than 2,000 FCFA. The reality contradicts our expectation.
\item Rewriting with \emph{although} + clause: \emph{``\textbf{Although} some informal workers in Yaoundé \textbf{work} 8 hours a day, they earn less than 2,000 FCFA daily.''}
\item Check the trap: we do \textbf{not} say \emph{Although they work 8 hours, but they earn...} ✗
\end{enumerate}
\end{exoresolu}

\begin{exercice}[Practice --- link the ideas]
Complete each sentence with the correct connector: \emph{and, but, so, although, despite, however, because, yet}.
\begin{enumerate}
\item The workers are tired, \dots\ they continue to smile at the customers.
\item \dots\ the rain, the farmers of Bafoussam worked all morning.
\item She forgot her boots, \dots\ she could not enter the workshop.
\item \dots\ he is young, he is the most diligent apprentice in the garage.
\item The office \textbf{closed} at 4 p.m.; \dots\ the accountant \textbf{stayed} until 6 p.m. to finish the report.
\item Integrity \dots\ punctuality are both necessary for a good cashier.
\end{enumerate}
\end{exercice}

\begin{corrige}[Exercice 1]
\begin{enumerate}
\item \textbf{but} (or \emph{yet}) --- contrast between tiredness and smiling.
\item \textbf{Despite} --- followed by a noun (\emph{the rain}); \emph{although} would need a clause (\emph{Although it rained...}).
\item \textbf{so} --- result: forgetting the boots causes the impossibility to enter.
\item \textbf{Although} --- concession at the beginning of the sentence, comma after the clause.
\item \textbf{however} --- contrast adverb after a semicolon, followed by a comma.
\item \textbf{and} --- simple addition of two qualities.
\end{enumerate}
\end{corrige}

\begin{savaistu}[Labour Day in Cameroon]
Cameroon celebrates \textbf{Labour Day} (la Fête du travail) on the 1st of May, like many countries in the world. On that day, workers' unions march in Yaoundé, Douala and the other regional capitals, and the authorities officially promote \cle{excellence at work}, better salaries and respect for workers' rights. It is a perfect occasion to use the vocabulary of this lesson: punctuality, commitment, integrity and team spirit!
\end{savaistu}

\begin{aretenir}[What you must remember]
\begin{itemize}
\item A \cle{work culture} = shared values, attitudes and behaviours at work: punctuality, honesty, teamwork, respect, diligence.
\item Key vocabulary: punctuality, commitment, integrity, productivity, excellence, professional ethics, team spirit, accountability.
\item To exchange opinions: \emph{I think that... / In my view... / I agree. / I see your point, but...}
\item Coordinating conjunctions (FANBOYS) link equal elements: \emph{and, but, or, so, for, yet}.
\item Subordinating conjunctions introduce a dependent clause: \emph{although, though, because, since}; \emph{despite / in spite of} + noun or \emph{-ing}.
\item Contrast adverbs \emph{however, nevertheless, still} start a sentence or follow a semicolon, with a comma.
\item Golden rule: never \emph{although} + \emph{but} in the same sentence.
\end{itemize}
\end{aretenir}

\coursec{Promoting Individual Civic Responsibility (Listening)}

In the previous section, we explored how a strong \cle{work culture} and the pursuit of \cle{excellence at work} shape professional success. We practised expressing opinions, agreeing and disagreeing politely, and using conjunctions to link ideas logically. Now, we widen our perspective from the workplace to the wider community. A healthy work culture cannot exist in isolation; it thrives when citizens fulfil their \cle{civic duties}. In this section, we will listen to a radio talk about individual civic responsibility, learn key vocabulary, sharpen our listening strategies, and master two important grammar points: reciprocal pronouns (\emph{each other} / \emph{one another}) and the indefinite pronoun \emph{each}.

\courssub{Pre-listening: Predicting Content from Titles and Pictures}

Before listening, effective learners activate their background knowledge. Look at the three pictures below and the title of the radio programme: \emph{``Building a Better Neighbourhood: Your Duty, My Duty''}.

\begin{center}
\begin{tabularx}{\linewidth}{|X|X|X|}\hline
\textbf{Picture A} & \textbf{Picture B} & \textbf{Picture C} \\ \hline
A citizen dropping a ballot paper into a transparent ballot box. \newline \emph{(Voting / \textit{Voter})} & A person handing a tax form to an officer at a counter. \newline \emph{(Paying taxes / \textit{Payer des impôts})} & Young volunteers sweeping streets and planting trees. \newline \emph{(Keeping the environment clean / \textit{Garder l'environnement propre})} \\ \hline
\end{tabularx}
\end{center}

\begin{exercice}{Predict and Match}
Match each picture (A, B, C) with the civic duty it illustrates. Then, with a partner, answer: \emph{``Why is each duty important for the community?''}
\begin{enumerate}
\item Picture A shows \dots
\item Picture B shows \dots
\item Picture C shows \dots
\end{enumerate}
\end{exercice}

\begin{corrige}{Exercice 1}
\begin{enumerate}
\item Picture A shows \textbf{voting} (participating in elections).
\item Picture B shows \textbf{paying taxes} (funding public services).
\item Picture C shows \textbf{keeping the environment clean / community service} (volunteer work).
\end{enumerate}
\textit{Sample reasons:} Voting chooses leaders; taxes pay for hospitals, roads, schools; a clean environment prevents disease and makes life pleasant.
\end{corrige}

\begin{definition}[Civic Responsibility]
\textbf{Civic responsibility} (la responsabilité civique) refers to the duties every citizen owes to the community and the State. These duties include:
\begin{itemize}
\item \cle{Voting} (le vote) in elections to choose representatives.
\item \cle{Paying taxes} (payer les impôts) to fund public services (health, education, infrastructure).
\item \cle{Respecting the law} (respecter la loi) and being a \cle{law-abiding citizen} (citoyen respectueux des lois).
\item Participating in \cle{community service} (service communautaire) or \cle{volunteer} (bénévole) work.
\item Showing \cle{patriotism} (patriotisme) by promoting national unity and protecting public property.
\end{itemize}
\end{definition}

\begin{aretenir}[Key Vocabulary: Civic Responsibility]
\begin{center}
\begin{tabularx}{\linewidth}{|l|X|l|}\hline
\rowcolor{popblueL}
\textbf{English} & \textbf{Meaning / Contexte} & \textbf{Français} \\ \hline
civic responsibility & duty to community/State & responsabilité civique \\
duty & moral/legal obligation & devoir \\
tax (n.) & compulsory payment to government & impôt \\
vote (v./n.) & express choice in election & voter / vote \\
community service & unpaid work for community & service communautaire \\
law-abiding citizen & person who obeys laws & citoyen respectueux des lois \\
patriotism & love for one's country & patriotisme \\
volunteer (n./v.) & person who works freely / to offer freely & bénévole / faire du bénévolat \\
clean-up campaign & organised cleaning event & campagne de nettoyage \\
neighbourhood (UK) / neighborhood (US) & local area / quarter & quartier \\ \hline
\end{tabularx}
\end{center}
\end{aretenir}

\courssub{Listening: A Radio Talk on Civic Responsibility}

\emph{Context:} You are going to listen to a 3-minute radio talk broadcast on CRTV (Cameroon Radio Television) during the morning programme ``Community Watch''. The host, Mr. Atangana, discusses the monthly \emph{``Clean-Up Saturday''} campaign in Yaoundé neighbourhoods.

\begin{methode}[How to Succeed in a Listening Task]
Follow these steps before, during, and after the audio:
\begin{enumerate}
\item \textbf{Read the questions first.} Identify the \cle{keywords} (who, what, when, where, why, how) you need to listen for.
\item \textbf{Predict the content.} Use the title, pictures, and vocabulary to guess what the speaker will say.
\item \textbf{First listening (Gist).} Listen for the \emph{main idea}. Do not try to understand every word. Ask: ``What is the general topic?''
\item \textbf{Second listening (Specific Information).} Focus on the keywords. Write short notes, \cle{abbreviations} (e.g., \emph{gov't} for government, \emph{camp'n} for campaign, \emph{vol's} for volunteers), or symbols.
\item \textbf{Third listening (Check).} Verify your answers. Fill in any gaps.
\item \textbf{Post-listening.} Summarise the talk in 2--3 sentences using your notes.
\end{enumerate}
\end{methode}

\begin{experience}[Listening Script -- ``Community Watch: Clean-Up Saturday'']
\textit{(Teacher reads aloud or plays audio. Students do not read the script during the first two listenings.)}

\textbf{Host (Mr. Atangana):} ``Good morning, dear listeners. Welcome to \emph{Community Watch}. I am your host, Atangana. Today, we talk about a duty that belongs to \emph{each} of us: keeping our neighbourhoods clean.

Last Saturday, the \emph{``Clean-Up Saturday''} campaign took place in several quarters of Yaoundé -- Mvog-Ada, Essos, and Nlongkak. The campaign starts \emph{since} the first Saturday of every month. Residents wake up early, \emph{although} it is a holiday. They gather at the quarter head's office \emph{at 7 a.m.} The delegate provides gloves, brooms, and bin bags.

I spoke to Madame Ngo, a volunteer in Mvog-Ada. She said: ``We do this \emph{for} our children. \emph{Each} citizen has a duty to protect the environment. \emph{Each of us} brings a tool. We work together; we help \emph{each other}. The young men carry heavy bins; the women sweep. Nobody watches the other do the work alone.''

The Divisional Officer praised the turnout. He noted that \emph{despite} the rain last month, people still came. \emph{However}, he warned that \emph{some} citizens still dump waste in streams. He urged \emph{each} household to sort waste: plastics, organic, glass.

Remember, listeners: a clean city is not the mayor's job alone. It is \emph{our} shared responsibility. When we clean \emph{one another}'s streets, we build a healthier city for \emph{each} family. Join us next month. This has been \emph{Community Watch}. Stay civic, stay clean.''
\end{experience}

\begin{exercice}{Listening for Gist (First Listening)}
Tick the correct answer.
\begin{enumerate}
\item What is the main topic of the radio talk?
\begin{itemize}
\item[a] The history of Yaoundé.
\item[b] A monthly clean-up campaign.
\item[c] Paying taxes on time.
\item[d] Voting in local elections.
\end{itemize}
\item Who is the host?
\begin{itemize}
\item[a] Madame Ngo.
\item[b] The Divisional Officer.
\item[c] Mr. Atangana.
\item[d] The quarter head.
\end{itemize}
\end{enumerate}
\end{exercice}

\begin{corrige}{Exercice 2}
\begin{enumerate}
\item \textbf{b} (A monthly clean-up campaign).
\item \textbf{c} (Mr. Atangana).
\end{enumerate}
\end{corrige}

\begin{exercice}{Listening for Specific Information (Second Listening)}
Complete the table below using information from the talk. Use \emph{abbreviations} and \emph{keywords} only.
\end{exercice}

\begin{popfigure}
\begin{center}
\begin{tabularx}{\linewidth}{|l|X|X|}\hline
\rowcolor{popblueL}
\textbf{Civic Duty / Action} & \textbf{Why it Matters (Reason)} & \textbf{Example in My Community} \\ \hline
Clean-Up Saturday campaign & & \\ \hline
Sorting waste (plastics, organic, glass) & & \\ \hline
Volunteering (gloves, brooms, bins) & & \\ \hline
Not dumping waste in streams & & \\ \hline
\end{tabularx}
\end{center}
\legende{Tableau de prise de notes : Devoir civique / Pourquoi c'est important / Exemple dans ma communauté}
\end{popfigure}

\begin{corrige}{Exercice 3}
\begin{center}
\begin{tabularx}{\linewidth}{|l|X|X|}\hline
\rowcolor{popblueL}
\textbf{Civic Duty / Action} & \textbf{Why it Matters (Reason)} & \textbf{Example in My Community} \\ \hline
Clean-Up Saturday campaign & For our children; protect environment; healthier city & Monthly in Mvog-Ada, Essos, Nlongkak \\ \hline
Sorting waste (plastics, organic, glass) & Easier recycling; less pollution & Household level urged by Div. Officer \\ \hline
Volunteering (gloves, brooms, bins) & Shared work; nobody works alone & Young men carry bins; women sweep \\ \hline
Not dumping waste in streams & Prevents water pollution; health & Still a problem noted by officer \\ \hline
\end{tabularx}
\end{center}
\end{corrige}

\begin{exercice}{Note-taking Practice (Third Listening)}
Listen again and fill in the missing conjunctions or pronouns in these extracts.
\begin{enumerate}
\item The campaign takes place \_\_\_\_\_\_\_\_ the first Saturday of every month. (\emph{since / for / although})
\item Residents gather \_\_\_\_\_\_\_\_ it is a holiday. (\emph{since / for / although})
\item \_\_\_\_\_\_\_\_ citizen has a duty to protect the environment. (\emph{Each / Each of / Every})
\item \_\_\_\_\_\_\_\_ us brings a tool. (\emph{Each / Each of / All})
\item We help \_\_\_\_\_\_\_\_. (\emph{each other / one another / ourselves})
\item When we clean \_\_\_\_\_\_\_\_ streets, we build a healthier city. (\emph{each other's / one another's / our})
\end{enumerate}
\end{exercice}

\begin{corrige}{Exercice 4}
\begin{enumerate}
\item \textbf{since} (point de départ dans le temps)
\item \textbf{although} (concession)
\item \textbf{Each} (chaque + singulier)
\item \textbf{Each of} (chacun de + pronom/pluriel)
\item \textbf{each other} (réciproque : deux groupes -- jeunes hommes / femmes -- mais ici sens collectif, \emph{one another} aussi possible)
\item \textbf{one another's} (réciproque possessif : plusieurs personnes/rues)
\end{enumerate}
\end{corrige}

\courssub{Post-listening Discussion}

Work in groups of three. Discuss the following questions. Use the vocabulary from the \textbf{aretenir} box.

\begin{enumerate}
\item What civic duties do \textbf{you} perform at school? (Examples: cleaning the classroom, respecting the school rules, helping a classmate, participating in the school garden.)
\item What civic duties do you see in your \emph{quarter} (neighbourhood)? (Examples: \emph{Clean-Up Saturday}, neighbourhood watch, contributing to the water well repair, respecting traffic signs.)
\item Madame Ngo says: ``We do this \emph{for} our children.'' What does she mean? Give a concrete example.
\item The Divisional Officer says some citizens \emph{still} dump waste in streams \emph{despite} the campaigns. Why do people do this? What can change this behaviour?
\end{enumerate}

\begin{exemplebox}[Sample Discussion Starter]
\textbf{Student A:} ``At school, I am a \emph{class monitor}. My \emph{duty} is to check if the classroom is clean before lessons. It is a \emph{civic responsibility} because a clean class helps \emph{each} student learn better.''
\textbf{Student B:} ``In my quarter, we have a \emph{neighbourhood watch}. Volunteers patrol at night. They help \emph{each other} to keep the area safe.''
\textbf{Student C:} ``Madame Ngo means we protect the environment \emph{for} the future generation. If we pollute now, our children will suffer.''
\end{exemplebox}

\courssub{Grammar Focus 1: Reciprocal Pronouns -- \emph{each other} vs \emph{one another}}

Reciprocal pronouns express a \cle{mutual action} or relationship: Person A does something to Person B, \emph{and} Person B does the same to Person A.

\begin{propriete}[Reciprocal Pronouns: Traditional Rule]
\begin{itemize}
\item \textbf{\emph{each other}} is used when \textbf{two} people (or things) are involved.
\item \textbf{\emph{one another}} is used when \textbf{more than two} people (or things) are involved.
\end{itemize}
\textit{Note:} In modern English, the distinction is often blurred; \emph{one another} is slightly more formal. For your exams, \textbf{apply the traditional rule}.
\end{propriete}

\begin{popfigure}
\begin{center}
\begin{tikzpicture}[
  person/.style={circle, draw=popdark, fill=popblueL, minimum width=1.8cm, font=\small},
  arrow/.style={->, >=Stealth, thick, draw=poporange},
  doublearrow/.style={<->, >=Stealth, thick, draw=popgreen},
  labelstyle/.style={font=\footnotesize, text=popdark}
]
% Two people: each other
\node[person] (A1) at (0,0) {Person A};
\node[person] (B1) at (3,0) {Person B};
\draw[doublearrow] (A1.east) -- (B1.west) node[midway, above, labelstyle] {\textbf{each other}};
\node[labelstyle] at (1.5, -1.3) {\textbf{Two people $\rightarrow$ each other}};
% Three+ people: one another
\node[person] (A2) at (6,0) {Person A};
\node[person] (B2) at (9,0) {Person B};
\node[person] (C2) at (7.5, 2.2) {Person C};
\draw[arrow, bend left=20] (A2) to node[midway, above left, labelstyle] {} (B2);
\draw[arrow, bend left=20] (B2) to node[midway, above right, labelstyle] {} (C2);
\draw[arrow, bend left=20] (C2) to node[midway, left, labelstyle] {} (A2);
\draw[arrow, bend right=20] (B2) to node[midway, below left, labelstyle] {} (A2);
\draw[arrow, bend right=20] (C2) to node[midway, below right, labelstyle] {} (B2);
\draw[arrow, bend right=20] (A2) to node[midway, right, labelstyle] {} (C2);
\node[labelstyle] at (7.5, -1.3) {\textbf{Three+ people $\rightarrow$ one another}};
\end{tikzpicture}
\end{center}
\legende{Schéma des pronoms réciproques : \emph{each other} (deux personnes) vs \emph{one another} (plus de deux)}
\end{popfigure}

\begin{exemplebox}[Examples from the Listening Text]
\begin{enumerate}
\item ``We work together; we help \textbf{each other}.'' \textit{(Context: two groups -- young men \emph{and} women -- helping mutually.)}
\item ``When we clean \textbf{one another}'s streets, we build a healthier city.'' \textit{(Context: many residents, many streets -- more than two entities.)}
\end{enumerate}
\end{exemplebox}

\begin{exercice}{Reciprocal Pronouns: Choose the Correct Form}
Complete the sentences with \emph{each other} or \emph{one another}.
\begin{enumerate}
\item The two candidates shook hands and congratulated \_\_\_\_\_\_\_\_ after the debate.
\item In a strong community, neighbours look after \_\_\_\_\_\_\_\_.
\item The five members of the committee sent emails to \_\_\_\_\_\_\_\_ to fix the meeting time.
\item My brother and I often lend \_\_\_\_\_\_\_\_ books.
\item During the clean-up, the volunteers passed the heavy bins to \_\_\_\_\_\_\_\_.
\end{enumerate}
\end{exercice}

\begin{corrige}{Exercice 5}
\begin{enumerate}
\item \textbf{each other} (two candidates).
\item \textbf{one another} (neighbours = many people).
\item \textbf{one another} (five members > 2).
\item \textbf{each other} (two people: my brother and I).
\item \textbf{one another} (volunteers = group > 2).
\end{enumerate}
\end{corrige}

\courssub{Grammar Focus 2: Indefinite Pronoun \emph{each} and Agreement}

\emph{Each} refers to \cle{every individual member} of a group \emph{separately}. It is \textbf{always singular}.

\begin{propriete}[\emph{each} + Singular Agreement]
\begin{itemize}
\item \textbf{\emph{each} + singular noun} $\rightarrow$ \textbf{singular verb}.
  \textit{Example:} \emph{Each citizen \textbf{has} a duty.} (NOT \emph{have})
\item \textbf{\emph{each of} + plural noun/pronoun} $\rightarrow$ \textbf{singular verb}.
  \textit{Example:} \emph{Each of them \textbf{brings} a tool.} (NOT \emph{bring})
\item \textbf{\emph{each}} as pronoun (subject) $\rightarrow$ \textbf{singular verb}.
  \textit{Example:} \emph{Each \textbf{is} responsible.}
\end{itemize}
\end{propriete}

\begin{attention}[Common Error: \emph{each} + Plural Noun/Verb]
\begin{itemize}
\item \textbf{Wrong:} \emph{Each citizens have a duty.} $\times$
\item \textbf{Wrong:} \emph{Each of the students are here.} $\times$
\item \textbf{Correct:} \emph{Each citizen \textbf{has} a duty.} \checkmark
\item \textbf{Correct:} \emph{Each of the students \textbf{is} here.} \checkmark
\end{itemize}
\textbf{Reason:} \emph{Each} isolates individuals $\Rightarrow$ singular concept.
\end{attention}

\begin{exoresolu}{Sentence Transformation: \emph{each} / \emph{each of}}
\textbf{Task:} Rewrite the sentences using \emph{each} or \emph{each of} + singular verb.
\begin{enumerate}
\item All citizens have a duty to vote. $\rightarrow$ \emph{Each citizen \dots}
\item All the volunteers brought a broom. $\rightarrow$ \emph{Each of the volunteers \dots}
\item Every student must respect the rules. $\rightarrow$ \emph{Each student \dots}
\end{enumerate}

\textbf{Solution:}
\begin{enumerate}
\item \emph{Each citizen \textbf{has} a duty to vote.} (\emph{citizen} singular, \emph{has} singular)
\item \emph{Each of the volunteers \textbf{brought} a broom.} (\emph{each of} + plural noun, verb singular)
\item \emph{Each student \textbf{must respect} the rules.} (\emph{student} singular, modal \emph{must} + base form)
\end{enumerate}
\end{exoresolu}

\begin{exercice}{Gap-fill: \emph{each} / \emph{each of} + Verb Agreement}
Complete with the correct form of the verb in brackets.
\begin{enumerate}
\item Each household \_\_\_\_\_\_\_\_ (sort) its waste into three bins.
\item Each of the neighbourhoods \_\_\_\_\_\_\_\_ (organise) its own clean-up day.
\item Each volunteer \_\_\_\_\_\_\_\_ (receive) gloves and a bag.
\item Each of us \_\_\_\_\_\_\_\_ (be) responsible for our street.
\item Each tax payer \_\_\_\_\_\_\_\_ (contribute) to the national budget.
\end{enumerate}
\end{exercice}

\begin{corrige}{Exercice 6}
\begin{enumerate}
\item \textbf{sorts} (Each household = singular)
\item \textbf{organises} (Each of the neighbourhoods = singular)
\item \textbf{receives} (Each volunteer = singular)
\item \textbf{is} (Each of us = singular)
\item \textbf{contributes} (Each tax payer = singular)
\end{enumerate}
\end{corrige}

\courssub{Integrated Practice: Civic Context Sentences}

Combine the vocabulary, reciprocal pronouns, and \emph{each} structures.

\begin{exercice}{Sentence Transformation}
Rewrite the sentences as indicated.
\begin{enumerate}
\item The two delegates help one another. (Change to \emph{each other})
\item Every citizen must pay taxes. (Start with \emph{Each citizen \dots})
\item The volunteers clean the streets together. (Use \emph{one another}'s)
\item All the quarter heads attended the meeting. (Start with \emph{Each of the quarter heads \dots})
\item We must respect each other's rights. (Use \emph{one another}'s -- formal/more than two)
\end{enumerate}
\end{exercice}

\begin{corrige}{Exercice 7}
\begin{enumerate}
\item The two delegates help \textbf{each other}.
\item \textbf{Each citizen must pay taxes.} (or \emph{Each citizen has to pay taxes.})
\item The volunteers clean \textbf{one another}'s streets.
\item \textbf{Each of the quarter heads attended the meeting.}
\item We must respect \textbf{one another}'s rights.
\end{enumerate}
\end{corrige}

\begin{savaistu}[Clean-Up Days in Cameroon: Collective Civic Action]
In Cameroon, the \emph{``Clean-Up Saturday''} (\emph{Samedi de salubrité}) is a powerful example of \cle{civic responsibility} in action. Instituted by presidential decree in the 1990s and reinforced by local authorities, it takes place on the \textbf{first Saturday of every month} (usually 6 a.m. -- 10 a.m.).
\begin{itemize}
\item \textbf{Yaoundé:} The Urban Community organises campaigns in all 7 subdivisions (Yaoundé I--VII). The Governor and Divisional Officers often join residents in quarters like Mvog-Ada, Essos, Nlongkak, and Melen.
\item \textbf{Douala:} The ``\emph{Keep Douala Clean}'' (\emph{Garder Douala Propre}) operations mobilise thousands of volunteers, city council workers, and military engineering units to clear drains, collect plastic waste, and unblock waterways before the rainy season.
\item \textbf{Legal basis:} Law No. 98/015 of 14 July 1998 on the organisation of local councils empowers mayors to enforce sanitation by-laws. Citizens who refuse to participate without valid reason can be fined.
\item \textbf{Impact:} These campaigns reduce flood risks, limit cholera/malaria outbreaks, and foster \emph{patriotism} and \emph{community spirit}. They illustrate \emph{one another} in action: \emph{each} person contributes, and \emph{all} benefit.
\end{itemize}
\textbf{Your turn:} Has your quarter organised a clean-up recently? What was your role? Write a short paragraph (50--60 words) using \emph{each}, \emph{each other} / \emph{one another}, and at least three vocabulary words from the \textbf{aretenir} box.
\end{savaistu}

\begin{exemplebox}[Sample Paragraph]
``Last month, my quarter organised a \emph{clean-up campaign}. \emph{Each} resident \emph{was} asked to participate. \emph{Each of us} brought a broom or a rake. We worked hard and helped \emph{one another} to clear the blocked drainage. The quarter head thanked \emph{each} volunteer. It showed true \emph{patriotism} and \emph{civic responsibility}.''
\end{exemplebox}

\begin{aretenir}[Summary: Key Points to Master]
\begin{itemize}
\item \textbf{Listening strategy:} Read questions $\rightarrow$ Predict $\rightarrow$ Listen for gist $\rightarrow$ Listen for details $\rightarrow$ Check.
\item \textbf{Vocabulary:} \emph{civic responsibility, duty, tax, vote, community service, law-abiding citizen, patriotism, volunteer, clean-up campaign}.
\item \textbf{Grammar -- Reciprocal pronouns:} \emph{each other} (2 people), \emph{one another} ($>$2 people).
\item \textbf{Grammar -- \emph{each}:} Always singular verb. \emph{Each citizen has...} / \emph{Each of them has...}
\item \textbf{Writing:} Use these structures to describe civic actions in your community.
\end{itemize}
\end{aretenir}

\coursec{Democracy, Tolerance and Basic Freedoms (Reading)}

In the previous section, we listened to texts about \cle{promoting individual civic responsibility}: we learnt that being a good citizen means voting, paying taxes, keeping our environment clean and respecting the law. But a responsible citizen also lives in a \cle{democracy}, a country where people can choose their leaders and express their opinions freely. In this section, we will read texts about exercising democracy, tolerance and basic freedoms, and we will master the prepositions of time, place and manner.

\courssub{Activating Prior Knowledge: What is Democracy?}

\textbf{Think and discuss.} Before reading, answer these questions orally with your neighbour:

\begin{enumerate}
\item Who is the Head of State of Cameroon? How was he chosen?
\item What types of elections do you know in Cameroon?
\item Have you ever seen a voter's card? What is it used for?
\end{enumerate}

In Cameroon, citizens elect their leaders through several types of \cle{elections}:

\begin{itemize}
\item the \textbf{presidential election} (every seven years), to choose the President of the Republic;
\item the \textbf{legislative elections}, to choose the members of the National Assembly (deputies);
\item the \textbf{municipal elections}, to choose mayors and municipal councillors;
\item the \textbf{senatorial elections}, to choose senators.
\end{itemize}

\begin{definition}[Democracy]
\cle{Democracy} is a system of government in which citizens choose their leaders through free and fair elections. In a democracy, people enjoy \cle{basic freedoms} such as freedom of speech and freedom of the press, and all citizens are equal before the law.
\end{definition}

\begin{definition}[Tolerance]
\cle{Tolerance} means accepting and respecting people whose opinions, cultures, languages or beliefs are different from ours. Tolerance does not mean agreeing with everybody; it means living together peacefully in spite of our differences.
\end{definition}

\courssub{Vocabulary: Democracy, Elections and Freedoms}

Learn these key words before reading the text.

\begin{center}
\begin{tabularx}{\linewidth}{|l|X|}\hline
\rowcolor{popblueL} \textbf{Word / Expression} & \textbf{Meaning (with French gloss)} \\ \hline
democracy & system where citizens choose their leaders (\emph{démocratie}) \\ \hline
election & organised vote to choose a leader (\emph{élection}) \\ \hline
ballot & the paper or act of voting; ballot box = \emph{urne} \\ \hline
voter registration & putting your name on the electoral list (\emph{inscription sur les listes électorales}) \\ \hline
polling station & place where people vote (\emph{bureau de vote}) \\ \hline
freedom of speech & right to express your opinion (\emph{liberté d'expression}) \\ \hline
freedom of the press & right of journalists to inform freely (\emph{liberté de la presse}) \\ \hline
tolerance & respecting differences (\emph{tolérance}) \\ \hline
human rights & basic rights of every person (\emph{droits de l'homme}) \\ \hline
discrimination & unfair treatment because of origin, sex or religion (\emph{discrimination}) \\ \hline
rule of law & the principle that everyone obeys the law (\emph{État de droit}) \\ \hline
\end{tabularx}
\end{center}

\courssub{Reading Text: ELECAM Launches a Voter Registration Campaign}

\textbf{Pre-reading question.} Look at the title. What do you think the article is about?

\begin{exemplebox}[Reading Text: ``Every Vote Counts'' (adapted newspaper article)]
\textbf{Yaoundé, 5 January.} -- Elections Cameroon (ELECAM) has launched a national campaign to encourage young people to register on the electoral lists. The campaign started on Monday and will last for three months, until the end of March.

According to the Director General of ELECAM, many young Cameroonians between 18 and 25 years old have never registered to vote. ``Democracy is not a spectator sport,'' he declared during a press conference at the ELECAM headquarters in Yaoundé. ``If you do not register, you cannot vote, and if you do not vote, you cannot complain about the leaders you have.''

Registration is free and simple. Citizens can register at any ELECAM office with their national identity card. Biometric kits have been sent to all the 360 subdivisions of the country so that people in rural areas can also register easily.

Civil society organisations have welcomed the campaign. ``Voting is a right, but it is also a duty,'' said Mrs Ngo Bell, a human rights activist. ``In a democracy, the ballot box is the most peaceful way to change things. We must also teach young people tolerance: after the elections, winners and losers must live together in the same country.''

Cameroon is a multicultural nation with more than 250 ethnic groups and two official languages, English and French. For many observers, this diversity is a strength, but it requires dialogue, mutual respect and the promotion of basic freedoms such as freedom of speech and freedom of the press.
\end{exemplebox}

\courssub{Reading Skills: Skimming, Scanning, Identifying Opinions}

\begin{methode}[Three Reading Strategies]
\begin{enumerate}
\item \textbf{Skimming} (lecture globale): read the title, the first sentence of each paragraph and the conclusion quickly to find the \cle{main idea}. Do not stop at unknown words.
\item \textbf{Scanning} (lecture sélective): move your eyes quickly over the text to find a \cle{specific detail}: a date, a number, a name, a place.
\item \textbf{Identifying the writer's opinion}: look for opinion markers such as \emph{``I think'', ``we must'', ``it is a strength''}, and quotation marks with reported speech. Distinguish \textbf{facts} (verifiable information) from \textbf{opinions} (what someone believes).
\end{enumerate}
\end{methode}

\textbf{Skimming task.} Read the text quickly and choose the main idea:
\begin{itemize}
\item[a)] ELECAM is organising the presidential election.
\item[b)] ELECAM wants young people to register on the electoral lists.
\item[c)] ELECAM is building new offices in Yaoundé.
\end{itemize}
The correct answer is \textbf{b)}.

\textbf{Scanning task.} Find these details in the text as fast as possible: How long will the campaign last? How many subdivisions received biometric kits? What document do you need to register? (\emph{Answers: three months; 360 subdivisions; the national identity card.})

\courssub{Comprehension Tasks}

\begin{exercice}[True or False]
Read the text again and write \textbf{True} or \textbf{False}. Justify with a line from the text.
\begin{enumerate}
\item The campaign will last for six months.
\item Many young Cameroonians have never registered to vote.
\item Registration costs 1\,000 FCFA.
\item Mrs Ngo Bell thinks voting is only a right, not a duty.
\item Cameroon has two official languages.
\end{enumerate}
\end{exercice}

\begin{corrige}[Exercise: True or False]
\begin{enumerate}
\item \textbf{False} -- ``The campaign \ldots will last for three months.''
\item \textbf{True} -- ``Many young Cameroonians between 18 and 25 years old have never registered to vote.''
\item \textbf{False} -- ``Registration is free and simple.''
\item \textbf{False} -- She said: ``Voting is a right, but it is also a duty.''
\item \textbf{True} -- ``Two official languages, English and French.''
\end{enumerate}
\end{corrige}

\begin{exercice}[Wh-Questions]
Answer with complete sentences.
\begin{enumerate}
\item Who launched the registration campaign?
\item When did the campaign start?
\item Where can citizens register?
\item Why does the Director General say democracy is not a spectator sport?
\item What must winners and losers do after elections, according to Mrs Ngo Bell?
\end{enumerate}
\end{exercice}

\begin{corrige}[Exercise: Wh-Questions]
\begin{enumerate}
\item Elections Cameroon (ELECAM) launched the campaign.
\item It started on Monday, 5 January.
\item Citizens can register at any ELECAM office.
\item Because if you do not register and vote, you cannot choose your leaders or complain about them.
\item They must live together peacefully in the same country.
\end{enumerate}
\end{corrige}

\begin{exoresolu}[Matching Vocabulary to Definitions]
Match each word to its definition: \textbf{ballot, discrimination, rule of law, polling station, freedom of speech}.
\begin{enumerate}
\item[a)] The place where people go to vote.
\item[b)] Unfair treatment of a person because of his or her ethnic group, sex or religion.
\item[c)] The right to express your opinions freely.
\item[d)] The paper used to vote, or the act of voting.
\item[e)] The principle that everybody, including leaders, must obey the law.
\end{enumerate}
\tcblower
\textbf{Solution:} 1--a) polling station; 2--b) discrimination; 3--c) freedom of speech; 4--d) ballot; 5--e) rule of law. \emph{Tip:} scan the vocabulary table and the text; the context always confirms the meaning.
\end{exoresolu}

\courssub{Grammar Focus 1: Prepositions of Time}

Prepositions are small words (\cle{in, on, at, since, for, during, by}) that connect a noun to the rest of the sentence. They are short but very important: a wrong preposition changes the meaning.

\begin{methode}[Choosing the Right Preposition of Time]
\begin{enumerate}
\item \textbf{at} + precise time: \emph{at 8 o'clock, at noon, at night, at the weekend}.
\item \textbf{on} + day or date: \emph{on Monday, on 5 January, on election day}.
\item \textbf{in} + month, year, season, part of the day: \emph{in January, in 2026, in the morning}.
\item \textbf{since} + starting point: \emph{since 2018, since Monday}.
\item \textbf{for} + duration: \emph{for three months, for seven years}.
\item \textbf{during} + an event or period (followed by a noun): \emph{during the campaign, during the press conference}.
\item \textbf{by} + deadline (not later than): \emph{register by 31 March}.
\end{enumerate}
\end{methode}

\begin{attention}[Frequent Errors]
\begin{itemize}
\item \emph{in Monday} is wrong. We say \textbf{on Monday} (day $\rightarrow$ \textbf{on}).
\item \textbf{since} + starting point: \emph{since January}. \textbf{for} + duration: \emph{for three months}. Never say \emph{since three months}.
\item No preposition before \emph{today, yesterday, tomorrow, last week, next year}: \emph{The campaign started last week} (not \emph{on last week}).
\end{itemize}
\end{attention}

\begin{exemplebox}[Prepositions of Time in the Reading Text]
\begin{itemize}
\item The campaign started \textbf{on} Monday. (day)
\item It will last \textbf{for} three months. (duration)
\item Citizens must register \textbf{by} the end of March. (deadline)
\item The Director spoke \textbf{during} a press conference. (event)
\item ELECAM has organised elections \textbf{since} its creation. (starting point)
\end{itemize}
\end{exemplebox}

\courssub{Grammar Focus 2: Prepositions of Place and of Manner}

\textbf{Prepositions of place} tell us where something is:

\begin{center}
\begin{tabularx}{\linewidth}{|l|X|}\hline
\rowcolor{popgreenL} \textbf{Preposition} & \textbf{Use and example} \\ \hline
in & inside a space or area: \emph{in Yaoundé, in Cameroon, in the office} \\ \hline
on & on a surface or line: \emph{on the electoral list, on the table} \\ \hline
at & a precise point: \emph{at the polling station, at the ELECAM office} \\ \hline
between & with two people or things: \emph{between the two candidates} \\ \hline
among & with three or more: \emph{among the 250 ethnic groups} \\ \hline
opposite & facing: \emph{The office is opposite the post office.} \\ \hline
\end{tabularx}
\end{center}

\textbf{Prepositions of manner} tell us how something is done:

\begin{itemize}
\item \textbf{by} + means of transport or action: \emph{by bus, by voting, by peaceful means};
\item \textbf{with} + instrument or feeling: \emph{with a pen, with respect};
\item \textbf{in} + language or way: \emph{in English, in a peaceful way};
\item \textbf{like} + comparison: \emph{He spoke like a true statesman}.
\end{itemize}

\begin{exemplebox}[A Numbered Example to Complete]
``Cameroon has more than 250 ethnic groups living together \textbf{in} the same country.'' -- We use \textbf{in} (preposition of place) because the country is an area, a large space. Other examples: \emph{Dialogue \textbf{among} communities is essential. The polling station is \textbf{opposite} the market. Voters expressed their choice \textbf{in} a peaceful way.}
\end{exemplebox}

\courssub{Cameroon: Diversity and Living Together}

\begin{popfigure}
\begin{center}
\begin{tikzpicture}[scale=0.9]
% Simplified outline of Cameroon
\draw[thick, popdark, fill=popcream] (2.2,4.6) -- (1.6,3.6) -- (0.6,3.4) -- (0.2,2.4) -- (0.5,1.2) -- (0.3,0.2) -- (1.4,0.0) -- (2.6,0.4) -- (3.4,0.2) -- (4.4,0.8) -- (5.2,1.6) -- (5.0,2.6) -- (4.2,3.2) -- (3.4,3.0) -- (3.0,3.8) -- (2.6,4.4) -- cycle;
% Linguistic areas
\fill[popblue!35] (0.35,1.5) -- (1.2,1.6) -- (1.3,2.6) -- (0.5,2.7) -- (0.25,2.2) -- cycle;
\fill[popblue!35] (0.55,0.9) -- (1.5,0.9) -- (1.5,0.15) -- (0.5,0.2) -- cycle;
% Cities
\fill[popink] (1.9,1.1) circle (0.07); \node[right, font=\scriptsize] at (2.0,1.1) {Yaoundé};
\fill[popink] (1.1,0.55) circle (0.07); \node[left, font=\scriptsize] at (1.0,0.55) {Douala};
\fill[popink] (0.85,2.1) circle (0.07); \node[left, font=\scriptsize] at (0.75,2.1) {Bamenda};
\fill[popink] (0.9,0.5) circle (0.07); \node[below, font=\scriptsize] at (0.9,0.42) {Buea};
\fill[popink] (2.3,3.5) circle (0.07); \node[right, font=\scriptsize] at (2.4,3.5) {Garoua};
\fill[popink] (2.5,4.35) circle (0.07); \node[right, font=\scriptsize] at (2.6,4.35) {Maroua};
\fill[popink] (1.7,2.6) circle (0.07); \node[right, font=\scriptsize] at (1.8,2.6) {Bafoussam};
\fill[popink] (4.3,1.7) circle (0.07); \node[right, font=\scriptsize] at (4.4,1.7) {Bertoua};
\fill[popink] (3.1,0.55) circle (0.07); \node[below, font=\scriptsize] at (3.1,0.45) {Ebolowa};
\fill[popink] (2.9,2.2) circle (0.07); \node[right, font=\scriptsize] at (3.0,2.2) {Ngaoundéré};
% North arrow
\draw[->, thick, popdark] (5.6,3.6) -- (5.6,4.4); \node[font=\scriptsize] at (5.6,4.6) {N};
% Legend
\fill[popblue!35] (5.9,1.9) rectangle (6.3,2.2); \node[right, font=\scriptsize] at (6.4,2.05) {English-speaking area};
\fill[popcream] (5.9,1.4) rectangle (6.3,1.7); \draw[popdark] (5.9,1.4) rectangle (6.3,1.7); \node[right, font=\scriptsize] at (6.4,1.55) {French-speaking area};
\fill[popink] (6.1,1.05) circle (0.07); \node[right, font=\scriptsize] at (6.4,1.05) {Regional capital};
\end{tikzpicture}
\end{center}
\legende{Simplified map of Cameroon: the 10 regional capitals and the two linguistic areas (North-West and South-West in blue). One country, two official languages, more than 250 ethnic groups: tolerance is essential. (Schematic map, not to scale.)}
\end{popfigure}

\begin{savaistu}[Did You Know?]
\cle{ELECAM} (Elections Cameroon) is the body in charge of organising, managing and supervising elections and referendums in Cameroon. It registers voters, prints ballot papers and announces provisional results. Cameroon is often called ``Africa in miniature'' because it gathers an extraordinary diversity of landscapes, climates, cultures and languages \textbf{in} one single country.
\end{savaistu}

\begin{popfigure}
\begin{center}
\begin{tikzpicture}
\begin{axis}[
  ybar, bar width=0.9cm,
  width=0.85\linewidth, height=6cm,
  xlabel={Age group}, ylabel={Registration rate (\%)},
  symbolic x coords={18-25,26-35,36-50,50+},
  xtick=data, ymin=0, ymax=100,
  ytick={0,20,40,60,80,100},
  nodes near coords, every node near coord/.append style={font=\scriptsize},
  ymajorgrids=true, grid style={popline!40},
]
\addplot[fill=poporange, draw=popdark] coordinates {(18-25,38) (26-35,55) (36-50,72) (50+,81)};
\end{axis}
\end{tikzpicture}
\end{center}
\legende{Voter registration rate by age group (fictitious but realistic data, in \%). Young people aged 18--25 register the least: this is why ELECAM targets them in its campaigns.}
\end{popfigure}

\textbf{Read the chart.} The bar chart shows that only 38\,\% of citizens aged 18--25 are registered, compared with 81\,\% of citizens over 50. The registration rate rises steadily with age: from 38\,\% to 55\,\%, then 72\,\% and finally 81\,\%. This gap of 43 points between the youngest and the oldest voters explains why the campaign in our reading text focuses on young voters.

\courssub{Discussion: Tolerance in a Multicultural Country}

\begin{exercice}[Pair Work: Speak Out]
Discuss these questions with your partner, then share your ideas with the class. Use prepositions of time, place and manner correctly.
\begin{enumerate}
\item Cameroon has more than 250 ethnic groups. Is this diversity a strength or a problem? Justify your answer.
\item How can we promote tolerance \textbf{in} our school and \textbf{among} our friends?
\item ``Freedom of speech means we can say anything, even insults.'' Do you agree? Why or why not?
\item Why is voting both a right and a duty?
\end{enumerate}
\end{exercice}

\begin{exercice}[Grammar Practice: Fill in the Blanks]
Complete with the correct preposition: \emph{in, on, at, since, for, during, by, between, among, with}.
\begin{enumerate}
\item The elections will take place \ldots\ October.
\item Polling stations open \ldots\ 8 a.m.
\item She has been a voter \ldots\ 2019.
\item He waited \ldots\ two hours to cast his ballot.
\item There was a long queue \ldots\ the polling station.
\item The journalist wrote the article \ldots\ English.
\item We must settle our differences \ldots\ dialogue, not violence.
\item The campaign ended \ldots\ 31 March.
\end{enumerate}
\end{exercice}

\begin{corrige}[Exercise: Fill in the Blanks]
1. \textbf{in} (month) -- 2. \textbf{at} (precise time) -- 3. \textbf{since} (starting point) -- 4. \textbf{for} (duration) -- 5. \textbf{at} (precise point) -- 6. \textbf{in} (language, manner) -- 7. \textbf{by} (means) -- 8. \textbf{by} (deadline).
\end{corrige}

\begin{aretenir}[Key Points to Remember]
\begin{itemize}
\item Democracy = citizens choose leaders through free and fair elections; tolerance = respecting differences.
\item Reading strategies: \textbf{skim} for the main idea, \textbf{scan} for details, distinguish \textbf{facts} from \textbf{opinions}.
\item Time: \textbf{at} + hour, \textbf{on} + day/date, \textbf{in} + month/year/season; \textbf{since} + starting point, \textbf{for} + duration, \textbf{by} + deadline, \textbf{during} + event.
\item Place: \textbf{in} (area), \textbf{on} (surface), \textbf{at} (point), \textbf{between} (two), \textbf{among} (three or more), \textbf{opposite} (facing).
\item Manner: \textbf{by} (means), \textbf{with} (instrument), \textbf{in} (language/way), \textbf{like} (comparison).
\end{itemize}
\end{aretenir}

Now that you can read and discuss texts about democracy and freedoms, and that you master prepositions of time, place and manner, you are ready for the next section: comparative adverbs, adverbs of negation and phrasal verbs, which will make your expression even richer and more precise.

\coursec{Comparative Adverbs, Adverbs of Negation and Phrasal Verbs}

In the previous section, we read texts about democracy, tolerance and basic freedoms. You probably noticed sentences like: \emph{``Citizens must participate \textbf{more actively} in public life''} or \emph{``We should \textbf{never give up} defending our rights.''} These sentences contain two powerful grammar tools: \cle{comparative adverbs} and \cle{phrasal verbs}. In this section, we will learn how to compare actions (who works \emph{faster}? who speaks \emph{more clearly}?), how to use strong negative adverbs like \emph{hardly} and \emph{never}, and how to master the phrasal verbs that English speakers use every day at work and in civic life.

\courssub{Revision: Adverbs of Manner}

\begin{definition}[Adverb of manner]
An \cle{adverb of manner} tells us \textbf{how} an action is done. It modifies a verb and usually answers the question \emph{How?} Most adverbs of manner are formed by adding \textbf{-ly} to an adjective.
\end{definition}

\begin{center}
\begin{tabularx}{\linewidth}{|l|X|X|}\hline
\rowcolor{popblueL} \textbf{Adjective} & \textbf{Adverb} & \textbf{Example} \\ \hline
slow & slowly & The old computer works \emph{slowly}. \\ \hline
careful & carefully & She fills in the forms \emph{carefully}. \\ \hline
happy & happily & (y $\rightarrow$ i) They voted \emph{happily}. \\ \hline
good & \textbf{well} (irregular!) & The mayor speaks English \emph{well}. \\ \hline
fast & \textbf{fast} (no change) & He runs \emph{fast}. \\ \hline
hard & \textbf{hard} (no change) & Cameroonians work \emph{hard}. \\ \hline
\end{tabularx}
\end{center}

\begin{attention}[Good $\neq$ well / Hard $\neq$ hardly]
\emph{Good} is an adjective; its adverb is \emph{well}: \og She is a \textbf{good} citizen. She behaves \textbf{well}. \fg{} Also, \emph{hard} (with effort) and \emph{hardly} (almost not) have completely different meanings: \og He works \textbf{hard} \fg{} = il travaille dur ; \og He \textbf{hardly} works \fg{} = il ne travaille presque pas.
\end{attention}

\courssub{Comparative and Superlative Forms of Adverbs}

Imagine two workers in a workshop in Douala. Paul finishes his task in two hours; John finishes the same task in one hour. We naturally want to \textbf{compare} their actions: \emph{John works \textbf{faster than} Paul.} This is the intuition behind the comparative of adverbs: we compare \textbf{the way} two actions are performed.

\begin{propriete}[Forming comparative and superlative adverbs]
\begin{itemize}
\item \textbf{Short adverbs} (one syllable: fast, hard, early, soon) take \textbf{-er} (comparative) and \textbf{the + -est} (superlative): \emph{fast $\rightarrow$ faster $\rightarrow$ the fastest}.
\item \textbf{Long adverbs} (almost all adverbs in \textbf{-ly}) take \textbf{more} (comparative) and \textbf{the most} (superlative): \emph{carefully $\rightarrow$ more carefully $\rightarrow$ the most carefully}.
\item The comparative is followed by \textbf{than}: \emph{She types more quickly than me.}
\end{itemize}
\end{propriete}

\begin{popfigure}
\begin{center}
\begin{tabular}{|l|l|l|l|}\hline
\rowcolor{popgreenL} \textbf{Adverb} & \textbf{Positive} & \textbf{Comparative} & \textbf{Superlative} \\ \hline
Short & fast & faster (than) & the fastest \\ \hline
Short & hard & harder (than) & the hardest \\ \hline
Short & early & earlier (than) & the earliest \\ \hline
Long (-ly) & carefully & more carefully (than) & the most carefully \\ \hline
Long (-ly) & efficiently & more efficiently (than) & the most efficiently \\ \hline
\rowcolor{poporangeL} Irregular & well & better & the best \\ \hline
\rowcolor{poporangeL} Irregular & badly & worse & the worst \\ \hline
\rowcolor{poporangeL} Irregular & far & farther / further & the farthest / furthest \\ \hline
\end{tabular}
\end{center}
\legende{Conjugation table of adverbs: positive, comparative and superlative forms. Short adverbs take -er/-est; -ly adverbs take more/most; well, badly and far are irregular.}
\end{popfigure}

\begin{aretenir}[Irregular adverbs to memorize]
\begin{itemize}
\item \cle{well} $\rightarrow$ \textbf{better} $\rightarrow$ \textbf{the best} : \emph{Our team played better than theirs.}
\item \cle{badly} $\rightarrow$ \textbf{worse} $\rightarrow$ \textbf{the worst} : \emph{He drives worse than his brother.}
\item \cle{far} $\rightarrow$ \textbf{farther/further} $\rightarrow$ \textbf{the farthest/furthest} : \emph{She walked further than the others.}
\end{itemize}
\end{aretenir}

\begin{exemplebox}[Two comparatives in one sentence]
\emph{``She works \textbf{harder} than her colleagues, so she was promoted \textbf{more quickly}.''}\par
Here we can spot two comparatives of adverbs: \textbf{harder} (short adverb + -er) compares the effort, and \textbf{more quickly} (more + -ly adverb) compares the speed of the promotion. Notice the structure: comparative + \emph{than} introduces the second element of the comparison.
\end{exemplebox}

\courssub{Adverbs of Negation: hardly, scarcely, never}

\begin{definition}[Adverbs of negation]
\cle{Hardly}, \cle{scarcely} and \cle{never} are negative adverbs. \emph{Hardly} and \emph{scarcely} mean \og almost not \fg{} (presque pas, \`a peine) ; \emph{never} means \og at no time \fg{} (jamais). A sentence with one of these adverbs is already negative: do \textbf{not} add another negation.
\end{definition}

\begin{itemize}
\item \emph{I \textbf{never} vote without reading the candidates' programmes.} (Je ne vote jamais sans\dots)
\item \emph{He \textbf{hardly} listens to other people's opinions.} (Il n'\'ecoute presque pas\dots)
\item \emph{There was \textbf{scarcely} enough time to discuss the project.} (Il n'y avait presque pas assez de temps\dots)
\end{itemize}

\begin{attention}[The double negation trap]
In standard English, two negatives in the same clause are \textbf{wrong}. Do not say: \og I \textbf{never don't} vote \fg{} ($\times$). Say instead: \og I \textbf{never vote} \fg{} or \og I \textbf{don't ever} vote \fg{}. The same rule applies with \emph{hardly}: \og I can't hardly wait \fg{} ($\times$) must become \og I \textbf{can hardly} wait \fg{} (Je peux \`a peine attendre).
\end{attention}

\begin{methode}[Translating \og \`a peine\dots que \fg{} : inversion after Hardly/Scarcely]
When \emph{Hardly} or \emph{Scarcely} begins the sentence, English uses an \textbf{inversion} with the past perfect, followed by \textbf{when}:
\begin{enumerate}
\item Start with \textbf{Hardly} or \textbf{Scarcely}.
\item Add \textbf{had} + subject + past participle.
\item Add \textbf{when} + simple past for the second action.
\end{enumerate}
Example: \emph{\textbf{Hardly had} the meeting \textbf{started when} the lights went off.} (\`A peine la r\'eunion avait-elle commenc\'e que la lumi\`ere s'est \'eteinte.)\par
Also: \emph{\textbf{Scarcely had} he arrived \textbf{when} the election results were announced.}
\end{methode}

\courssub{Phrasal Verbs}

\begin{definition}[Phrasal verb]
A \cle{phrasal verb} is a verb combined with a particle (a preposition or an adverb) whose meaning is \textbf{different from the original verb}. Example: \emph{give} = donner, but \emph{give up} = \textbf{stop, abandon} (abandonner). You cannot guess the meaning from the verb alone: you must learn each phrasal verb as a new word.
\end{definition}

\begin{popfigure}
\begin{center}
\begin{tikzpicture}[
  grow=right,
  level 1/.style={sibling distance=45mm, level distance=35mm},
  level 2/.style={sibling distance=14mm, level distance=42mm},
  every node/.style={font=\small},
  base/.style={rectangle, rounded corners, draw=popblue, fill=popblueL, thick, font=\small\bfseries, minimum width=2.2cm, text centered},
  pv/.style={rectangle, rounded corners, draw=poporange, fill=poporangeL, minimum width=2.5cm, text centered},
  sens/.style={rectangle, draw=poppurple, fill=poppurpleL, rounded corners, font=\footnotesize, minimum width=3cm, text centered}
]
\node[base] {take}
  child { node[pv] {take part in} child { node[sens] {participate (participer \`a)} } }
  child { node[pv] {take off} child { node[sens] {remove / leave (enlever, d\'ecoller)} } }
  child { node[pv] {take up} child { node[sens] {start a hobby (se mettre \`a)} } };
\begin{scope}[yshift=-4cm]
\node[base] {put}
  child { node[pv] {put off} child { node[sens] {postpone (remettre)} } }
  child { node[pv] {put on} child { node[sens] {wear (mettre, porter)} } }
  child { node[pv] {put out} child { node[sens] {extinguish (\'eteindre)} } };
\end{scope}
\begin{scope}[yshift=-8cm]
\node[base] {stand}
  child { node[pv] {stand for} child { node[sens] {represent / tolerate} } }
  child { node[pv] {stand up} child { node[sens] {rise (se lever)} } }
  child { node[pv] {stand by} child { node[sens] {support (soutenir)} } };
\end{scope}
\end{tikzpicture}
\end{center}
\legende{Phrasal verb tree: one base verb (take, put, stand) combined with different particles gives completely different meanings.}
\end{popfigure}

\begin{aretenir}[Essential phrasal verbs for work and citizenship]
\begin{itemize}
\item \cle{carry out} = to perform, to do (effectuer) : \emph{The council carried out a survey on public health.}
\item \cle{look after} = to take care of (s'occuper de) : \emph{We must look after our environment.}
\item \cle{put off} = to postpone (remettre \`a plus tard) : \emph{Never put off your civic duties.}
\item \cle{give up} = to stop, to abandon (abandonner) : \emph{He gave up smoking for his health.}
\item \cle{stand for} = to represent (repr\'esenter) : \emph{What does this political party stand for?}
\item \cle{take part in} = to participate (participer \`a) : \emph{All citizens should take part in elections.}
\item \cle{bring up} = to raise children or a topic (\'elever, soulever) : \emph{She brought up an important question at the meeting.}
\item \cle{get along with} = to have a good relationship with (s'entendre avec) : \emph{He gets along with all his colleagues.}
\item \cle{follow up} = to pursue, to check (faire un suivi) : \emph{The committee will follow up on the recommendations.}
\item \cle{hand in} = to submit (remettre, d\'eposer) : \emph{Please hand in your application before Friday.}
\end{itemize}
\end{aretenir}

\begin{savaistu}[One phrasal verb, two meanings]
\emph{Stand for} can mean \og represent \fg{} : \emph{What does the flag stand for?} (Que repr\'esente le drapeau ?) But it can also mean \og tolerate \fg{}, usually in negative sentences: \emph{I won't \textbf{stand for} injustice!} (Je ne tol\'ererai pas l'injustice !) Context tells you which meaning is correct.
\end{savaistu}

\begin{exoresolu}[Rewriting with comparative adverbs and phrasal verbs]
Rewrite the following sentences using the word in brackets, without changing the meaning.\par
1. Amina types quickly. Blandine types very quickly. (than)\par
2. The workers decided to postpone the strike. (put off)\par
3. Every citizen must participate in community work. (take part in)\par
4. Our old machine works efficiently, but the new one works very efficiently. (more)
\tcblower
\textbf{Solution:}\par
1. \emph{Blandine types \textbf{more quickly than} Amina.} (or: \emph{Amina types \textbf{less quickly than} Blandine} -- \emph{quickly} is a long adverb $\rightarrow$ more quickly.)\par
2. \emph{The workers decided to \textbf{put off} the strike.}\par
3. \emph{Every citizen must \textbf{take part in} community work.}\par
4. \emph{The new machine works \textbf{more efficiently than} the old one.}
\end{exoresolu}

\begin{exercice}[Practice: work and citizenship]
\textbf{A.} Complete with the comparative of the adverb in brackets.\par
1. Since the training, the secretary works \dots\dots{} (efficiently).\par
2. Drivers in big cities should drive \dots\dots{} (carefully).\par
3. Our team arrived \dots\dots{} (early) than the visitors.\par
4. After his mistake, he behaved \dots\dots{} (badly) than before.\par
\textbf{B.} Replace the underlined expression with a phrasal verb: \emph{carry out, look after, give up, stand for}.\par
1. The government must \emph{perform} reforms to protect basic freedoms.\par
2. A good citizen \emph{takes care of} public property.\par
3. The candidate explained what his party \emph{represents}.\par
4. She decided to \emph{stop} wasting time and registered to vote.\par
\textbf{C.} Translate: \og \`A peine le pr\'esident avait-il fini son discours que la foule a applaudi. \fg{}
\textbf{D.} Choose the correct option.\par
1. I have \emph{never / never not} seen such a peaceful election.\par
2. She \emph{hardly / can't hardly} knows her neighbours.\par
3. The delegates discussed the issue \emph{more calmly / calmer} than expected.
\end{exercice}

\begin{corrige}[Exercice : Practice]
\textbf{A.} 1. more efficiently \quad 2. more carefully \quad 3. earlier \quad 4. worse.\par
\textbf{B.} 1. carry out $\rightarrow$ \emph{The government must \textbf{carry out} reforms...} \quad 2. \emph{... \textbf{looks after} public property.} \quad 3. \emph{... what his party \textbf{stands for}.} \quad 4. \emph{... to \textbf{give up} wasting time...}\par
\textbf{C.} \emph{\textbf{Hardly had} the president finished his speech \textbf{when} the crowd applauded.}\par
\textbf{D.} 1. never \quad 2. hardly \quad 3. more calmly.
\end{corrige}

\begin{aretenir}[Key points of this section]
\begin{itemize}
\item Short adverbs: \textbf{-er / the -est} (faster, the fastest); long adverbs in -ly: \textbf{more / the most} (more carefully, the most carefully).
\item Irregular: \emph{well $\rightarrow$ better $\rightarrow$ best}; \emph{badly $\rightarrow$ worse $\rightarrow$ worst}; \emph{far $\rightarrow$ further $\rightarrow$ furthest}.
\item \emph{Hardly, scarcely, never} are already negative: \textbf{no double negation}.
\item Inversion: \emph{Hardly had + subject + past participle + when...}
\item A phrasal verb = verb + particle with a \textbf{new meaning}: learn it as one word.
\end{itemize}
\end{aretenir}

\coursec{Writing: Compositions on Democracy, Tolerance and Freedoms}

In the previous sections of this unit, you have studied the grammar and vocabulary tools needed for this writing task: \cle{coordinating and subordinating conjunctions} (\emph{since, for, yet, although, in spite of, despite, still, however, nevertheless}), \cle{pronouns} (\emph{each, each other, one another}), \cle{prepositions and prepositional phrases} of time, place and manner, \cle{comparative adverbs}, \cle{adverbs of negation} (\emph{hardly, scarcely, never}) and \cle{phrasal verbs} (\emph{stand up for, put up with, look after, give up}). You have also built vocabulary on \cle{work culture}, \cle{civic responsibility}, \cle{democracy}, \cle{tolerance} and \cle{basic freedoms}. Now you will assemble all these resources into a complete, exam-ready composition.

\courssub{Understanding the Task: Types of Compositions}

Before writing a single word, you must identify \emph{exactly what type of text} the examination subject requires. At the Probatoire and Baccalauréat technique, three formats appear regularly. Misidentifying the format costs heavy marks.

\begin{definition}[The three common composition types at Probatoire/Baccalauréat technique]
\begin{itemize}
\item \textbf{Argumentative essay} (\emph{l'essai argumentatif}): you present a clear thesis on a given topic and defend it with logical arguments and concrete examples. Typical prompt: \emph{``Why every young Cameroonian should vote.''} or \emph{``Is tolerance necessary in a modern society?''}
\item \textbf{Formal letter} (\emph{la lettre formelle}): addressed to an authority (Divisional Officer, Minister, Editor of a newspaper) or an institution. It denounces a problem (corruption, tribalism, violation of freedoms) and proposes solutions. It \emph{must} contain: sender's address, recipient's address, date, subject line, formal salutation (\emph{``Sir,'' / ``Madam,''}), structured body, formal closing (\emph{``Yours faithfully,''}) and signature.
\item \textbf{Speech} (\emph{le discours}): a text written to be delivered orally on a public occasion (Youth Day, Human Rights Day, school assembly). It opens with a protocol greeting (\emph{``Mr Chairman, Distinguished Guests, Fellow Students...''}), uses rhetorical devices (repetition, direct address, rhetorical questions), and ends with a strong \cle{call to action} (\emph{``Let us all stand up for tolerance today!''}).
\end{itemize}
\end{definition}

\begin{attention}[Read the instructions twice -- underline key constraints]
Many candidates lose marks because they write an essay when the subject asks for a letter, or they ignore the word limit. \textbf{Always underline}: (1) the \emph{type of text} required, (2) the \emph{word count} (usually \textbf{150--200 words} for Probatoire, \textbf{200--250 words} for Baccalauréat), (3) the \emph{exact topic}. Respect them \emph{exactly}.
\end{attention}

\courssub{Planning: Brainstorming and Building an Outline}

A successful composition is never improvised. Good writers spend \textbf{10 minutes planning} before writing. This investment saves time later and guarantees coherence.

\begin{methode}[Step-by-step planning for an argumentative essay]
\begin{enumerate}
\item \textbf{Analyse the prompt.} Identify the key terms. Example: \emph{``Promoting tolerance in my school''} $\rightarrow$ key terms: \emph{promoting} (action), \emph{tolerance} (core value), \emph{school} (context).
\item \textbf{Brainstorm freely (3 minutes).} Write down every idea, example, word, or personal experience that comes to mind. Do not judge yet. For \emph{tolerance in school}: sharing food at break, respecting different religions, no tribal jokes, helping new students, speaking kindly, accepting defeat in sports, mixed group work in workshops, celebrating cultural days.
\item \textbf{Select and order (3 minutes).} Choose the \textbf{three strongest, most distinct ideas}. Order them strategically: \emph{simplest/most obvious $\rightarrow$ deeper/more persuasive $\rightarrow$ strongest/broadest impact}.
\item \textbf{Write the thesis statement (1 minute).} Formulate your position in one clear sentence. \emph{``Promoting tolerance in my school is the foundation of peaceful learning and academic success.''}
\item \textbf{Sketch the outline (3 minutes).} Use the structure below: Introduction (topic + thesis), Body 1 (Argument 1 + Example), Body 2 (Argument 2 + Example), Body 3 (Argument 3 + Example), Conclusion (restate thesis + final message).
\end{enumerate}
\end{methode}

\begin{popfigure}
\begin{center}
\begin{tikzpicture}[
  box/.style={draw=popblue, thick, fill=popblueL, rounded corners, text width=3.8cm, align=center, minimum height=1.5cm, font=\small},
  boxg/.style={box, fill=popgreenL, draw=popgreen},
  boxp/.style={box, fill=poppinkL, draw=poppink},
  fleche/.style={-{Stealth[length=3mm]}, very thick, poporange}]
\node[box, align=center] (intro) at (0,0){\textbf{Introduction}\\ \textasciitilde 20--30 words\\ Topic presentation + Thesis statement};
\node[boxg, align=center] (b1) at (5.8,2.4){\textbf{Body Paragraph 1}\\ \textasciitilde 40--50 words\\ Argument 1 + Concrete example};
\node[boxg, align=center] (b2) at (5.8,0){\textbf{Body Paragraph 2}\\ \textasciitilde 40--50 words\\ Argument 2 + Concrete example};
\node[boxg, align=center] (b3) at (5.8,-2.4){\textbf{Body Paragraph 3}\\ \textasciitilde 40--50 words\\ Argument 3 + Concrete example};
\node[boxp, align=center] (conc) at (11.6,0){\textbf{Conclusion}\\ \textasciitilde 20--30 words\\ Restate thesis + Call to action / Final thought};
\draw[fleche] (intro.east) -- (b1.west);
\draw[fleche] (intro.east) -- (b2.west);
\draw[fleche] (intro.east) -- (b3.west);
\draw[fleche] (b1.east) -- (conc.160);
\draw[fleche] (b2.east) -- (conc.west);
\draw[fleche] (b3.east) -- (conc.200);
\end{tikzpicture}
\end{center}
\legende{Standard argumentative essay structure for Probatoire/Baccalauréat technique (target: 150--200 words). Each body paragraph = ONE main idea + ONE developed example.}
\end{popfigure}

\courssub{Paragraph Structure: The Building Blocks}

An essay is only as strong as its paragraphs. Master the \cle{perfect paragraph} formula.

\begin{propriete}[The rule of the perfect paragraph]
A high-scoring paragraph contains exactly four elements in this order:
\begin{enumerate}
\item \textbf{Topic sentence} (\emph{phrase sujet}): states the \emph{single main idea} of the paragraph. No surprise, no question.
\item \textbf{Supporting sentences} (2--3): explain the idea, give reasons, and provide a \textbf{concrete, contextualised example} (from your school, quarter, Cameroon).
\item \textbf{Connector/Transition}: a linking word showing the logical relation to the next idea (\emph{moreover, however, therefore}).
\item \textbf{Concluding/Linking sentence}: summarises the paragraph's point or bridges to the next paragraph.
\end{enumerate}
\textbf{Golden rule:} One paragraph = One idea only.
\end{propriete}

\begin{exemplebox}[Annotated model paragraph]
\emph{\textbf{[Topic sentence]} Tolerance makes our school a better place to learn.} \emph{\textbf{[Support 1]} In Government Technical High School Bamenda, students come from all ten regions of Cameroon.} \emph{\textbf{[Support 2 + Example]} Although we speak different languages at home, we work together in the mechanical workshop and help \textbf{one another} during practical lessons. For example, during the last inter-class sports competition, students from different villages formed one team and won the trophy.} \emph{\textbf{[Concluding + Transition]} \textbf{Therefore}, accepting our differences makes us stronger \textbf{and prepares us for national unity.}}
\end{exemplebox}

\begin{attention}[Common paragraph errors]
\begin{itemize}
\item \textbf{No topic sentence} $\rightarrow$ the reader guesses the main idea.
\item \textbf{Two ideas in one paragraph} $\rightarrow$ split into two paragraphs.
\item \textbf{Example missing or vague} $\rightarrow$ ``For example, in my school...'' (name the place, the event, the people).
\item \textbf{Connector used as topic sentence} $\rightarrow$ ``Moreover, tolerance is good.'' (``Moreover'' connects, it does not state an idea).
\end{itemize}
\end{attention}

\courssub{Linking Ideas: Recycling Your Connectors}

Connectors are the \emph{signposts} of your logic. The examiner checks for \textbf{variety} and \textbf{accuracy}. You have learned them in the grammar lessons of this unit -- now \textbf{recycle} them deliberately.

\begin{center}
\begin{tabularx}{\linewidth}{|l|X|}\hline
\rowcolor{popblueL} \textbf{Logical Function} & \textbf{Connectors (vary them!)} \\ \hline
Addition / Reinforcement & moreover, in addition, furthermore, also, besides, what is more \\ \hline
Contrast / Concession & although, even though, however, nevertheless, yet, still, despite, in spite of, whereas, while \\ \hline
Cause / Reason & because, since, as, for, due to the fact that \\ \hline
Consequence / Result & therefore, so, consequently, as a result, that is why, hence \\ \hline
Exemplification & for example, for instance, such as, namely, a case in point is \\ \hline
Conclusion / Summary & in conclusion, to sum up, finally, in short, ultimately, all things considered \\ \hline
\end{tabularx}
\end{center}

\begin{methode}[How to vary connectors in your essay]
\begin{enumerate}
\item Map each logical move in your outline to a \textbf{different} connector from the table.
\item \textbf{Introduction $\rightarrow$ Body 1}: \emph{First of all,} / \emph{To begin with,} / \emph{The primary reason is that}
\item \textbf{Body 1 $\rightarrow$ Body 2}: \emph{Moreover,} / \emph{In addition,} / \emph{Furthermore,}
\item \textbf{Inside Body 2 (contrast)}: \emph{Although some argue...,} / \emph{Despite this,}
\item \textbf{Body 2 $\rightarrow$ Body 3}: \emph{What is more,} / \emph{Beyond that,}
\item \textbf{Body 3 $\rightarrow$ Conclusion}: \emph{In conclusion,} / \emph{Ultimately,} / \emph{All things considered,}
\end{enumerate}
\end{methode}

\begin{attention}[The ``And/But'' trap]
Writing \emph{``Voting is a right and it is a duty and it is important''} shows poor control. Write: \emph{``Voting is a right; \textbf{moreover}, it is a civic duty. \textbf{However}, many young people neglect it.''} \textbf{Never use the same connector twice} in a 180-word composition.
\end{attention}

\courssub{Using the Unit's Grammar Accurately -- The Examiner's Checklist}

The marking grid explicitly rewards \textbf{correct, natural use} of the grammar points taught in this unit. Insert them \emph{where they fit logically}, not artificially.

\begin{center}
\begin{tabularx}{\linewidth}{|l|l|X|}\hline
\rowcolor{popblueL} \textbf{Grammar Point} & \textbf{Required Form} & \textbf{Natural Integration Example} \\ \hline
Pronouns & \emph{each, each other, one another} & \emph{Each citizen has a duty. Students must respect \textbf{each other}. We should listen to \textbf{one another}.} \\ \hline
Prepositions & time (\emph{in, on, at}), place (\emph{in, at, on}), manner (\emph{with, by, in}) & \emph{Elections take place \textbf{in} October \textbf{at} the polling station. We must vote \textbf{with} dignity \textbf{in} a peaceful manner.} \\ \hline
Comparative adverbs & \emph{more/less + adverb + than}, \emph{as + adverb + as} & \emph{Young people participate \textbf{more actively than} before. We must listen \textbf{more carefully} to opposing views.} \\ \hline
Adverbs of negation & \emph{never, hardly, scarcely} (mid-position) & \emph{We must \textbf{never} accept tribal discrimination. Some citizens can \textbf{hardly} express their opinions freely.} \\ \hline
Phrasal verbs & \emph{stand up for, put up with, look after, give up, speak out against} & \emph{We must \textbf{stand up for} our rights, \textbf{put up with} differences, \textbf{look after} vulnerable peers, \textbf{give up} violence, \textbf{speak out against} injustice.} \\ \hline
Conjunctions & \emph{since, for, yet, although, despite, in spite of, however, nevertheless} & \emph{\textbf{Since} the youth are the majority, \textbf{although} some doubt their power, \textbf{nevertheless} each vote counts.} \\ \hline
\end{tabularx}
\end{center}

\begin{aretenir}[Grammar insertion strategy]
\begin{itemize}
\item \textbf{Plan the grammar} during your 10-minute outline: write one target structure per body paragraph in the margin.
\item \textbf{Do not force} all structures into one sentence. Spread them naturally.
\item \textbf{Accuracy > Complexity}: a simple correct sentence scores higher than a complex incorrect one.
\end{itemize}
\end{aretenir}

\courssub{Model Composition Analysis: ``Why Every Young Cameroonian Should Vote''}

Study this model (182 words) as a template. It integrates \textbf{all required grammar}, \textbf{varied connectors}, \textbf{concrete Cameroonian context}, and \textbf{clear paragraphing}.

\begin{exemplebox}[Model composition -- Argumentative Essay]
\emph{Democracy gives every citizen the power to choose their leaders. \textbf{In my opinion}, voting is not only a right but also a duty for every young Cameroonian.}

\emph{\textbf{First of all}, voting allows young people to shape their future. \textbf{Since} the youth represent the majority of the population, their votes can change the country's direction. \textbf{For example}, when students vote massively, they can elect leaders who will build more technical schools and create apprenticeships.}

\emph{\textbf{Moreover}, voting teaches civic responsibility. \textbf{Although} some young people think their single vote is useless, \textbf{each} vote counts. A young citizen who votes learns to \textbf{stand up for} his ideas and to respect the choice of \textbf{others}.}

\emph{\textbf{However}, voting must be done peacefully. We should \textbf{never} fight because of politics; instead, we must \textbf{put up with} different opinions and \textbf{speak out against} violence.}

\emph{\textbf{In conclusion}, every young Cameroonian should register and vote, \textbf{because} voting builds democracy, responsibility and peace. \textbf{Let us all take part in the future of our country!}}
\end{exemplebox}

\begin{exoresolu}[Detailed analysis of the model composition]
\textbf{Task:} Identify (a) the thesis statement, (b) the three arguments with their examples, (c) six different connectors, (d) three phrasal verbs, (e) one adverb of negation, (f) two pronouns (\emph{each/each other/one another}), (g) two prepositional phrases (time/place/manner), (h) one comparative adverb.
\tcblower
\textbf{Solution.}
(a) \textbf{Thesis}: \emph{``voting is not only a right but also a duty for every young Cameroonian''} (end of introduction).
(b) \textbf{Arguments}:
\begin{enumerate}
\item \emph{Shape future} -- Example: students elect leaders who build technical schools/create apprenticeships.
\item \emph{Teaches civic responsibility} -- Example: voter learns to stand up for ideas, respect others' choice.
\item \emph{Must be peaceful} -- Example: put up with differences, speak out against violence.
\end{enumerate}
(c) \textbf{Connectors}: \emph{In my opinion, First of all, Since, For example, Moreover, Although, However, instead, In conclusion, because} (9 different ones used).
(d) \textbf{Phrasal verbs}: \emph{stand up for, put up with, speak out against}.
(e) \textbf{Adverb of negation}: \emph{never} (\emph{``We should never fight...''}).
(f) \textbf{Pronouns}: \emph{each} (\emph{``each vote''}), \emph{others} (\emph{``choice of others''} -- reciprocal sense).
(g) \textbf{Prepositional phrases}: \emph{in the country's direction} (place/manner), \emph{with} dignity (manner -- implied in ``peacefully''), \emph{at} the polling station (place -- typical context).
(h) \textbf{Comparative adverb}: \emph{more actively} (implied in ``participate more actively than before'' -- not explicit in this model but required in your writing).
\end{exoresolu}

\courssub{Guided Writing: Argumentative Essay (Exam Practice)}

\begin{exemplebox}[Probatoire-type subject]
\emph{``Write a composition of 150--200 words on: \textbf{The importance of freedom of speech in a democracy}.''}
\textbf{Plan (10 min):}
\begin{itemize}
\item \textbf{Thesis}: Freedom of speech is the pillar of democracy.
\item \textbf{Arg 1}: Citizens can criticise bad governance $\rightarrow$ Example: press exposes corruption.
\item \textbf{Arg 2}: Informed citizens make better choices $\rightarrow$ Example: radio debates before elections.
\item \textbf{Arg 3}: Freedom has limits (respect others) $\rightarrow$ Example: hate speech laws.
\item \textbf{Conclusion}: Defend free speech responsibly.
\end{itemize}
\textbf{Grammar targets to insert}: \emph{although, moreover, however, never, stand up for, each other, more freely than, in + month, at + place}.
\end{exemplebox}

\begin{exercice}[Guided Writing -- Argumentative Essay]
Write a composition of \textbf{150--200 words} on: \emph{``Promoting tolerance in my school / community.''}
\textbf{Requirements:}
\begin{enumerate}
\item Follow the 5-paragraph structure (Intro, 3 Body, Conclusion).
\item Use \textbf{at least four different connectors} from the table.
\item Include \textbf{one phrasal verb}, \textbf{one adverb of negation} (\emph{never/hardly/scarcely}), \textbf{the pronoun \emph{each other} or \emph{one another}}.
\item Give \textbf{one concrete Cameroonian example} per body paragraph (name a place, event, or situation).
\item Count your words and write the total at the end.
\end{enumerate}
\end{exercice}

\begin{corrige}[Guided Writing -- Sample Answer (178 words)]
\emph{Tolerance means accepting people who are different from us. \textbf{In my opinion}, promoting tolerance in my school is the best way to live and learn in peace.}

\emph{\textbf{First}, tolerance creates friendship among students. \textbf{In} my school \textbf{in} Yaoundé, we come from different regions and religions, \textbf{yet} we share our food at break time and help \textbf{each other} with homework. \textbf{For instance}, during the last cultural week, students from the Far North and the South-West organised a joint dance show that won first prize.}

\emph{\textbf{Moreover}, tolerance improves our academic results. \textbf{Since} we work in mixed groups in the electrical workshop, everyone brings different skills, and we succeed \textbf{more easily than} when we work alone. \textbf{However}, tolerance is not always easy. We must \textbf{never} laugh at a classmate's accent or tribe; instead, we must \textbf{put up with} our differences and \textbf{stand up for} any student who is rejected.}

\emph{\textbf{In conclusion}, tolerance makes my school a family. \textbf{If} each student respects \textbf{one another}, our school will become an example for the whole community. \textbf{Let us build that future today!}}
\end{corrige}

\courssub{Guided Writing: Formal Letter and Speech (Other Exam Formats)}

The Probatoire/Baccalauréat may ask for a \textbf{letter} or a \textbf{speech}. The argumentative core remains; only the \textbf{format} changes.

\begin{exemplebox}[Formal Letter -- Structure \& Sample Opening]
\textbf{Subject:} \emph{``Write a letter to the Divisional Officer denouncing tribal discrimination in your quarter and proposing solutions.''}

\emph{[Your Address]\\
\relax[Date]\\
The Divisional Officer\\
\relax[Division]\\
\relax[Region]\\
\textbf{Subject: Denouncing Tribal Discrimination in [Quarter Name]}\\
\textbf{Sir,}\\
\emph{I am writing to draw your attention to the growing tribal discrimination in our quarter... [Body paragraphs using same argument structure] ... I urge you to organise quarterly dialogue meetings and sanction offenders.}\\
\textbf{Yours faithfully,}\\
\relax[Signature]\\
\relax[Your Full Name]}
\end{exemplebox}

\begin{exemplebox}[Speech -- Structure \& Sample Opening]
\textbf{Occasion:} Youth Day, 11 February. \textbf{Subject:} \emph{``The Role of Youth in Defending Basic Freedoms.''}

\emph{\textbf{Mr Chairman, Distinguished Authorities, Dear Friends,}\\
\emph{Today, as we celebrate Youth Day, I want to speak about the duty of every young Cameroonian to defend our basic freedoms. \textbf{First of all}, freedom of expression allows us to \textbf{speak out against} injustice... [Body paragraphs] ... \textbf{In conclusion}, let us \textbf{never} be silent. \textbf{Let us stand up for} our rights and build a Cameroon where \textbf{each} voice counts! \textbf{Thank you.}}}
\end{exemplebox}

\begin{exercice}[Format Conversion Practice]
Take the sample answer on \emph{``Promoting tolerance in my school''} (178 words) and:
\begin{enumerate}
\item Rewrite it as a \textbf{formal letter to the Principal} (add addresses, date, subject, salutation, closing). Keep 150--200 words.
\item Rewrite it as a \textbf{speech for the School Assembly} (add protocol greeting, rhetorical question, call to action). Keep 150--200 words.
\end{enumerate}
\end{exercice}

\courssub{Peer Correction and Self-Assessment: The Examiner's Grid}

After writing, exchange compositions with a partner. Use this official-style grid. Be strict -- generosity hides errors.

\begin{center}
\begin{tabularx}{\linewidth}{|l|X|c|}\hline
\rowcolor{popblueL} \textbf{Criterion} & \textbf{Guiding Questions} & \textbf{Max} \\ \hline
\textbf{Content} (Ideas) & Are ideas relevant, developed, and supported by \textbf{concrete Cameroonian examples}? Is the thesis clear? & /4 \\ \hline
\textbf{Organisation} & Clear 5-paragraph structure? Logical flow? Topic sentences? Effective conclusion? & /4 \\ \hline
\textbf{Connectors} & \textbf{At least 4 different connectors} used correctly? Varied functions (addition, contrast, cause, conclusion)? & /3 \\ \hline
\textbf{Grammar Accuracy} & Target structures used \textbf{correctly}: pronouns (\emph{each/each other}), prepositions (time/place/manner), comparative adverbs, adverbs of negation, phrasal verbs, conjunctions? Tenses consistent? & /4 \\ \hline
\textbf{Vocabulary Range} & Words from unit: \emph{democracy, tolerance, freedom, civic responsibility, discrimination, governance, ballot, petition, advocacy, unity, diversity, harmony}? No repetition of basic words? & /3 \\ \hline
\textbf{Mechanics} & Spelling, punctuation, capitalisation? Word count respected (150--200 / 200--250)? Legible handwriting? & /2 \\ \hline
\textbf{TOTAL} &  & \textbf{/20} \\ \hline
\end{tabularx}
\end{center}

\begin{methode}[Self-correction routine (do this alone before handing in)]
\begin{enumerate}
\item \textbf{Count words.} If $<$ 150 or $>$ 200 (Probatoire), adjust.
\item \textbf{Circle every connector.} Are there $\ge$ 4 different ones? Any repeated? Replace repeats.
\item \textbf{Highlight target grammar:} pronouns (yellow), prepositions (green), comparative adverbs (blue), negation adverbs (pink), phrasal verbs (orange), conjunctions (purple). \textbf{Missing one? Insert it now.}
\item \textbf{Check each paragraph:} Topic sentence? Example? Transition?
\item \textbf{Read aloud.} Does it sound natural? Fix awkward phrasing.
\end{enumerate}
\end{methode}

\begin{savaistu}[How your composition is marked at Probatoire/Baccalauréat technique]
The composition carries \textbf{20 marks} (often 20\% of the English paper). Markers use a grid similar to the one above. Key insights from past marking reports:
\begin{itemize}
\item \textbf{Content + Organisation} = 8 marks. A clear thesis + 3 developed arguments with \textbf{local examples} (your school, your quarter, Cameroon) scores highly.
\item \textbf{Grammar} = 4 marks. \textbf{Accuracy} of the \emph{unit's specific structures} is checked first. Errors in basic tenses (present simple, past simple) are penalised heavily.
\item \textbf{Vocabulary} = 3 marks. Using \emph{unit-specific words} (\emph{civic responsibility, ballot, advocacy, tolerance, discrimination, freedom of speech, governance}) shows mastery. ``Good'', ``bad'', ``nice'' lower the score.
\item \textbf{Mechanics} = 2 marks. Easy marks -- don't lose them through carelessness.
\item \textbf{Quality beats quantity.} A well-organised 160-word essay with correct grammar and rich vocabulary \textbf{outscores} a 250-word essay full of errors and repetition.
\end{itemize}
\end{savaistu}

\begin{aretenir}[Key Points to Remember for Exam Success]
\begin{itemize}
\item \textbf{Identify the format} (essay, letter, speech) -- underline it in the subject.
\item \textbf{Plan 10 minutes:} brainstorm $\rightarrow$ select 3 ideas $\rightarrow$ order $\rightarrow$ thesis $\rightarrow$ outline.
\item \textbf{Structure:} 5 paragraphs (Intro, Body$\times$3, Conclusion). One idea per paragraph.
\item \textbf{Paragraph anatomy:} Topic sentence + 2--3 supports (with \textbf{concrete example}) + Transition.
\item \textbf{Connectors:} Minimum 4 different ones, varied functions, no repetition.
\item \textbf{Grammar showcase:} Naturally insert pronouns, prepositions, comparative adverbs, negation adverbs, phrasal verbs, conjunctions -- one per paragraph planned in advance.
\item \textbf{Vocabulary:} Use the unit's lexical field (democracy, tolerance, freedoms, civic responsibility...).
\item \textbf{Self-correct:} Word count, connectors, grammar targets, paragraph structure, mechanics.
\item \textbf{Peer correct:} Use the grid. Be the strict examiner you want to impress.
\end{itemize}
\end{aretenir}

\coursec{Integration, Evaluation and Remediation}

In the previous section, you learned how to plan, draft and correct a composition about democracy, tolerance and basic freedoms. Writing is the last skill of the lesson, but learning does not stop there. Now it is time to \cle{integrate} everything: to combine speaking, listening, reading, grammar and writing in one real-life project, to check your progress through a \cle{summative evaluation} (évaluation sommative), and to repair your weaknesses through \cle{remediation} (remédiation). A good technician never delivers a machine without testing it; a good learner never finishes a unit without testing himself or herself.

\courssub{Integration Activity: A Real-Life Citizenship Project}

\begin{definition}[Integration activity]
An \cle{integration activity} is a task that combines several language skills (speaking, listening, reading, grammar and writing) around one meaningful, real-life situation. Its goal is not to test one rule, but to check that you can \emph{use} English to act in society.
\end{definition}

\textbf{The project: ``Our Community, Our Responsibility''.} In groups of four, you will prepare a short public presentation (5 minutes) for the class, as if you were addressing your village or neighbourhood committee. The project has five steps, and each step uses one skill of the lesson:

\begin{enumerate}
\item \textbf{Reading:} read a short text about a community problem (dirty market square, water shortage, absenteeism at work) and identify the causes.
\item \textbf{Listening:} listen to a recorded message from a local leader about civic responsibility and take notes.
\item \textbf{Speaking:} discuss in your group and agree on two realistic solutions; use connectors such as \emph{although, however, since, moreover}.
\item \textbf{Grammar:} check your script for correct pronouns (\emph{each, each other, one another}), prepositions of time, place and manner, and phrasal verbs (\emph{carry out, clean up, look after}).
\item \textbf{Writing:} produce a one-page written proposal that summarises your presentation.
\end{enumerate}

\begin{popfigure}
\begin{center}
\begin{tikzpicture}
\foreach \i/\name/\col in {90/Speaking/popblue, 162/Listening/popgreen, 234/Reading/poporange, 306/Grammar/poppurple, 18/Writing/poppink}{
  \fill[\col!70] (0,0) -- ({\i-36}:2.2) arc ({\i-36}:{\i+36}:2.2) -- cycle;
  \node[font=\small\bfseries, text=popdark] at (\i:1.45) {\name};
}
\fill[popcream] (0,0) circle (0.75);
\node[font=\small\bfseries, text=popdark, align=center] at (0,0) {Civic\\project};
\end{tikzpicture}
\end{center}
\legende{The five skills of the lesson integrated around one citizenship project. Each skill feeds the others: reading and listening give ideas, speaking and grammar shape them, and writing fixes them.}
\end{popfigure}

\begin{methode}[How to succeed in an integration task]
\begin{enumerate}
\item Read the situation carefully and underline the task: \emph{who} are you, \emph{to whom} do you speak, \emph{what} must you produce?
\item Gather ideas from the documents (text, recording) before inventing your own.
\item Share the work in the group: a secretary writes, a timekeeper watches the clock, everyone speaks.
\item Use at least three connectors and two phrasal verbs from the lesson.
\item Rehearse the presentation once, then correct the script before the final performance.
\end{enumerate}
\end{methode}

\courssub{Summative Evaluation}

\begin{definition}[Summative evaluation]
A \cle{summative evaluation} (évaluation sommative) is a final test given at the end of a lesson or unit. It measures what you have mastered and gives a mark. It is different from a \emph{formative} evaluation, which only helps you to progress during learning.
\end{definition}

The evaluation of this lesson has three parts, for a total of 20 marks.

\begin{exemplebox}[Typical marking scheme (barème type)]
\begin{center}
\begin{tabularx}{\linewidth}{|l|X|c|}
\hline
\rowcolor{popblueL} \textbf{Part} & \textbf{Content} & \textbf{Marks} \\
\hline
Grammar quiz & 10 items: conjunctions (\emph{although, despite, however, since, for, yet}), pronouns (\emph{each, each other, one another}), prepositions of time/place/manner, comparative adverbs, adverbs of negation (\emph{hardly, scarcely, never}), phrasal verbs & /10 \\
\hline
Vocabulary test & 5 items: work culture, civic responsibility, democracy and freedoms & /5 \\
\hline
Short composition & 80--100 words: ``Why should every citizen vote?'' & /5 \\
\hline
\textbf{Total} & & \textbf{/20} \\
\hline
\end{tabularx}
\end{center}
\end{exemplebox}

\begin{exoresolu}[Sample grammar item (conjunctions)]
Complete the sentence: \emph{``........... it was raining heavily, the workers finished the bridge on time.''} (a) Despite \quad (b) Although \quad (c) However
\tcblower
The correct answer is \textbf{(b) Although}, because it is followed by a full clause (subject + verb: \emph{it was raining}). \emph{Despite} must be followed by a noun or -ing form (\emph{Despite the rain}), and \emph{However} introduces a new sentence, usually after a semicolon or a full stop. Typical distractors test exactly this difference.
\end{exoresolu}

\begin{attention}[The most expensive mistake: neglecting proofreading]
Many learners lose 2 or 3 marks not because they do not know the rule, but because they do not re-read. Always keep \textbf{the last 5 minutes} to check: (1) subject--verb agreement (\emph{each citizen has}, not \emph{have}); (2) prepositions of time (\emph{in} 2026, \emph{on} Monday, \emph{at} 8 o'clock); (3) spelling of key words (\emph{responsibility, democracy, freedom}); (4) double negation, which is wrong in standard English (\emph{I don't know nothing} $\rightarrow$ \emph{I don't know anything}).
\end{attention}

\courssub{Remediation: Repairing Your Weaknesses}

After the test, the teacher returns the scripts with the errors highlighted. Remediation means working \emph{precisely} on your own errors, not repeating the whole lesson.

\begin{methode}[How to remediate effectively]
\begin{enumerate}
\item \textbf{Identify the error:} name it exactly (``I wrote \emph{although\ldots but}'').
\item \textbf{Review the rule:} go back to the grammar table of the lesson and read the correct model aloud.
\item \textbf{Redo three targeted exercises} on that single point, then check with the correction.
\item \textbf{Get re-tested:} ask a classmate or the teacher to give you one new item on the same point. If you succeed, the gap is closed.
\end{enumerate}
\end{methode}

\textbf{Targeted exercise 1 -- Double negation.} Correct: \emph{``The voters didn't say nothing.''} $\rightarrow$ \emph{The voters didn't say anything.} (Standard English allows only one negation per clause.)

\textbf{Targeted exercise 2 -- Although / but.} Correct: \emph{``Although the work was hard, but they finished it.''} $\rightarrow$ \emph{Although the work was hard, they finished it.} (Never use \emph{although} and \emph{but} in the same sentence: one connector is enough.)

\textbf{Targeted exercise 3 -- Prepositions of time.} Fill in: the meeting is \emph{on} Friday \emph{at} 10 a.m., \emph{in} December. (Days $\rightarrow$ \emph{on}; clock times $\rightarrow$ \emph{at}; months and years $\rightarrow$ \emph{in}.)

\courssub{Peer Tutoring and Self-Evaluation}

\begin{definition}[Peer tutoring]
\cle{Peer tutoring} (tutorat entre pairs) is a strategy where a stronger learner explains one grammar point to a partner, using his or her own examples. Explaining forces you to organise your knowledge: teaching is the best way of learning.
\end{definition}

In pairs, learner A chooses one point (for example \emph{despite} vs \emph{in spite of}), explains the rule in simple English, gives two personal examples (``\emph{Despite the noise, I sleep well.}''), then learner B must produce one new sentence. After five minutes, roles are exchanged.

\begin{popfigure}
\begin{center}
\begin{tikzpicture}
\node[font=\small\bfseries, text=popdark] at (4.6,3.4) {Self-evaluation grid -- Lesson 4};
\begin{scope}
\draw[popline, fill=popblueL] (0,2.4) rectangle (9.2,3.0);
\node[font=\small\bfseries, text=popdark, anchor=west] at (0.1,2.7) {Competence};
\node[font=\small\bfseries, text=popdark] at (7.0,2.7) {I can do it well};
\node[font=\small\bfseries, text=popdark] at (8.6,2.7) {I need practice};
\foreach \y/\txt in {1.8/{I can use conjunctions (although, despite, however, since, for, yet)}, 1.2/{I can use pronouns (each, each other, one another)}, 0.6/{I can use prepositions of time, place and manner}, 0.0/{I can use comparative adverbs, phrasal verbs and adverbs of negation}, -0.6/{I can write a composition on democracy and freedoms}}{
  \draw[popline] (0,\y-0.3) rectangle (9.2,\y+0.3);
  \node[font=\small, text=popdark, anchor=west] at (0.1,\y) {\txt};
  \draw[popdark] (6.85,\y-0.12) rectangle (7.15,\y+0.12);
  \draw[popdark] (8.45,\y-0.12) rectangle (8.75,\y+0.12);
}
\end{scope}
\end{tikzpicture}
\end{center}
\legende{Tick one box per line. If you tick ``I need practice'', go back to the remediation method: identify the error, review the rule, redo three exercises, get re-tested.}
\end{popfigure}

\courssub{Extension for the Probatoire and Bac Technique}

\begin{attention}[Extension -- signalled as exam preparation]
The following task goes one step further: it prepares you directly for the writing paper of the Probatoire and Baccalauréat technique, where a formal letter is a frequent subject.
\end{attention}

\textbf{Task:} Write a formal letter (120--150 words) to your Sub-Divisional Officer about a community problem (for example, the dirty market square). Use the vocabulary of the lesson (\emph{civic responsibility, community, to carry out, to clean up}) and at least three connectors (\emph{since, although, however, moreover}).

\begin{savaistu}[The formal letter at the Probatoire technique]
The \cle{formal letter} is one of the most frequent writing tasks at the Probatoire technique. Its structure is fixed and earns easy marks: (1) your address and the date at the top right; (2) the recipient's address at the left; (3) the salutation (\emph{Dear Sir/Madam}); (4) the body in two or three short paragraphs; (5) the closing formula (\emph{Yours faithfully,}) and your signature. Examiners reward a clean layout as much as correct grammar.
\end{savaistu}

\begin{exercice}[Final self-check]
1. Choose the correct option: \emph{``Each student must respect one \ldots\ldots'' (other / another / each other).}
2. Rewrite without a mistake: \emph{``Despite of the rain, the meeting took place on 5 p.m.''}
3. In six lines, explain why voting is a civic duty, using \emph{since} and \emph{however}.
\end{exercice}

\begin{corrige}[Exercice Final self-check]
1. \emph{another} (``one another'' is the reciprocal pronoun; \emph{each other} cannot follow \emph{one} here).
2. \emph{Despite the rain, the meeting took place at 5 p.m.} (\emph{despite} never takes \emph{of} -- that is \emph{in spite of}; clock time takes \emph{at}, not \emph{on}).
3. Sample answer: \emph{Voting is a civic duty since it allows every citizen to choose the leaders of the community. Many people fought for this freedom; however, some citizens still do not register. If we all vote, our democracy becomes stronger.}
\end{corrige}

\begin{aretenir}[To remember]
Integration combines the five skills in one real project; the summative evaluation measures your mastery (grammar /10, vocabulary /5, composition /5); remediation repairs your exact errors in four steps; peer tutoring and the ``I can\ldots'' grid make you an autonomous learner -- the true mark of a citizen and of an excellent worker.
\end{aretenir}

\newpage

\coursec{Activité d'intégration}

\coursec{Citizenship and Human Rights -- Integration Activity}

\courssub{Objectifs de l'activité}

This integration activity (\emph{activité d'intégration}) brings together all the knowledge of the lesson \emph{Citizenship and Human Rights}. At the end of the activity, each student will be able to:

\begin{itemize}
\item analyse authentic documents (a photograph, a short article, statistics) and identify a citizenship problem in the school (\emph{identifier un problème de citoyenneté}) ;
\item use the vocabulary of \cle{work culture}, \cle{excellence at work}, \cle{civic responsibility}, \cle{democracy}, \cle{tolerance} and \cle{basic freedoms} in a real production task ;
\item write a 150-word campaign speech using at least 3 contrast conjunctions (\emph{although, however, nevertheless, yet, despite\ldots}), 2 pronouns (\emph{each, each other, one another}), 3 prepositions of time or place, 1 adverb of negation (\emph{hardly, scarcely, never}) and 2 phrasal verbs ;
\item deliver the speech orally with fluency and correct pronunciation (\emph{présenter le discours à l'oral}) ;
\item design a poster with a clear slogan (\emph{une affiche avec un slogan}) ;
\item evaluate the production of other groups with an assessment grid (\emph{grille d'évaluation}) and accept peer feedback.
\end{itemize}

\courssub{Situation}

\textbf{Citizenship in Action: A Campaign for Our School}

You are students of G.T.H.S. (\emph{Government Technical High School}) Nkol-Eton, Yaoundé. The school has 1\,200 students. Next month, the school will organise the election of class delegates (\emph{délégués de classe}) and the annual ``Clean School Day''. The principal is worried, and the Parent-Teacher Association has collected the following information.

\textbf{Document 1 -- Photograph.} A photograph of the school yard taken last Friday at 4 p.m.: plastic bags and empty bottles on the ground, papers near the dustbin, two students walking past the rubbish without picking it up.

\textbf{Document 2 -- Short article (school magazine).} ``Last year, only 40\,\% of the students of G.T.H.S. Nkol-Eton voted in the class delegate elections. Many students said: `Voting is useless, nothing changes.' Some candidates insulted each other during the campaign, and two groups quarrelled after the results. However, the elected delegates obtained a new water tap and repaired fans for three classrooms.''

\textbf{Document 3 -- Statistics (fictional data).}

\begin{center}
\begin{tabularx}{\linewidth}{|l|X|X|}\hline
\rowcolor{popblueL} \textbf{Indicator} & \textbf{Last year} & \textbf{This year (target)} \\ \hline
Participation in delegate elections & 40\,\% (480 students) & 75\,\% (900 students) \\ \hline
Students who clean the yard regularly & 25\,\% (300 students) & 60\,\% (720 students) \\ \hline
Reported quarrels during campaign week & 12 cases & 0 case \\ \hline
\end{tabularx}
\end{center}

The principal launches a competition: \emph{``Citizenship in Action''}. Each group of 4 students must prepare a campaign to convince the school community to vote, to keep the yard clean, to respect one another and to develop a culture of excellence at work. The best campaign will be presented during the general assembly (\emph{assemblée générale}).

\courssub{Consignes}

Work in groups of 4. Follow the steps in order.

\begin{enumerate}
\item \textbf{Document analysis.} Look at the photograph, read the article and study the table. Answer in full sentences: (a) What problems can you identify in the school? (b) How many students did \emph{not} vote last year? (c) What did the elected delegates obtain? (d) Why is 12 quarrels a problem for tolerance and democracy?
\item \textbf{Problem identification.} Choose \emph{one} main citizenship problem for your campaign (low participation, dirty yard, intolerance during elections, or lack of work culture). Justify your choice in two sentences, using one contrast conjunction.
\item \textbf{Writing the speech.} Write a campaign speech of about \cle{150 words}. Your speech must contain: at least \textbf{3 contrast conjunctions} (\emph{although, though, yet, however, nevertheless, despite, in spite of, still}), \textbf{2 pronouns} from the lesson (\emph{each, each other, one another}), \textbf{3 prepositions} of time or place (\emph{in, on, at, during, before, after, next to, in front of\ldots}), \textbf{1 adverb of negation} (\emph{hardly, scarcely, never}) and \textbf{2 phrasal verbs} (\emph{stand up for, look after, pick up, put forward, carry out, take part in}).
\item \textbf{Oral presentation.} Choose one or two speakers. Present your speech to the class in 2 minutes. Speak clearly, look at the audience (\emph{regarder l'auditoire}) and use gestures.
\item \textbf{Poster.} Produce an A4 poster with a short slogan (maximum 10 words), one drawing or symbol, and the name of your group.
\item \textbf{Peer evaluation.} While other groups present, fill in the assessment grid below. Give a mark from 0 to 4 for each criterion and one piece of advice (\emph{un conseil}).
\end{enumerate}

\begin{center}
\begin{tabularx}{\linewidth}{|X|c|c|}\hline
\rowcolor{popblueL} \textbf{Criterion} & \textbf{Mark (0--4)} & \textbf{Advice} \\ \hline
Content: the problem is clear, the message is convincing & & \\ \hline
Grammar: conjunctions, pronouns, prepositions, adverbs & & \\ \hline
Vocabulary: words of citizenship, work culture, freedoms & & \\ \hline
Oral fluency: pronunciation, rhythm, confidence & & \\ \hline
Poster: slogan, clarity, creativity & & \\ \hline
\end{tabularx}
\end{center}

\courssub{Ressources à mobiliser}

\begin{aretenir}[Toolbox for the campaign]
\begin{itemize}
\item \textbf{Contrast conjunctions:} \emph{although / even though} + clause; \emph{despite / in spite of} + noun or \emph{-ing}; \emph{however, nevertheless, still, yet} to oppose two ideas. Example: ``\emph{Although} voting takes time, it changes our school.''
\item \textbf{Pronouns:} \emph{each} (every single person, singular verb), \emph{each other} (reciprocal, two people), \emph{one another} (reciprocal, more than two). Example: ``Candidates must respect \emph{one another}.''
\item \textbf{Prepositions of time:} \emph{at} + hour (\emph{at 8 a.m.}), \emph{on} + day/date (\emph{on Friday, on 12 June}), \emph{in} + month/year/part of day (\emph{in June, in the morning}), \emph{during} + noun (\emph{during the campaign}). \textbf{Prepositions of place:} \emph{in} (inside), \emph{on} (surface), \emph{at} (point), \emph{next to}, \emph{in front of}, \emph{behind}, \emph{near}.
\item \textbf{Adverbs of negation:} \emph{never, hardly, scarcely} (negative meaning: do not add another ``not''). Example: ``We \emph{never} clean the yard.'' / ``We \emph{hardly} see any change.''
\item \textbf{Phrasal verbs:} \emph{stand up for} (défendre), \emph{look after} (prendre soin de), \emph{pick up} (ramasser), \emph{put forward} (proposer), \emph{carry out} (mener à bien), \emph{take part in} (participer à).
\item \textbf{Comparative adverbs:} \emph{more actively, better, harder, more peacefully than}.
\end{itemize}
\end{aretenir}

\courssub{Démarche et corrigé commenté}

\begin{exoresolu}[Model answer -- Group ``Clean Vote'']
\textbf{Énoncé.} Produce the complete group production: document analysis, problem identification, 150-word speech and slogan, respecting all the language constraints.

\tcblower

\textbf{Step 1 -- Document analysis.} (a) The problems are: low participation in elections, a dirty yard, and intolerance between candidates. (b) Last year, $1\,200 - 480 = 720$ students did not vote, that is 60\,\%. (c) The elected delegates obtained a new water tap and repaired fans for three classrooms: this proves that voting is useful. (d) Twelve quarrels show a lack of tolerance; democracy needs peaceful debate, not violence.

\textbf{Step 2 -- Problem identification.} Our group chooses \emph{low participation in the delegate elections}, because it is the root of the other problems: \emph{although} students complain, they \emph{never} vote, so nothing changes.

\textbf{Step 3 -- Model speech (about 150 words).}

\emph{``Dear students of G.T.H.S. Nkol-Eton,}

\emph{Last year, only 40\,\% of us voted \textbf{in} the delegate elections. \textbf{Although} many students complain about dirty classrooms and broken fans, they \textbf{never} vote. \textbf{However}, voting works: our delegates obtained a new water tap \textbf{in front of} Block B.}

\emph{\textbf{Each} of us has a voice, and we must respect \textbf{one another} \textbf{during} the campaign. Candidates should debate peacefully; they should not insult \textbf{each other}. \textbf{Despite} our differences, we are one school.}

\emph{\textbf{Nevertheless}, words are not enough. Let us \textbf{take part in} the elections, \textbf{stand up for} our rights and \textbf{look after} our yard. Let us \textbf{pick up} the rubbish \textbf{next to} the dustbin and work harder than before.}

\emph{\textbf{On} Friday, vote! Together, we can build a clean, tolerant and excellent school.''}

(150 words; constraints verified: 4 contrast conjunctions -- \emph{although, however, nevertheless, despite}; 3 pronouns -- \emph{each, one another, each other}; 5 prepositions -- \emph{in, in front of, during, next to, on}; 1 adverb of negation -- \emph{never}; 4 phrasal verbs -- \emph{take part in, stand up for, look after, pick up}.)

\textbf{Step 4 -- Oral presentation.} The speaker reads slowly, stresses the key words (\emph{vote, respect, together}) and looks at the audience. The group answers one question from the class.

\textbf{Step 5 -- Slogan.} ``\emph{Your vote, your school, your future!}'' (6 words, with a comma rhythm easy to remember.)

\textbf{Step 6 -- Peer evaluation.} Example of feedback: Content 4/4; Grammar 3/4 (one wrong preposition: \emph{in Friday} $\rightarrow$ \emph{on Friday}); Vocabulary 4/4; Fluency 3/4 (speak louder); Poster 4/4. Total: 18/20.
\end{exoresolu}

\begin{attention}[Common mistakes to avoid]
\begin{itemize}
\item Do not use two negatives: say ``We \emph{never} vote'' or ``We do \emph{not} vote'', never ``We don't never vote''.
\item After \emph{despite / in spite of}, use a noun or \emph{-ing}, not a full clause: ``\emph{Despite} the rain'', not ``Despite it rains''.
\item \emph{Each} takes a singular verb: ``\emph{Each} student \emph{has} a voice.''
\item Prepositions of time: \emph{at} + hour, \emph{on} + day, \emph{in} + month or year.
\item \emph{Hardly} and \emph{scarcely} are negative in meaning; never write ``not hardly'' or ``not scarcely''.
\item Phrasal verbs are inseparable in this context: \emph{take part in}, not \emph{take in part}.
\end{itemize}
\end{attention}

\courssub{Bilan de l'activité}

This activity has shown that grammar and vocabulary are not isolated rules: they are tools to act as citizens (\emph{citoyens}). By analysing real documents, the students practised reading figures and identifying a social problem; by writing the speech, they reused contrast conjunctions, pronouns, prepositions, adverbs of negation and phrasal verbs in a meaningful context; by speaking in public and evaluating their peers, they developed confidence, tolerance and critical thinking (\emph{esprit critique}). The remediation phase targets the errors observed in the grids -- mainly double negation, the structure after \emph{despite}, and prepositions of time -- through short corrective exercises before the final evaluation. The campaign also proves that democracy begins at school: voting, respecting one another and working with excellence are concrete behaviours that each student can adopt every day.

\newpage

\coursec{Méthodes et exercices résolus}

\coursec{Lesson 04: Citizenship and Human Rights}

\courssub{Methods and Solved Exercises}

% ============================================================
% PAIR 1 : SPEAKING & GRAMMAR (CONJUNCTIONS)
% ============================================================
\begin{methode}[Expressing Work Culture and Excellence using Conjunctions]
    \textbf{Objective:} Link ideas logically when speaking about professional attitudes (cause, contrast, addition, condition).
    \begin{enumerate}
        \item \textbf{Identify the relationship} between the two ideas you want to connect.
              \begin{itemize}
                  \item \cle{Addition}: \emph{and, also, moreover, furthermore, in addition}
                  \item \cle{Contrast}: \emph{but, yet, although, even though, despite, in spite of, however, nevertheless, still}
                  \item \cle{Cause/Reason}: \emph{because, since, for, as, due to, owing to}
                  \item \cle{Condition}: \emph{if, unless, provided that}
                  \item \cle{Purpose}: \emph{to, in order to, so that}
              \end{itemize}
        \item \textbf{Choose the correct structure}:
              \begin{itemize}
                  \item \emph{Although / Even though} + \textbf{Subject + Verb} (\textit{Although he was tired, he finished the report.})
                  \item \emph{Despite / In spite of} + \textbf{Noun / Gerund (-ing)} (\textit{Despite his tiredness / Despite being tired, he finished.})
                  \item \emph{Since / As / For} often start a sentence or come mid-sentence for reason. \emph{For} is formal/literary, never starts a sentence in modern English.
                  \item \emph{However / Nevertheless / Moreover / Furthermore} are \cle{conjunctive adverbs}. Use a semicolon (;) or full stop (.) before them and a comma (,) after. (\textit{He was late; however, he apologized.})
              \end{itemize}
        \item \textbf{Practise aloud}: Simulate a job interview or a team meeting. Use at least 3 different conjunction types per minute.
        \item \textbf{Self-correct}: Record yourself. Check if \emph{although} is followed by a clause and \emph{despite} by a noun/gerund.
    \end{enumerate}
    \textbf{Key Vocabulary (Vocabulaire clé)}: \textit{work ethic (éthique professionnelle), punctuality (ponctualité), accountability (responsabilité/reddition de comptes), teamwork (travail d'équipe), deadline (échéance), excellence (excellence), integrity (intégrité), proactive (proactif).}
\end{methode}

\begin{exoresolu}[Speaking Simulation: The Ideal Employee]
    \textbf{Task:} Complete the dialogue between a Supervisor (Mr. Ngono) and an Intern (Marie) using the correct conjunction from the box. Then, practise the dialogue with a partner.
    \vspace{0.3cm}
    \begin{center}
        \begin{tabularx}{\linewidth}{|l|X|}\hline
            \textbf{Box:} & \emph{although, despite, since, however, moreover, in order to, unless} \\ \hline
        \end{tabularx}
    \end{center}
    \vspace{0.3cm}
    \textbf{Dialogue:}
    \begin{enumerate}
        \item Mr. Ngono: Marie, you submitted the report early \_\_\_\_\_\_\_\_\_\_ the deadline was Friday.
        \item Marie: \_\_\_\_\_\_\_\_\_\_ I want to build a reputation for reliability, I prioritize tasks.
        \item Mr. Ngono: \_\_\_\_\_\_\_\_\_\_ your young age, you show great maturity. \_\_\_\_\_\_\_\_\_\_, you communicate well with the team.
        \item Marie: Thank you, sir. I believe we \textbf{fail} \_\_\_\_\_\_\_\_\_\_ we support \emph{each other}.
        \item Mr. Ngono: Exactly. Maintain this attitude \_\_\_\_\_\_\_\_\_\_ go far.
    \end{enumerate}
    \tcblower
    \textbf{Solution détaillée :}
    \begin{enumerate}
        \item \textbf{although} (or \emph{even though}). \textit{Reason:} Introduces a contrast clause (Subject + Verb: "the deadline was Friday"). \textit{Fr: Bien que.}
        \item \textbf{Since} (or \emph{As / Because}). \textit{Reason:} Introduces a reason/cause clause at the start of the sentence. \textit{Fr: Puisque / Comme.}
        \item \textbf{Despite} (or \emph{In spite of}). \textit{Reason:} Followed by a noun phrase ("your young age"), not a full clause. \textit{Fr: Malgré.}
        \item \textbf{Moreover} (or \emph{Furthermore / In addition}). \textit{Reason:} Adds a new positive point. Starts a new sentence, followed by a comma. \textit{Fr: De plus / En outre.}
        \item \textbf{unless}. \textit{Reason:} Introduces a negative condition ("if not"). "We fail \textbf{unless} we support..." = "We fail if we do not support...". \textit{Fr: À moins que / Si... pas.}
        \item \textbf{in order to}. \textit{Reason:} Expresses purpose, followed by base verb ("go"). Structure: Imperative + \emph{in order to} + base verb. \textit{Fr: Afin de / Pour.}
    \end{enumerate}
    \textbf{Corrected Dialogue:}
    \begin{quote}
    Mr. Ngono: Marie, you submitted the report early \textbf{although} the deadline was Friday.\\
    Marie: \textbf{Since} I want to build a reputation for reliability, I prioritize tasks.\\
    Mr. Ngono: \textbf{Despite} your young age, you show great maturity. \textbf{Moreover}, you communicate well with the team.\\
    Marie: Thank you, sir. I believe we \textbf{fail unless} we support \emph{each other}.\\
    Mr. Ngono: Exactly. Maintain this attitude \textbf{in order to} go far.
    \end{quote}
    \textbf{Conclusion:} Mastering the syntax of conjunctions (clause vs. phrase) is essential for fluent professional speech.
\end{exoresolu}

% ============================================================
% PAIR 2 : LISTENING & GRAMMAR (RECIPROCAL PRONOUNS)
% ============================================================
\begin{methode}[Using Reciprocal Pronouns for Civic Responsibility]
    \textbf{Objective:} Correctly use \cle{each other} and \cle{one another} to express mutual actions, and \cle{each} as a determiner/pronoun.
    \begin{enumerate}
        \item \textbf{Distinguish the forms}:
              \begin{itemize}
                  \item \textbf{Each other} / \textbf{One another}: \cle{Reciprocal pronouns}. Action goes both ways between two+ people. \textit{They respect \textbf{each other}.} (Standard for 2 or more). \textit{The community members help \textbf{one another}.} (Often preferred for groups/general statements).
                  \item \textbf{Each}: \cle{Distributive}. Focuses on individuals separately. \textit{\textbf{Each} citizen has a duty.} (Determiner + singular noun). \textit{\textbf{Each of us} must vote.} (Pronoun + \emph{of} + plural pronoun/noun, verb usually singular).
              \end{itemize}
        \item \textbf{Listening Strategy}: Identify keywords signaling mutual obligation (\textit{mutual, shared, together, community, solidarity, respect, help, support}).
        \item \textbf{Grammar Check}: 
              \begin{itemize}
                  \item Never say "each others" (no 's').
                  \item Possessive: \textit{each other's / one another's} (\textit{They checked each other's work.}).
                  \item \emph{Each} takes a \textbf{singular verb}. \textit{Each of the students \textbf{is} responsible.}
              \end{itemize}
    \end{enumerate}
    \textbf{Key Vocabulary}: \textit{civic duty (devoir civique), volunteer (bénévole), community service (service communautaire), solidarity (solidarité), accountability (responsabilisation), mutual respect (respect mutuel), ballot (bulletin de vote), polling station (bureau de vote).}
\end{methode}

\begin{exoresolu}[Listening Comprehension \& Grammar Gap-Fill]
    \textbf{Task:} Listen to the short talk "Building a Better Neighbourhood" (script provided below) and fill in the gaps with \emph{each other}, \emph{one another}, \emph{each}, or \emph{each of}.
    \vspace{0.2cm}
    \textbf{Script:} "Good morning. A strong community depends on mutual respect. Citizens must listen to \_\_\_\_\_\_\_\_\_\_ (1) during town hall meetings. \_\_\_\_\_\_\_\_\_\_ (2) person brings a unique view. \_\_\_\_\_\_\_\_\_\_ (3) the representatives, the mayor included, must serve the people. We should care for \_\_\_\_\_\_\_\_\_\_ (4) in times of crisis. If \_\_\_\_\_\_\_\_\_\_ (5) neighbour watches the street, we are all safer. Let us promise to support \_\_\_\_\_\_\_\_\_\_ (6)."
    \tcblower
    \textbf{Solution détaillée :}
    \begin{enumerate}
        \item \textbf{each other} (or \emph{one another}). \textit{Reason:} Reciprocal action "listen to [one another]". Subject "Citizens" (plural). \textit{Fr: les uns les autres.}
        \item \textbf{Each}. \textit{Reason:} Determiner before singular noun "person". Verb "brings" is singular. \textit{Fr: Chaque.}
        \item \textbf{Each of}. \textit{Reason:} Pronoun structure "Each of [the representatives]". Verb "must serve" (modal + base form, but logic is singular distributive). \textit{Fr: Chacun de.}
        \item \textbf{one another} (or \emph{each other}). \textit{Reason:} Reciprocal "care for [one another]" within a group/community context. \textit{Fr: les uns les autres.}
        \item \textbf{each}. \textit{Reason:} Determiner before singular noun "neighbour". "If \textbf{each} neighbour watches..." = "If every single neighbour watches...". \textit{Fr: Chaque.}
        \item \textbf{one another} (or \emph{each other}). \textit{Reason:} Reciprocal "support [one another]" - general group pledge. \textit{Fr: les uns les autres.}
    \end{enumerate}
    \textbf{Grammar Note:} In (3), "Each of the representatives... must serve" - the verb follows the modal "must" (base form), but conceptually singular. In standard grammar: "Each of them \textbf{is} here."
    \textbf{Conclusion:} Reciprocal pronouns build the language of solidarity; distributive \emph{each} emphasizes individual duty within the collective.
\end{exoresolu}

% ============================================================
% PAIR 3 : READING & GRAMMAR (PREPOSITIONS)
% ============================================================
\begin{methode}[Mastering Prepositions and Prepositional Phrases (Time, Place, Manner)]
    \textbf{Objective:} Identify and use complex prepositional phrases accurately in texts about democracy and rights.
    \begin{enumerate}
        \item \textbf{Categorize the phrase} by function:
              \begin{itemize}
                  \item \cle{Time}: \emph{at noon, on Monday, in 2025, during the election, for two hours, since morning, until midnight, by the deadline, before voting, after the debate}.
                  \item \cle{Place}: \emph{at the polling station, in the capital, on the ballot paper, under the law, between parties, among voters, near the border}.
                  \item \cle{Manner}: \emph{with integrity, by secret ballot, in a peaceful way, without fear, like a democrat, through dialogue}.
              \end{itemize}
        \item \textbf{Spot the structure}: Preposition + Noun Phrase (often with adjectives/determiners). \textit{In \textbf{free and fair} elections.}
        \item \textbf{Avoid translation traps}:
              \begin{itemize}
                  \item \textit{Participer à} $\to$ \textbf{take part \emph{in}} (not \emph{to}).
                  \item \textit{Voter pour} $\to$ \textbf{vote \emph{for}} (candidate) / \textbf{vote \emph{on}} (issue).
                  \item \textit{Respecter les droits} $\to$ \textbf{respect \emph{for}} (noun) / \textbf{respect} (verb, no prep).
              \end{itemize}
        \item \textbf{Reading Tip}: Underline all prepositional phrases in a text. Ask: "When? Where? How?" to classify them.
    \end{enumerate}
    \textbf{Key Vocabulary}: \textit{freedom of speech (liberté d'expression), rule of law (État de droit), human rights (droits de l'homme), tolerance (tolérance), pluralism (pluralisme), election (élection), candidate (candidat), ballot box (urne), constitution (constitution).}
\end{methode}

\begin{exoresolu}[Reading Analysis: "The Pillars of Democracy"]
    \textbf{Task:} Read the extract. Underline 5 prepositional phrases and classify them (Time, Place, Manner). Then answer the questions.
    \vspace{0.2cm}
    \textbf{Text:} "In a true democracy, power belongs \textbf{to the people}. Elections are held \textbf{every five years} \textbf{at designated centres}. Citizens vote \textbf{by secret ballot} \textbf{without intimidation}. The government rules \textbf{according to the constitution} \textbf{for the common good}."
    \vspace{0.2cm}
    \textbf{Questions:}
    \begin{enumerate}
        \item Identify the prepositional phrase indicating \textbf{Manner} for "vote".
        \item Identify the phrase indicating \textbf{Time} for "held".
        \item Rewrite: "The government rules according to the constitution" using a different prepositional phrase of manner (use \emph{in accordance with}).
        \item Find the phrase showing \textbf{Place}.
    \end{enumerate}
    \tcblower
    \textbf{Solution détaillée :}
    \textbf{Classification:}
    \begin{center}
        \begin{tabularx}{\linewidth}{|l|X|l|}\hline
            \textbf{Phrase} & \textbf{Function} & \textbf{Head Noun} \\ \hline
            to the people & Place (abstract destination) / Possession & people \\
            every five years & \textbf{Time} (Frequency/Duration) & years \\
            at designated centres & \textbf{Place} (Location) & centres \\
            by secret ballot & \textbf{Manner} (Method/Instrument) & ballot \\
            without intimidation & \textbf{Manner} (Condition/Absence) & intimidation \\
            according to the constitution & \textbf{Manner} (Standard/Rule) & constitution \\
            for the common good & \textbf{Manner} (Purpose/Beneficiary) & good \\ \hline
        \end{tabularx}
    \end{center}
    \vspace{0.3cm}
    \textbf{Answers:}
    \begin{enumerate}
        \item \textbf{by secret ballot} (and \emph{without intimidation}).
        \item \textbf{every five years}.
        \item "The government rules \textbf{in accordance with} the constitution."
        \item \textbf{at designated centres}.
    \end{enumerate}
    \textbf{Justification:} \emph{By + -ing/noun} shows method. \emph{Every + duration} shows frequency. \emph{At + specific location} shows place. \emph{In accordance with} is a formal multi-word preposition synonymous with \emph{according to}.
    \textbf{Conclusion:} Prepositional phrases add precision to civic discourse; confusing \emph{in/on/at} for time/place is a major error source.
\end{exoresolu}

% ============================================================
% PAIR 4 : GRAMMAR (COMPARATIVE ADVERBS \& NEGATION)
% ============================================================
\begin{methode}[Comparing Actions and Using Negative Adverbs]
    \textbf{Objective:} Form comparative adverbs correctly and use \cle{hardly}, \cle{scarcely}, \cle{never} with correct syntax (inversion).
    \begin{enumerate}
        \item \textbf{Comparative Adverbs} (Comparing \emph{how} actions happen):
              \begin{itemize}
                  \item \textbf{One syllable} (fast, hard, early, late): add \textbf{-er}. \textit{fast $\to$ faster}. \textit{He works \textbf{faster than} me.}
                  \item \textbf{-ly adverbs} (carefully, efficiently, democratically): use \textbf{more / less}. \textit{She speaks \textbf{more clearly than} him.}
                  \item \textbf{Irregulars}: \textit{well $\to$ better; badly $\to$ worse; little $\to$ less; much $\to$ more}.
                  \item \textbf{Equality}: \textbf{as + adverb + as}. \textit{They voted \textbf{as peacefully as} expected.}
              \end{itemize}
        \item \textbf{Adverbs of Negation} (\emph{hardly, scarcely, barely, never, rarely, seldom}):
              \begin{itemize}
                  \item They have \textbf{negative meaning}. \textbf{NEVER} use double negative (\textit{Incorrect: "I didn't hardly know." Correct: "I \textbf{hardly} knew."}).
                  \item \textbf{Inversion Rule (Formal/Written)}: When starting a sentence with \emph{Never, Rarely, Seldom, Hardly, Scarcely, Barely} $\to$ \textbf{Auxiliary + Subject + Verb}.
                      \begin{itemize}
                          \item \textit{Never \textbf{have I seen} such corruption.}
                          \item \textit{Hardly \textbf{did he arrive} when the meeting started.} (\emph{Hardly/Scarcely/Barely} + \textbf{Past Perfect} ... \emph{when} + \textbf{Past Simple}).
                      \end{itemize}
              \end{itemize}
    \end{enumerate}
    \textbf{Key Vocabulary}: \textit{efficiently (efficacement), transparently (transparentement), corruptly (corrompu), freely (librement), fairly (équitablement), rigorously (rigoureusement).}
\end{methode}

\begin{exoresolu}[Grammar Transformation: Civic Actions]
    \textbf{Task:} Rewrite the sentences using the word in brackets. Do not change the meaning.
    \begin{enumerate}
        \item The officials counted the votes very slowly. (quickly) $\to$ The officials counted the votes \_\_\_\_\_\_\_\_\_\_.
        \item Young people vote less frequently than older people. (often) $\to$ Young people \_\_\_\_\_\_\_\_\_\_.
        \item He almost didn't understand the law. (hardly) $\to$ He \_\_\_\_\_\_\_\_\_\_ the law.
        \item I have never seen such tolerance. (Never) $\to$ \_\_\_\_\_\_\_\_\_\_ such tolerance.
        \item She speaks more persuasively than her opponent. (as ... as) $\to$ Her opponent doesn't speak \_\_\_\_\_\_\_\_\_\_ she does.
    \end{enumerate}
    \tcblower
    \textbf{Solution détaillée :}
    \begin{enumerate}
        \item \textbf{more quickly} (or \emph{quicker} - informal). \textit{Reason:} "Quickly" is -ly adverb $\to$ \emph{more quickly}. Comparative structure. \textit{Fr: plus rapidement.}
        \item \textbf{don't vote as often as} older people. \textit{Reason:} "Less frequently" = negative comparison of equality. "Not as often as". \textit{Fr: ne votent pas aussi souvent que.}
        \item \textbf{hardly understood}. \textit{Reason:} \emph{Hardly} = almost not. Negative meaning $\to$ positive verb (\emph{understood}, not \emph{didn't understand}). Past tense narrative. \textit{Fr: a à peine compris.}
        \item \textbf{Never have I seen}. \textit{Reason:} Fronted negative adverb \emph{Never} triggers \textbf{Subject-Auxiliary Inversion} (Present Perfect: have I seen). Formal style. \textit{Fr: Jamais je n'ai vu.}
        \item \textbf{as persuasively as}. \textit{Reason:} Transforming comparative (\emph{more ... than}) to negative equality (\emph{not as ... as}). Adverb \emph{persuasively} kept. \textit{Fr: aussi persuasivement que.}
    \end{enumerate}
    \textbf{Conclusion:} Inversion after negative adverbs is a hallmark of advanced writing (Bac level). Double negatives (\emph{not hardly}) are strictly forbidden.
\end{exoresolu}

% ============================================================
% PAIR 5 : GRAMMAR (PHRASAL VERBS)
% ============================================================
\begin{methode}[Identifying and Using Phrasal Verbs for Civic Contexts]
    \textbf{Objective:} Recognize separable/inseparable phrasal verbs and guess meaning from context.
    \begin{enumerate}
        \item \textbf{Structure}: Verb + Particle (preposition/adverb). Meaning is often \cle{idiomatic} (not literal).
        \item \textbf{Types}:
              \begin{itemize}
                  \item \textbf{Intransitive} (No object): \textit{The meeting \textbf{broke up} early. / Please \textbf{speak up}.}
                  \item \textbf{Transitive Separable} (Object can go between or after): \textit{\textbf{Bring up} the topic / \textbf{Bring} the topic \textbf{up}.} \textbf{Rule:} Pronoun object \textbf{MUST} go in middle (\textit{Bring \textbf{it} up}, not *Bring up it).
                  \item \textbf{Transitive Inseparable} (Object \textbf{only} after): \textit{\textbf{Look into} the matter.} (Not *Look the matter into). \textit{\textbf{Stand for} (represent/tolerate).}
                  \item \textbf{Three-word verbs} (Verb + Part + Prep): \textit{\textbf{Look forward to}, \textbf{put up with}, \textbf{come up with}.} Always inseparable.
              \end{itemize}
        \item \textbf{Strategy for Exams}:
              \begin{enumerate}
                  \item Identify the verb + particle.
                  \item Check if there is an object (noun/pronoun).
                  \item Test separability: Try putting a pronoun (\emph{it, them}) in the middle. If it works $\to$ Separable. If impossible $\to$ Inseparable.
                  \item Use context (democracy, rights, work) to deduce meaning.
              \end{enumerate}
    \end{enumerate}
    \textbf{Key Phrasal Verbs (Unit Vocabulary)}:
    \begin{itemize}
        \item \textbf{Stand up for} (defend/support rights) - Inseparable.
        \item \textbf{Speak out} (express opinion publicly) - Intransitive.
        \item \textbf{Look into} (investigate) - Inseparable.
        \item \textbf{Bring about} (cause change) - Separable.
        \item \textbf{Give in} (surrender/yield) - Intransitive.
        \item \textbf{Carry out} (execute/implement) - Separable.
        \item \textbf{Put up with} (tolerate) - Inseparable (3 words).
        \item \textbf{Turn out} (attend/result) - Intransitive.
    \end{itemize}
\end{methode}

\begin{exoresolu}[Phrasal Verbs in Context: Civic Engagement]
    \textbf{Task:} Complete sentences with the correct form of the phrasal verbs from the box. Indicate if Separable (S) or Inseparable (I).
    \vspace{0.2cm}
    \begin{center}
        \begin{tabularx}{\linewidth}{|X|}\hline
            \emph{stand up for, look into, bring about, carry out, speak out, put up with, turn out, give in} \\ \hline
        \end{tabularx}
    \end{center}
    \begin{enumerate}
        \item Citizens must \_\_\_\_\_\_\_\_\_\_ corruption in the workplace. (S/I?)
        \item The committee will \_\_\_\_\_\_\_\_\_\_ the allegations of fraud. (S/I?)
        \item We cannot \_\_\_\_\_\_\_\_\_\_ injustice any longer. (S/I?)
        \item The new law aims to \_\_\_\_\_\_\_\_\_\_ social change. (S/I?) $\to$ \_\_\_\_\_\_\_\_\_\_ \textbf{it} \_\_\_\_\_\_\_\_\_\_.
        \item Many voters \_\_\_\_\_\_\_\_\_\_ \_\_\_\_\_\_\_\_\_\_ for the referendum. (S/I?)
        \item The union refused to \_\_\_\_\_\_\_\_\_\_ to pressure. (S/I?)
    \end{enumerate}
    \tcblower
    \textbf{Solution détaillée :}
    \begin{enumerate}
        \item \textbf{stand up for} \textbf{(I)}. \textit{Meaning:} Defend. \textit{Test:} "Stand up for \textbf{it}" (OK). "Stand \textbf{it} up for" (Impossible). Preposition \emph{for} needs object.
        \item \textbf{look into} \textbf{(I)}. \textit{Meaning:} Investigate. \textit{Test:} "Look into \textbf{it}" (OK). "Look \textbf{it} into" (Wrong).
        \item \textbf{put up with} \textbf{(I)}. \textit{Meaning:} Tolerate. Three-word verb. \textit{Test:} "Put up with \textbf{it}".
        \item \textbf{bring about} \textbf{(S)}. \textit{Meaning:} Cause. \textit{Test with pronoun:} "Bring \textbf{it} about." (Correct). "Bring about \textbf{it}" (Awkward/Poetic). Rewrite: \textbf{Bring it about}.
        \item \textbf{turned out}. \textbf{(I)}. \textit{Meaning:} Attended / Resulted. Intransitive. Past tense narrative.
        \item \textbf{give in}. \textbf{(I)}. \textit{Meaning:} Surrender. Intransitive. \textit{Refused to give in.}
    \end{enumerate}
    \textbf{Conclusion:} The "Pronoun Test" (\emph{it/them}) is the fastest way to determine separability in exams. Learn the 3-word verbs (\emph{put up with, look forward to}) as fixed blocks.
\end{exoresolu}

% ============================================================
% PAIR 6 : WRITING (COMPOSITION)
% ============================================================
\begin{methode}[Writing an Argumentative Composition on Democracy and Freedoms]
    \textbf{Objective:} Structure a coherent 200-250 word essay for the Probatoire/Bac.
    \begin{enumerate}
        \item \textbf{Analyse the Prompt}: Identify \textbf{Topic} (Democracy, Tolerance, Freedoms), \textbf{Task} (Discuss, Advantages/Disadvantages, Solutions, Opinion), \textbf{Format} (Essay, Letter, Article).
        \item \textbf{Plan (5-10 min)}: Use \cle{Connector Map}.
              \begin{itemize}
                  \item \textbf{Introduction} ($\approx$ 40 words): Hook (Quote/Stat/Fact) $\to$ Define key terms $\to$ \textbf{Thesis Statement} (Your clear position/roadmap).
                  \item \textbf{Body Paragraph 1} ($\approx$ 70 words): \textbf{Argument 1} (Topic Sentence) $\to$ Explanation/Example $\to$ Connector $\to$ Argument 2 $\to$ Example. \textit{Connectors: Furthermore, Moreover, In addition, For instance.}
                  \item \textbf{Body Paragraph 2} ($\approx$ 70 words): \textbf{Counter-argument / Contrast} $\to$ Refutation. \textit{Connectors: However, Although, On the other hand, Despite.}
                  \item \textbf{Conclusion} ($\approx$ 30 words): \textbf{Restate Thesis} (different words) $\to$ Summary of main points $\to$ Final thought / Call to action / Future implication. \textit{Connectors: In conclusion, To sum up, Ultimately.}
              \end{itemize}
        \item \textbf{Draft}: Write quickly. Focus on \cle{Grammar Targets}: Conjunctions, Passive voice (\textit{Rights are guaranteed}), Conditionals (\textit{If we tolerate...}), Phrasal verbs (\textit{stand up for}).
        \item \textbf{Proofread (5 min)}: Check \textbf{Subject-Verb Agreement}, \textbf{Tense Consistency}, \textbf{Prepositions}, \textbf{Spelling}, \textbf{Word Count}.
    \end{enumerate}
    \textbf{Checklist Bac:} \fbox{Thesis?} \fbox{Connectors?} \fbox{Vocab (Unit words)?} \fbox{Grammar (Conjunctions, Preps, Adverbs)?} \fbox{Paragraphing?} \fbox{Word count?}
\end{methode}

\begin{exoresolu}[Model Composition: Tolerance and Freedom]
    \textbf{Topic:} "Tolerance is the foundation of a free society." Write an essay of about 200 words discussing this statement. Illustrate with examples from your country (Cameroon).
    \tcblower
    \textbf{Model Essay:}
    \vspace{0.2cm}
    \noindent \textbf{Tolerance: The Pillar of Liberty}
    \vspace{0.2cm}
    \noindent Voltaire famously said, "I disapprove of what you say, but I will defend to the death your right to say it." This quote captures the essence of tolerance: accepting differences \textbf{in order to} preserve freedom. In Cameroon, a nation of over 250 ethnic groups, tolerance is not merely a virtue but a necessity for national unity.
    \vspace{0.2cm}
    \noindent \textbf{Firstly}, tolerance allows democracy to function. \textbf{Without} the acceptance of opposing views, elections become battles rather than choices. \textbf{For instance}, during the 2018 presidential election, peaceful transitions relied on candidates \textbf{speaking out} respectfully and supporters \textbf{putting up with} results they disliked. \textbf{Moreover}, freedom of expression thrives only when citizens \textbf{stand up for} the rights of those they disagree with. \textbf{If} we silence minority voices, we erode the \cle{rule of law} for everyone.
    \vspace{0.2cm}
    \noindent \textbf{However}, tolerance has limits. \textbf{Although} we must respect beliefs, we \textbf{never} tolerate hate speech or violence. \textbf{Despite} constitutional guarantees, some groups \textbf{hardly} respect the rights of others. The government must \textbf{look into} such violations \textbf{and} \textbf{bring about} justice swiftly.
    \vspace{0.2cm}
    \noindent \textbf{In conclusion}, tolerance is the glue holding our diverse society together. It demands active effort, not passive indifference. \textbf{Unless} every Cameroonian \textbf{carries out} their civic duty to respect difference, our freedoms remain fragile. Let us \textbf{bring up} the next generation to value unity in diversity.
    \vspace{0.3cm}
    \textbf{Analysis of Target Structures Used:}
    \begin{itemize}
        \item \textbf{Conjunctions}: \emph{in order to, without, for instance, moreover, if, however, although, despite, and, unless}.
        \item \textbf{Prepositional Phrases}: \emph{In Cameroon, during the election, for everyone, with limits, In conclusion}.
        \item \textbf{Comparative/Negation}: \emph{never tolerate, hardly respect}.
        \item \textbf{Phrasal Verbs}: \emph{speaking out, putting up with, stand up for, look into, bring about, carry out, bring up}.
        \item \textbf{Vocabulary}: \cle{ethnic groups, national unity, peaceful transitions, rule of law, hate speech, civic duty, unity in diversity}.
    \end{itemize}
    \textbf{Word Count:} $\approx$ 215 words. \textbf{Grade Target:} 18-20/20.
\end{exoresolu}

% ============================================================
% PAIR 7 : INTEGRATION (MIXED GRAMMAR \& VOCAB)
% ============================================================
\begin{methode}[Integration Strategy: Mixed Grammar Editing for Exam Readiness]
    \textbf{Objective:} Simulate the "Integration Activities / Evaluation" section. Combine all unit grammar points in an error-correction / gap-fill task.
    \begin{enumerate}
        \item \textbf{Scan for Triggers}:
              \begin{itemize}
                  \item \emph{Although / Despite} $\to$ Check Clause vs Noun Phrase.
                  \item \emph{Each / Each of} $\to$ Check Singular Verb.
                  \item \emph{Each other / One another} $\to$ Check Reciprocal context.
                  \item \emph{In / On / At / By / With} $\to$ Check Time/Place/Manner collocation.
                  \item \emph{Hardly / Never / Scarcely} $\to$ Check No Double Negative + Inversion (if fronted).
                  \item \emph{Phrasal Verb + Pronoun} $\to$ Check Separability (Middle position).
              \end{itemize}
        \item \textbf{Systematic Editing Pass}:
              \begin{enumerate}
                  \item Pass 1: Verbs (Tense, Agreement, Form).
                  \item Pass 2: Connectors/Conjunctions (Logic, Punctuation).
                  \item Pass 3: Prepositions (Dependent prepositions: \emph{interested in, good at, responsible for}).
                  \item Pass 4: Vocabulary Precision (Collocations: \emph{make a decision, not do}).
              \end{enumerate}
    \end{enumerate}
\end{methode}

\begin{exoresolu}[Integration Test: Error Correction \& Gap Fill]
    \textbf{Task A:} Correct the errors in the following sentences (Grammar/Vocab).
    \begin{enumerate}
        \item Although the rain, the voters queued patiently.
        \item Each of the candidate have their own manifesto.
        \item The delegates discussed about the new law.
        \item We must stand up the rights of minorities.
        \item He didn't hardly know the procedure.
        \item The committee brought the new regulations up. (Check pronoun placement).
    \end{enumerate}
    \textbf{Task B:} Fill gaps with ONE word (Preposition / Conjunction / Adverb / Particle).
    "Democracy thrives \_\_\_ (1) citizens participate actively. \_\_\_ (2) voting every five years, they must hold leaders accountable \_\_\_ (3) the year. Leaders should listen \_\_\_ (4) the people and not give \_\_\_ (5) to corruption. \_\_\_ (6) citizen has a role. We must respect \_\_\_ (7) \_\_\_ (8) differences."
    \tcblower
    \textbf{Solution détaillée :}
    \textbf{Task A Corrections:}
    \begin{enumerate}
        \item \textbf{Despite the rain,} (or \emph{Although it rained,}). \textit{Error:} \emph{Although} + Clause needed. \emph{Despite} + Noun Phrase.
        \item \textbf{Each of the candidates has his/her own manifesto.} \textit{Errors:} \emph{candidate} $\to$ plural after \emph{Each of}; \emph{have} $\to$ singular \emph{has}; \emph{their} $\to$ singular \emph{his/her} (or singular \emph{their} accepted in modern usage, but \emph{has} is key).
        \item \textbf{The delegates discussed the new law.} \textit{Error:} \emph{Discuss} is transitive (no preposition). \emph{Discuss about} is wrong. (But \emph{talk about} / \emph{debate}).
        \item \textbf{We must stand up \emph{for} the rights...} \textit{Error:} \emph{Stand up for} (inseparable 3-word verb). Missing particle \emph{for}.
        \item \textbf{He hardly knew the procedure.} \textit{Error:} Double negative \emph{didn't hardly}. \emph{Hardly} is negative itself.
        \item \textbf{The committee brought \emph{them} up.} \textit{Correction:} Separable phrasal verb. Noun object "the new regulations" can go after, but Pronoun object \textbf{must} go in middle: \emph{Brought them up}.
    \end{enumerate}
    \textbf{Task B Gap Fill:}
    \begin{enumerate}
        \item \textbf{when} / \textbf{if} / \textbf{whenever} (Time conjunction).
        \item \textbf{Besides} / \textbf{In addition to} / \textbf{Apart from} (Prepositional phrase addition). Or \emph{Besides} (preposition).
        \item \textbf{throughout} / \textbf{during} / \textbf{all} (Preposition time/duration).
        \item \textbf{to} (Dependent prep: \emph{listen to}).
        \item \textbf{in} (Phrasal verb: \emph{give in to} - gap is particle: \textbf{in}. "give \textbf{in} to corruption").
        \item \textbf{Each} / \textbf{Every} (Distributive determiner, singular verb \emph{has}).
        \item \textbf{each} (Reciprocal pronoun part 1).
        \item \textbf{other} (Reciprocal pronoun part 2: \emph{each other}).
    \end{enumerate}
    \textbf{Conclusion:} Integration exercises test the ability to switch rapidly between grammar modules. Accuracy on "small words" (prepositions, particles, conjunctions) distinguishes top grades.
\end{exoresolu}

% ============================================================
% ATTENTION BOX
% ============================================================
\begin{attention}[Les 5 erreurs qui coûtent des points au Bac]
    \begin{enumerate}
        \item \textbf{Confusion \emph{Although} vs \emph{Despite/In spite of}.} \\ 
            \textit{Incorrect:} "Although his injury, he played." \\
            \textit{Correct:} "\textbf{Despite} his injury..." OR "Although \textbf{he was injured}..." \\
            \textbf{Reflexe:} \emph{Although} = Sujet + Verbe ; \emph{Despite} = Nom / Gérondif.
        \item \textbf{Double négation avec \emph{Hardly / Scarcely / Never}.} \\
            \textit{Incorrect:} "I \textbf{didn't hardly} see him." / "He \textbf{doesn't never} come." \\
            \textit{Correct:} "I \textbf{hardly} saw him." / "He \textbf{never} comes." \\
            \textbf{Reflexe:} Un seul négatif par proposition. Si l'adverbe est en tête $\to$ Inversion (\emph{Never have I...}).
        \item \textbf{Mauvaise placement du pronom avec \emph{Phrasal Verbs} séparables.} \\
            \textit{Incorrect:} "Turn off \textbf{it}." / "Bring up \textbf{them}." \\
            \textit{Correct:} "Turn \textbf{it} off." / "Bring \textbf{them} up." \\
            \textbf{Reflexe:} Objet pronominal $\to$ \textbf{OBLIGATOIREMENT} au milieu.
        \item \textbf{Accord du verbe après \emph{Each of} / \emph{Every} / \emph{Each}.} \\
            \textit{Incorrect:} "Each of the students \textbf{are} here." / "Every candidate \textbf{have} a card." \\
            \textit{Correct:} "Each of the students \textbf{is} here." / "Every candidate \textbf{has} a card." \\
            \textbf{Reflexe:} \emph{Each/Every} = Singulier. \emph{Each of} + Pluriel + Verbe Singulier.
        \item \textbf{Prépositions dépendantes oubliées ou erronées (Collocations).} \\
            \textit{Incorrect:} "Participate \textbf{to} / Discuss \textbf{about} / Interested \textbf{on} / Responsible \textbf{of}." \\
            \textit{Correct:} "Participate \textbf{in} / Discuss \textbf{\_\_\_} (direct) / Interested \textbf{in} / Responsible \textbf{for}." \\
            \textbf{Reflexe:} Apprendre les blocs : \emph{listen to, wait for, look at, look into, stand up for, vote for/on}.
    \end{enumerate}
\end{attention}

\newpage

\coursec{Exercices}

\courssub{Je vérifie mes connaissances}

\begin{aretenir}[Les conjonctions de coordination et de subordination -- Rappel]
\textbf{Conjonctions de subordination} : introduisent une proposition subordonnée.
\begin{itemize}
\item \textbf{Since} / \textbf{As} / \textbf{Because} : cause (puisque, comme, parce que). \emph{Fr: puisque / comme / parce que}
\item \textbf{Although} / \textbf{Though} / \textbf{Even though} : concession (bien que). \emph{Fr: bien que / quoique}
\item \textbf{Since} / \textbf{For} : aussi temps (depuis que / car).
\end{itemize}
\textbf{Adverbes de liaison / Conjonctions de coordination} :
\begin{itemize}
\item \textbf{However} / \textbf{Nevertheless} / \textbf{Still} / \textbf{Yet} : opposition (cependant, néanmoins, pourtant). \emph{Fr: cependant / néanmoins / pourtant}
\item \textbf{Therefore} / \textbf{Consequently} / \textbf{Thus} : conséquence (donc, par conséquent).
\end{itemize}
\textbf{Locutions prépositionnelles} :
\begin{itemize}
\item \textbf{Despite} / \textbf{In spite of} + nom / gérondif (malgré). \emph{Fr: malgré}
\end{itemize}
\emph{Attention :} \textbf{Although} + sujet + verbe ; \textbf{Despite / In spite of} + nom / \emph{-ing}. Ne pas dire \emph{Although... but...} (redondance).
\end{aretenir}

\begin{exercice}[QCM: Conjunctions and Prepositions]
Choose the correct answer (A, B, C or D).
\begin{enumerate}
\item \ldots the rain was heavy, the workers continued building the road in Bamenda.\\
A. Despite \quad B. Although \quad C. However \quad D. Since
\item The team worked hard; \ldots, they did not finish the project on time.\\
A. although \quad B. because \quad C. however \quad D. since
\item The students have been cleaning the school yard \ldots two hours.\\
A. since \quad B. at \quad C. for \quad D. in
\item \ldots his lack of experience, the young mechanic repaired the engine perfectly.\\
A. In spite of \quad B. Although \quad C. Yet \quad D. However
\end{enumerate}
\end{exercice}

\begin{exercice}[True or False -- Justify]
Say whether the following statements are True or False, and justify your answer with the grammar rule.
\begin{enumerate}
\item ``Although the voters were many, but they voted peacefully.'' This sentence is correct.
\item ``Each other'' is used for two people, while ``one another'' is traditionally used for more than two.
\item ``Hardly'' and ``scarcely'' are negative adverbs, so we must not add another negative word in the same sentence.
\item We say ``good at work'' and ``interested on politics''.
\end{enumerate}
\end{exercice}

\begin{exercice}[Definitions]
Give a short definition in English of the following words and expressions related to citizenship and work culture.
\begin{enumerate}
\item civic responsibility
\item tolerance
\item excellence at work
\item basic freedoms
\end{enumerate}
\end{exercice}

\begin{exercice}[Comparative Adverbs -- Direct Application]
Put the adverb in brackets in the correct comparative or superlative form.
\begin{enumerate}
\item Amina speaks English \ldots (fluently) than her brother.
\item Of all the workers, Moussa works \ldots (hard).
\item The new secretary types \ldots (quickly) than the old one.
\item Our team arrived \ldots (early) of all the teams at the workshop.
\end{enumerate}
\end{exercice}

\courssub{Je m'entraîne}

\begin{methode}[Stratégie: Choisir la bonne conjonction]
\begin{enumerate}
\item Identifier la relation logique: cause (\emph{because, since, as}), concession (\emph{although, despite/in spite of}), opposition (\emph{however, yet, nevertheless}), temps (\emph{since, for}).
\item Vérifier la structure: \emph{Although} + S + V ; \emph{Despite/In spite of} + N / V-\emph{ing} ; \emph{However} + , + S + V (après . ou ;).
\item \emph{Since} = point de départ (date) ; \emph{For} = durée (période).
\end{enumerate}
\end{methode}

\begin{exercice}[Fill in the Blanks: Conjunctions]
Complete the sentences with the correct conjunction or linking expression: \emph{although, despite, however, since, yet, nevertheless, in spite of, still, for, because}.
\begin{enumerate}
\item \ldots the workers were tired, they finished the bridge on time.
\item \ldots the strike, the nurses of the Yaoundé Central Hospital continued to help patients.
\item The mayor of Douala promised new markets; \ldots, nothing has been done so far.
\item The mechanic has worked in this garage \ldots 2015.
\item The farmers of the West Region worked hard; \ldots, the harvest was poor because of the drought.
\item \ldots of the difficulties, the students of G.T.H.S. Bafoussam passed their practical exam.
\item It was raining heavily; \ldots, the football match went on.
\item \ldots Cameroon is a bilingual country, every citizen should learn both English and French.
\item The elections were peaceful; \ldots, some young people did not register on the electoral lists.
\item She is very busy at the workshop; \ldots, she always finds time to vote.
\end{enumerate}
\end{exercice}

\begin{attention}[Pièges fréquents: Each other vs One another]
\begin{itemize}
\item \textbf{Each other} = relation réciproque entre \textbf{deux} personnes/groupes. \emph{Fr: l'un l'autre}
\item \textbf{One another} = relation réciproque entre \textbf{plus de deux} personnes/groupes. \emph{Fr: les uns les autres}
\item \textbf{Erreur courante:} \emph{each others} (pas de \emph{s} à \emph{other}).
\item \textbf{Note:} En anglais moderne, l'usage tend à les rendre interchangeables, mais l'examen technique attend la distinction traditionnelle.
\end{itemize}
\end{attention}

\begin{exercice}[Each Other / One Another]
Rewrite the following sentences using \emph{each other} or \emph{one another}, and justify your choice.
\begin{enumerate}
\item The two candidates respect their opponent. $\rightarrow$ The two candidates respect \ldots
\item The members of the cooperative help all the members.
\item My friend and I write letters to my friend.
\item The students of the class lend books to all their classmates.
\item The husband and the wife trust the partner.
\item The five neighbours clean the quarter with all the neighbours.
\end{enumerate}
\end{exercice}

\begin{aretenir}[Phrasal Verbs clés -- Citoyenneté et Travail]
\begin{center}
\begin{tabularx}{\linewidth}{|l|X|X|}\hline
\textbf{Phrasal Verb} & \textbf{Meaning (English)} & \textbf{Sens (Français)} \\ \hline
carry out & to perform, execute (a task, duty, plan) & accomplir, exécuter \\ \hline
give up & to stop trying; to abandon & abandonner, renoncer \\ \hline
stand for & to represent, symbolise; to tolerate (neg.) & représenter, signifier ; tolérer (nég.) \\ \hline
take part in & to participate in & participer à \\ \hline
put off & to postpone, delay & reporter, remettre à plus tard \\ \hline
\end{tabularx}
\end{center}
\emph{Rappel:} Verbe + particule(s). Sens souvent idiomatique. Séparables ou non selon l'objet (pronom $\rightarrow$ séparation obligatoire).
\end{aretenir}

\begin{exercice}[Phrasal Verbs: Match and Use]
\textbf{A.} Match each phrasal verb (1--5) to its meaning (a--e).
\begin{enumerate}
\item carry out \quad a. to stop doing something; to abandon
\item give up \quad b. to represent; to be a symbol of
\item stand for \quad c. to participate in
\item take part in \quad d. to postpone; to delay
\item put off \quad e. to do; to perform (a task, a duty)
\end{enumerate}
\textbf{B.} Use each phrasal verb in a sentence of your own about citizenship or work culture in Cameroon.
\end{exercice}

\begin{exercice}[Prepositions: Time, Place and Manner]
Choose the correct preposition (\emph{in, on, at, since, for, by, with}).
\begin{enumerate}
\item Cameroonians have voted in municipal elections \ldots many years.
\item The polling stations open \ldots 8 a.m. \ldots election day.
\item The meeting will take place \ldots the town hall of Buea.
\item She has been a member of the trade union \ldots 2018.
\item The workers go to the factory \ldots bus every morning.
\item He repaired the machine \ldots great care.
\item The independence of Cameroon was proclaimed \ldots 1960.
\item We discussed human rights \ldots the classroom yesterday.
\end{enumerate}
\end{exercice}

\courssub{Je me prépare au Bac}

\begin{exercice}[Reading Comprehension -- Type Probatoire]
Read the following press article carefully, then answer the questions.

\textbf{Young People and the Electoral Lists in Cameroon}

Every year, ELECAM (Elections Cameroon) organises campaigns to encourage young Cameroonians to register on the electoral lists. In many towns, such as Douala, Yaoundé and Garoua, mobile teams visit schools, markets and neighbourhoods. They explain that voting is not only a right but also a civic duty. However, many young people still do not register. Some say they are too busy with their studies or their work; others believe that their vote cannot change anything. Despite these difficulties, ELECAM registered thousands of new voters last year. ``Democracy needs the participation of everyone,'' declared one official in Bamenda. ``Although the process takes time, each registration is a victory for the nation.'' Since 2010, the number of registered voters has increased, yet the participation of young people remains low. Civil society organisations now carry out door-to-door campaigns, and they never give up, because they stand for democracy and basic freedoms.

\begin{enumerate}
\item Say whether the following statements are True or False. Justify with a quotation from the text.
\begin{enumerate}
\item ELECAM teams only work in Yaoundé.
\item All young people refuse to register on the electoral lists.
\item The number of registered voters has grown since 2010.
\end{enumerate}
\item Answer these questions in your own words.
\begin{enumerate}
\item What two reasons do some young people give for not registering?
\item What do civil society organisations do to promote registration?
\end{enumerate}
\item Find in the text:
\begin{enumerate}
\item one sentence containing the conjunction \emph{although};
\item one sentence containing an adverb of negation;
\item two phrasal verbs and give their meaning.
\end{enumerate}
\item Vocabulary: find words in the text that mean: (a) to write one's name on an official list; (b) an obligation of a citizen.
\end{enumerate}
\end{exercice}

\begin{exercice}[Grammar in Context -- Error Correction and Transformation]
\textbf{A. Error correction.} Each of the following sentences contains one mistake from this lesson (double negation, \emph{although\ldots but}, wrong preposition, wrong pronoun, wrong comparative). Rewrite each sentence correctly.
\begin{enumerate}
\item Although the workers were tired, but they finished the bridge on time.
\item The two brothers don't never help each others.
\item She works more hardly than her colleagues.
\item I have lived in Douala since five years.
\item The students didn't scarcely understand the lesson.
\item The members of the association respect each other (there are twelve members).
\item He is interested on politics and good in his work.
\item Despite of the rain, the campaign continued.
\item Hardly he had arrived when the meeting started.
\item The secretary types more quicklier than the director.
\end{enumerate}
\textbf{B. Sentence transformation.} Rewrite each sentence beginning with the word(s) given, using inversion where necessary.
\begin{enumerate}
\item He hardly ever participates in community work. $\rightarrow$ Hardly ever \ldots
\item She had scarcely finished her speech when the audience applauded. $\rightarrow$ Scarcely \ldots
\item The workers had never seen such a modern machine before. $\rightarrow$ Never before \ldots
\end{enumerate}
\end{exercice}

\begin{exercice}[Composition -- Type Baccalauréat Technique]
Write a composition of \textbf{180--200 words} on the following topic:

\textbf{``Tolerance is the key to peaceful co-existence in Cameroon.''}

In your essay, you should:
\begin{enumerate}
\item introduce the topic and explain what tolerance means;
\item give at least two concrete examples from life in Cameroon (school, workplace, neighbourhood, social media);
\item explain how tolerance promotes democracy and basic freedoms;
\item conclude with your personal opinion and a recommendation for young Cameroonians.
\end{enumerate}

Use at least \textbf{three conjunctions or linking expressions} from the lesson (\emph{although, despite, however, nevertheless, since, yet}), \textbf{two phrasal verbs} and \textbf{one adverb of negation}.

\textbf{Marking guide (out of 20):} task fulfilment and ideas: 8 marks; organisation and linking: 4 marks; grammar accuracy (conjunctions, prepositions, pronouns, adverbs): 4 marks; vocabulary and spelling: 4 marks.
\end{exercice}



\newpage

\coursec{Corrigés des exercices}

\coursec{Leçon 04 : Citizenship and Human Rights}

\courssub{Je vérifie mes connaissances}

\begin{corrige}[Exercice 1 -- QCM : Conjunctions and Prepositions]
\begin{enumerate}
\item \textbf{B. Although} -- La phrase contient une proposition subordonnée avec sujet + verbe (\emph{the rain was heavy}), donc il faut une conjonction de subordination de concession : \emph{Although}. \emph{Despite} exigerait un nom (\emph{Despite the heavy rain}) ; \emph{However} ne peut pas ouvrir une subordonnée ; \emph{Since} marquerait la cause, ce qui est illogique ici. \emph{Fr: Bien que la pluie fût forte, les ouvriers ont continué...}
\item \textbf{C. however} -- Après un point-virgule, on utilise un adverbe de liaison d'opposition : \emph{however} (suivi d'une virgule). \emph{Although} et \emph{because}/\emph{since} introduisent des subordonnées, incompatibles avec le point-virgule. \emph{Fr: L'équipe a travaillé dur ; cependant, elle n'a pas fini à temps.}
\item \textbf{C. for} -- \emph{For} + durée (\emph{for two hours}) avec le present perfect continu (\emph{have been cleaning}). \emph{Since} s'emploie avec un point de départ (\emph{since 10 a.m.}). \emph{Fr: Les élèves nettoient la cour depuis deux heures.}
\item \textbf{A. In spite of} -- Suivi d'un nom (\emph{his lack of experience}), il faut la locution prépositionnelle \emph{In spite of} (ou \emph{Despite}). \emph{Although} exigerait une proposition (\emph{Although he lacked experience}). \emph{Fr: Malgré son manque d'expérience, le jeune mécanicien a réparé le moteur parfaitement.}
\end{enumerate}
\textbf{Points de vigilance :} bien distinguer \emph{although} + S + V et \emph{despite/in spite of} + nom ; \emph{since} = point de départ, \emph{for} = durée.
\end{corrige}

\begin{corrige}[Exercice 2 -- True or False -- Justify]
\begin{enumerate}
\item \textbf{False.} \emph{Although} et \emph{but} ne se combinent jamais (redondance de la concession). Correction : \emph{Although the voters were many, they voted peacefully.} ou \emph{The voters were many, but they voted peacefully.}
\item \textbf{True.} Règle traditionnelle exigée à l'examen : \emph{each other} pour deux personnes ou entités ; \emph{one another} pour plus de deux. (L'usage moderne les rend souvent interchangeables, mais la distinction reste attendue au Probatoire/Bac technique.)
\item \textbf{True.} \emph{Hardly, scarcely, barely, rarely, never} ont déjà une valeur négative : la double négation est interdite. \emph{They didn't hardly work} $\rightarrow$ \emph{They hardly worked.}
\item \textbf{False.} Les prépositions correctes sont \emph{good \textbf{at} work} (bon en) et \emph{interested \textbf{in} politics} (intéressé par). \emph{Interested on} n'existe pas.
\end{enumerate}
\end{corrige}

\begin{corrige}[Exercice 3 -- Definitions]
Réponses indicatives (toute formulation correcte en anglais est acceptée) :
\begin{enumerate}
\item \textbf{Civic responsibility:} The duty of every citizen to take part in community life, respect the laws, vote and contribute to the common good. \emph{Fr: responsabilité civique}
\item \textbf{Tolerance:} Accepting and respecting opinions, beliefs or ways of life that are different from one's own, without hostility. \emph{Fr: tolérance}
\item \textbf{Excellence at work:} Doing one's professional tasks with high quality, punctuality, honesty and a constant wish to improve. \emph{Fr: excellence au travail}
\item \textbf{Basic freedoms:} The fundamental rights of every person, such as freedom of speech, of assembly, of religion and of movement. \emph{Fr: libertés fondamentales}
\end{enumerate}
\textbf{Points de vigilance :} une définition doit être une phrase complète, sans reprendre le mot défini comme seule explication, et rester dans le champ lexical de la leçon.
\end{corrige}

\begin{corrige}[Exercice 4 -- Comparative Adverbs]
\begin{enumerate}
\item \textbf{more fluently} -- Les adverbes en \emph{-ly} forment le comparatif avec \emph{more} : \emph{more fluently than}. \emph{Fr: plus couramment que}
\item \textbf{the hardest} -- \emph{Hard} (adverbe = avec effort) est irrégulier au sens où il ne prend pas \emph{-ly} ; superlatif : \emph{the hardest}. Ne pas confondre avec \emph{hardly} (= à peine). \emph{Fr: le plus assidûment}
\item \textbf{more quickly} -- Adverbe en \emph{-ly} $\rightarrow$ \emph{more quickly than}. (Forme familière : \emph{quicker}, à éviter à l'écrit d'examen.)
\item \textbf{the earliest} -- \emph{Early} suit la règle des adverbes courts : \emph{y} $\rightarrow$ \emph{i} + \emph{-est} ; superlatif : \emph{the earliest}. \emph{Fr: le plus tôt}
\end{enumerate}
\textbf{Points de vigilance :} ne jamais écrire \emph{more harder} ou \emph{more quicklier} (double marque du comparatif) ; vérifier la présence de \emph{the} au superlatif.
\end{corrige}

\courssub{Je m'entraîne}

\begin{corrige}[Exercice 5 -- Fill in the Blanks: Conjunctions]
\begin{enumerate}
\item \textbf{Although} -- concession + sujet + verbe (\emph{the workers were tired}).
\item \textbf{Despite} / \textbf{In spite of} -- concession + nom (\emph{the strike}).
\item \textbf{however} / \textbf{nevertheless} -- opposition après point-virgule, suivie d'une virgule.
\item \textbf{since} -- point de départ dans le temps (\emph{since 2015}) avec le present perfect.
\item \textbf{yet} / \textbf{nevertheless} / \textbf{still} -- opposition entre l'effort et le mauvais résultat.
\item \textbf{In spite of} / \textbf{Despite} -- concession + nom (\emph{the difficulties}).
\item \textbf{still} / \textbf{yet} / \textbf{nevertheless} -- opposition : le match a continué malgré la pluie.
\item \textbf{Since} / \textbf{As} / \textbf{Because} -- cause en tête de phrase (\emph{puisque le Cameroun est bilingue...}).
\item \textbf{however} / \textbf{nevertheless} / \textbf{still} -- opposition après point-virgule.
\item \textbf{still} / \textbf{yet} / \textbf{nevertheless} -- opposition ; \emph{still} se place avant le verbe principal.
\end{enumerate}
\textbf{Points de vigilance :} après \emph{however/nevertheless} en tête de seconde proposition, mettre une virgule ; \emph{despite} ne prend jamais \emph{of} ; \emph{since} + date ou événement, \emph{for} + durée.
\end{corrige}

\begin{corrige}[Exercice 6 -- Each Other / One Another]
\begin{enumerate}
\item \textbf{each other} -- deux candidats. \emph{The two candidates respect each other.}
\item \textbf{one another} -- les membres d'une coopérative (plus de deux). \emph{The members of the cooperative help one another.}
\item \textbf{each other} -- deux personnes (\emph{my friend and I}). \emph{My friend and I write letters to each other.}
\item \textbf{one another} -- les élèves d'une classe (groupe). \emph{The students of the class lend books to one another.}
\item \textbf{each other} -- deux personnes (le mari et la femme). \emph{The husband and the wife trust each other.}
\item \textbf{one another} -- cinq voisins (plus de deux). \emph{The five neighbours clean the quarter with one another.}
\end{enumerate}
\textbf{Points de vigilance :} jamais de \emph{s} à \emph{other} (\emph{each others} est faux) ; compter le nombre de participants avant de choisir.
\end{corrige}

\begin{corrige}[Exercice 7 -- Phrasal Verbs: Match and Use]
\textbf{A. Correspondances :} 1--e (\emph{carry out} = to perform a task), 2--a (\emph{give up} = to abandon), 3--b (\emph{stand for} = to represent), 4--c (\emph{take part in} = to participate), 5--d (\emph{put off} = to postpone).

\textbf{B. Phrases personnelles (exemples) :}
\begin{enumerate}
\item The election officials \textbf{carried out} the voter registration campaign efficiently. \emph{Fr: ont mené à bien}
\item Civil society organisations never \textbf{give up} promoting human rights. \emph{Fr: n'abandonnent jamais}
\item The national flag \textbf{stands for} unity and sovereignty. \emph{Fr: représente}
\item Every citizen should \textbf{take part in} community development activities. \emph{Fr: participer à}
\item The management \textbf{put off} the meeting until next Monday because of the public holiday. \emph{Fr: a reporté}
\end{enumerate}
\textbf{Points de vigilance :} le sens d'un phrasal verb est idiomatique (on ne peut pas le deviner mot à mot) ; avec un pronom objet, les verbes séparables imposent la coupure : \emph{carry it out}, \emph{put it off}.
\end{corrige}

\begin{corrige}[Exercice 8 -- Prepositions: Time, Place and Manner]
\begin{enumerate}
\item \textbf{for} -- durée : \emph{for many years}.
\item \textbf{at} 8 a.m. ; \textbf{on} election day -- heure précise $\rightarrow$ \emph{at} ; jour $\rightarrow$ \emph{on}.
\item \textbf{at} / \textbf{in} the town hall -- \emph{at} pour le lieu comme point, \emph{in} pour l'intérieur du bâtiment (les deux acceptés).
\item \textbf{since} -- point de départ : \emph{since 2018} (avec le present perfect).
\item \textbf{by} -- moyen de transport : \emph{by bus} (sans article).
\item \textbf{with} -- manière : \emph{with great care}.
\item \textbf{in} -- année : \emph{in 1960}.
\item \textbf{in} -- lieu clos : \emph{in the classroom}.
\end{enumerate}
\textbf{Points de vigilance :} mémoriser les trios \emph{at/on/in} du temps (heure / jour / mois-année) ; \emph{by} + transport sans article ; \emph{since} + point de départ, \emph{for} + durée.
\end{corrige}

\courssub{Je me prépare au Bac}

\begin{corrige}[Exercice 9 -- Reading Comprehension -- Type Probatoire]
\begin{enumerate}
\item \textbf{True or False (justification par citation obligatoire)}
\begin{enumerate}
\item \textbf{False.} ``In many towns, such as Douala, Yaoundé and Garoua, mobile teams visit schools, markets and neighbourhoods.'' (Les équipes ne travaillent pas \emph{uniquement} à Yaoundé.)
\item \textbf{False.} ``However, many young people still do not register.'' (\emph{Many} ne signifie pas \emph{all} : certains jeunes s'inscrivent, et le texte précise que des milliers de nouveaux électeurs ont été enregistrés.)
\item \textbf{True.} ``Since 2010, the number of registered voters has increased...''
\end{enumerate}
\item \textbf{Questions (réponses en phrases complètes, mots personnels)}
\begin{enumerate}
\item Some young people say they are too busy with their studies or their work, and others think that their vote cannot change anything.
\item They carry out door-to-door campaigns to encourage people to register, and they never give up.
\end{enumerate}
\item \textbf{Repérage grammatical dans le texte}
\begin{enumerate}
\item Conjonction \emph{although} : ``Although the process takes time, each registration is a victory for the nation.''
\item Adverbe de négation : \emph{never} dans ``...and they never give up...'' (on note aussi \emph{still} et \emph{yet} comme connecteurs).
\item Phrasal verbs : \textbf{carry out} (to perform, accomplish -- accomplir) ; \textbf{give up} (to stop trying -- abandonner) ; \textbf{stand for} (to represent, defend -- représenter, défendre).
\end{enumerate}
\item \textbf{Vocabulaire}
\begin{enumerate}
\item (a) \textbf{register} -- s'inscrire (sur une liste officielle) ; (b) \textbf{civic duty} -- devoir civique.
\end{enumerate}
\end{enumerate}
\textbf{Barème indicatif (sur 10) :} Q1 : 3 pts (0,5 pour la réponse + 0,5 pour la citation) ; Q2 : 2 pts ; Q3 : 3 pts ; Q4 : 2 pts.\\
\textbf{Points de vigilance :} toute réponse True/False sans citation du texte ne vaut que la moitié des points ; recopier la citation exactement, entre guillemets ; répondre aux questions ouvertes par des phrases complètes sans recopier tout le paragraphe.
\end{corrige}

\begin{corrige}[Exercice 10 -- Grammar in Context]
\textbf{A. Correction des erreurs}
\begin{enumerate}
\item \textbf{Although the workers were tired, they finished the bridge on time.} -- \emph{Although} et \emph{but} sont incompatibles (supprimer \emph{but}).
\item \textbf{The two brothers never help each other.} -- Double négation interdite (\emph{don't never} $\rightarrow$ \emph{never}) ; \emph{each other} sans \emph{s}.
\item \textbf{She works harder than her colleagues.} -- \emph{Hardly} signifie ``à peine'' ; l'adverbe ``avec effort'' est \emph{hard}, comparatif \emph{harder}.
\item \textbf{I have lived in Douala for five years.} -- \emph{For} + durée ; \emph{since} exige un point de départ (\emph{since 2019}).
\item \textbf{The students scarcely understood the lesson.} -- Double négation : \emph{scarcely} suffit, verbe à la forme affirmative.
\item \textbf{The members of the association respect one another.} -- Douze membres (plus de deux) $\rightarrow$ \emph{one another}.
\item \textbf{He is interested in politics and good at his work.} -- \emph{interested in}, \emph{good at}.
\item \textbf{Despite the rain, the campaign continued.} -- \emph{Despite} ne prend jamais \emph{of} (contrairement à \emph{in spite of}).
\item \textbf{Hardly had he arrived when the meeting started.} -- \emph{Hardly} en tête de phrase impose l'inversion sujet--auxiliaire : \emph{Hardly had he arrived...}
\item \textbf{The secretary types more quickly than the director.} -- Pas de double comparatif : \emph{more quickly} (ou \emph{quicker}), jamais \emph{more quicklier}.
\end{enumerate}
\textbf{B. Transformations avec inversion}
\begin{enumerate}
\item \textbf{Hardly ever does he participate in community work.} -- Inversion avec l'auxiliaire \emph{do} au présent : \emph{does he participate}.
\item \textbf{Scarcely had she finished her speech when the audience applauded.} -- Inversion au past perfect : \emph{had she finished}.
\item \textbf{Never before had the workers seen such a modern machine.} -- Inversion : \emph{had the workers seen}.
\end{enumerate}
\textbf{Points de vigilance :} après un adverbe négatif ou restrictif en tête de phrase (\emph{hardly, scarcely, never, rarely}), l'inversion sujet--auxiliaire est \emph{obligatoire} ; sans auxiliaire apparent, introduire \emph{do/does/did}.
\end{corrige}

\begin{corrige}[Exercice 11 -- Composition -- Type Baccalauréat Technique]
\textbf{Modèle de réponse (indicatif, environ 195 mots) :}

\emph{Tolerance: The Foundation of Peaceful Co-existence in Cameroon}

Tolerance means accepting and respecting differences in opinions, cultures and beliefs without hostility. In Cameroon, a country with more than 250 ethnic groups and two official languages, tolerance is not a luxury but a necessity for national unity.

\textbf{Although} our diversity is a strength, tensions sometimes arise. In schools, students from different regions must learn together; tolerance allows them to share ideas instead of fighting. Similarly, in workplaces such as the Douala port or the Yaoundé hospitals, colleagues with various backgrounds collaborate every day. \textbf{Despite} occasional misunderstandings, mutual respect guarantees productivity and excellence at work.

Moreover, tolerance sustains democracy. When citizens accept opposing views, elections remain peaceful and the freedoms of speech and assembly are protected. \textbf{Since} 1990, multipartism has required politicians to \textbf{carry out} debates without violence, and civil society organisations \textbf{stand for} human rights; they \textbf{never} \textbf{give up} promoting dialogue.

In conclusion, I firmly believe that young Cameroonians must cultivate tolerance. \textbf{Nevertheless}, this requires daily effort. I recommend that every student \textbf{take part in} intercultural clubs and refuse to spread hate speech on social media. Together, we can build a Cameroon where diversity means richness, not division.

\textbf{Barème indicatif (sur 20) :}
\begin{itemize}
\item Réalisation de la tâche et idées (introduction, deux exemples concrets, lien tolérance--démocratie, opinion personnelle et recommandation) : 8 pts
\item Organisation et articulation (paragraphes, connecteurs variés) : 4 pts
\item Correction grammaticale (conjonctions, prépositions, pronoms réciproques, adverbes, inversion éventuelle) : 4 pts
\item Lexique et orthographe (vocabulaire de la citoyenneté et du travail) : 4 pts
\end{itemize}
\textbf{Points de vigilance :} respecter la fourchette 180--200 mots ; employer au minimum trois connecteurs de la leçon, deux phrasal verbs et un adverbe de négation, en les utilisant correctement (pas de \emph{although... but}, pas de double négation, pas de \emph{despite of}) ; conclure par une opinion personnelle claire et une recommandation adressée aux jeunes Camerounais.
\end{corrige}

\newpage

\coursec{Fiche bilan}

\courssub{Mind Map of the Lesson}

\begin{popfigure}
\begin{center}
\begin{tikzpicture}[mindmap, grow cyclic,
  every node/.style={concept, font=\small},
  root concept/.append style={concept color=popblue, font=\bfseries\small, minimum size=2.6cm},
  level 1/.append style={level distance=3.4cm, sibling angle=60},
  level 1 concept/.append style={font=\scriptsize, minimum size=2.2cm}]
\node[root concept]{Citizenship and Human Rights}
  child[concept color=poporange]{ node {Work culture and excellence at work} }
  child[concept color=popgreen]{ node {Civic responsibility (Listening)} }
  child[concept color=poppurple]{ node {Democracy, tolerance, freedoms (Reading)} }
  child[concept color=popgold]{ node {Conjunctions: although, despite, however\ldots} }
  child[concept color=poppink]{ node {Pronouns: each, each other, one another} }
  child[concept color=popblue!70]{ node {Adverbs, negation, phrasal verbs, prepositions} };
\end{tikzpicture}
\end{center}
\legende{Carte mentale de la leçon : les six branches de l'unité \emph{Citizenship and Human Rights}.}
\end{popfigure}

\begin{bilan}
\textbf{1. Key vocabulary (à connaître par c\oe ur)}
\begin{itemize}
\item \textbf{Work culture:} \cle{punctuality} (ponctualité), \cle{commitment} (engagement), \cle{excellence} (excellence), \cle{team spirit} (esprit d'équipe), \cle{discipline} (discipline), \cle{efficiency} (efficacité).
\item \textbf{Civic responsibility:} \cle{citizen} (citoyen), \cle{to vote} (voter), \cle{to pay taxes} (payer ses impôts), \cle{to respect the law} (respecter la loi), \cle{community service} (service communautaire).
\item \textbf{Democracy and freedoms:} \cle{freedom of speech} (liberté d'expression), \cle{tolerance} (tolérance), \cle{human rights} (droits de l'homme), \cle{election} (élection), \cle{to demonstrate peacefully} (manifester pacifiquement).
\end{itemize}

\textbf{2. Grammar essentials}
\begin{center}
\begin{tabularx}{\linewidth}{|l|X|X|}\hline
\rowcolor{popblueL} \textbf{Point} & \textbf{Rule} & \textbf{Example} \\ \hline
Coordinating conjunctions & and, but, or, so, \cle{yet}, \cle{for} & He worked hard, \emph{yet} he stayed humble. \\ \hline
Subordinating conjunctions & \cle{although}, \cle{since}, because, while & \emph{Although} she was tired, she voted. \\ \hline
Concession & \cle{despite} / \cle{in spite of} + noun or -ing & \emph{Despite} the rain, they demonstrated. \\ \hline
Contrast linkers & \cle{however}, \cle{nevertheless}, \cle{still} & It was difficult; \emph{however}, they succeeded. \\ \hline
Reciprocal pronouns & \cle{each other} (2 people), \cle{one another} (more than 2) & Citizens must respect \emph{one another}. \\ \hline
\cle{each} & + singular noun and singular verb & \emph{Each} citizen has rights. \\ \hline
Prepositions & time: at, on, in, since, for; place: at, in, on; manner: by, with, in & The meeting is \emph{on} Monday \emph{at} 8 a.m. \\ \hline
Comparative adverbs & more + adv. + than; irregular: well $\to$ better, badly $\to$ worse & She speaks \emph{more clearly than} before. \\ \hline
Adverbs of negation & \cle{hardly}, \cle{scarcely}, \cle{never} (already negative: no ``not'') & He \emph{never} cheats at work. \\ \hline
Phrasal verbs & verb + particle: \cle{stand up for}, \cle{look after}, \cle{put off}, \cle{carry out}, \cle{give up} & We must \emph{stand up for} our rights. \\ \hline
\end{tabularx}
\end{center}

\textbf{3. Express methods}
\begin{itemize}
\item \textbf{Speaking:} greet $\to$ give your opinion (\emph{I think that\ldots}) $\to$ justify with \emph{because/since} $\to$ conclude.
\item \textbf{Listening:} read the questions first $\to$ listen for key words $\to$ note names, dates, actions.
\item \textbf{Reading:} skim for the general idea $\to$ scan for details $\to$ identify linkers (however, although) to follow the argument.
\item \textbf{Writing a composition:} introduction (present the topic) $\to$ 2--3 developed paragraphs with linkers $\to$ conclusion with your personal opinion. Use at least one conjunction of concession and one phrasal verb.
\end{itemize}
\end{bilan}

\begin{aretenir}[Checklist avant le Bac]
\begin{itemize}
\item $\square$ I can name at least six words about work culture and excellence.
\item $\square$ I can use \emph{although}, \emph{despite} and \emph{however} correctly in a sentence.
\item $\square$ I know that \emph{despite/in spite of} is followed by a noun or a verb in -ing, never by a full clause.
\item $\square$ I can distinguish \emph{each other} (two people) from \emph{one another} (more than two).
\item $\square$ I remember that \emph{each} takes a singular verb.
\item $\square$ I can choose the right preposition of time (at/on/in), place and manner.
\item $\square$ I can form the comparative of adverbs: \emph{more slowly than}, \emph{better}, \emph{worse}.
\item $\square$ I never use ``not'' with \emph{hardly}, \emph{scarcely} or \emph{never}.
\item $\square$ I know at least five phrasal verbs linked to citizenship (stand up for, give up, carry out\ldots).
\item $\square$ I can write a 150-word composition on democracy or tolerance with an introduction, a body and a conclusion.
\item $\square$ I can talk for one minute about civic responsibility using linkers.
\end{itemize}
\end{aretenir}

\findecours

\end{document}
